% Combined TeX source bundle

% Title: Arabic translation drafts

% Note: constituent files are preserved below in sources/. Some are standalone

% documents with their own preambles; this combined file is for inspection,

% comparison, and source packet rather than guaranteed one-command compilation.





% ============================================================

% Source: translations/ar/chinese/li-ye-ceyuan-haijing-vols1-3_arabic.tex

% ============================================================



% =============================================================================

% 測圓海鏡 — Chinese–Arabic Parallel Translation Edition

% Ceyuan Haijing (Sea Mirror of Circle Measurements), Volumes 1–3

% Original: Li Ye (Li Zhi, 1192–1279)

% Translation: Chinese → Arabic

% =============================================================================

\documentclass[11pt,a4paper]{article}



% -------- Core Packages --------

\usepackage{fontspec}

\usepackage{polyglossia}

\setmainlanguage{english}

\setotherlanguage{arabic}

\usepackage{amsmath,amssymb,amsthm}

\usepackage{geometry}

\usepackage{graphicx}

\usepackage{xcolor}

\usepackage{titlesec}

\usepackage{fancyhdr}

\usepackage{enumitem}

\usepackage{array,booktabs,longtable,multirow}

\usepackage{hyperref}

\usepackage{tocloft}

\usepackage{indentfirst}

\usepackage{xeCJK}

\usepackage{paracol}

\usepackage{tocloft}



% -------- CJK Font Setup --------

\setCJKmainfont{Microsoft YaHei}

\setCJKsansfont{Microsoft YaHei}

\setCJKmonofont{Microsoft YaHei}



% -------- Latin Font Setup --------

\setmainfont{Times New Roman}

\setsansfont{Arial}



% -------- Arabic Font Setup --------

\newfontfamily\arabicfont[Script=Arabic]{Arial}

\TeXXeTstate=1

\newcommand{\ar}[1]{\arabicfont\beginR #1\endR}



% -------- Page Geometry --------

\geometry{

  paperwidth=210mm,paperheight=297mm,

  left=20mm,right=20mm,top=25mm,bottom=25mm,

  headheight=14pt,footskip=30pt

}



% -------- Colors --------

\definecolor{titlecolor}{RGB}{45,55,72}

\definecolor{sectioncolor}{RGB}{55,65,81}

\definecolor{notecolor}{RGB}{120,53,15}

\definecolor{rulecolor}{RGB}{180,160,140}

\definecolor{prefacecolor}{RGB}{80,60,40}

\definecolor{problemcolor}{RGB}{30,60,90}

\definecolor{arabiclabel}{RGB}{139,69,19}

\definecolor{chineselabel}{RGB}{25,55,95}



% -------- Column Ratio --------

\columnratio{0.48,0.48}



% -------- Section Formatting --------

\titleformat{\section}

  {\normalfont\Large\bfseries\color{sectioncolor}}

  {\thesection}{1em}{}

\titleformat{\subsection}

  {\normalfont\large\bfseries\color{sectioncolor}}

  {\thesubsection}{1em}{}



% -------- Header/Footer --------

\pagestyle{fancy}

\fancyhf{}

\fancyhead[L]{\small\textsc{Ceyuan Haijing — Chinese–Arabic Parallel Edition}}

\fancyhead[R]{\small\thepage}

\renewcommand{\headrulewidth}{0.4pt}

\renewcommand{\footrulewidth}{0pt}



% -------- Custom Parallel Commands --------

\newcommand{\cn}[1]{\textnormal{#1}}

\newcommand{\Problem}[1]{\noindent\textbf{\color{problemcolor}#1}}

\newcommand{\Answer}[1]{\par\noindent\textbf{\color{problemcolor}答} #1\par}

\newcommand{\Method}[1]{\par\noindent\textbf{\color{problemcolor}法曰} #1\par}

\newcommand{\Workings}[1]{\par\noindent\textbf{\color{problemcolor}草曰} #1\par}

\newcommand{\Remark}[1]{\par\noindent\textit{〔按〕}#1\par}



% Arabic equivalents

\newcommand{\ProblemAr}[1]{\noindent\textbf{\color{problemcolor}\ar{#1}}}

\newcommand{\AnswerAr}{\par\noindent\textbf{\color{problemcolor}\ar{الجواب:}} }

\newcommand{\MethodAr}{\par\noindent\textbf{\color{problemcolor}\ar{الطريقة:}} }

\newcommand{\WorkingsAr}{\par\noindent\textbf{\color{problemcolor}\ar{الحل:}} }

\newcommand{\RemarkAr}{\par\noindent\ar{〔تعليق〕}}



% Parallel environment helpers

\newenvironment{paralleltext}{%

  \begin{paracol}{2}%

}{%

  \end{paracol}%

}



\newcommand{\lefttext}[1]{#1}

\newcommand{\righttext}[1]{\switchcolumn #1}



% Hyperref

\hypersetup{

  colorlinks=true,

  linkcolor=sectioncolor,

  citecolor=sectioncolor,

  urlcolor=sectioncolor,

  pdfencoding=auto

}



\setcounter{tocdepth}{2}



% =============================================================================

\begin{document}



% -------- Title Page --------

\begin{titlepage}

\centering

\vspace*{1.5cm}



{\fontsize{36}{44}\selectfont\CJKfamily{sf}\color{titlecolor}

測圓海鏡}



\vspace{0.6cm}



{\fontsize{16}{20}\selectfont\color{sectioncolor}

\ar{مرآة البحر لقياس الدائرة}}



\vspace{0.3cm}



{\fontsize{14}{18}\selectfont\color{sectioncolor}

Ceyuan Haijing — Sea Mirror of Circle Measurements}



\vspace{0.4cm}



{\large\color{rulecolor}\rule{0.6\textwidth}{0.4pt}}



\vspace{0.6cm}



{\LARGE 卷一 — 卷三}



{\large\color{sectioncolor}

\ar{المجلد الأول — المجلد الثالث}}



{\large Volumes I – III}



\vspace{1.2cm}



{\large\color{sectioncolor}

元 \quad 李冶 \quad 撰}



\vspace{0.3cm}



{\normalsize\color{sectioncolor}

\ar{لِـ يـِه (لي جِـه)، 1192–1279 م}}



{\normalsize\color{sectioncolor}

(Li Ye / Li Zhi, 1192–1279 CE)}



\vspace{1.5cm}



{\color{rulecolor}\rule{0.4\textwidth}{0.4pt}}



\vspace{0.8cm}



{\normalsize

以天元術解勾股容圓一百七十問}



\vspace{0.2cm}



{\normalsize\color{sectioncolor}

\ar{حلّ مئة وسؤالٍ وسبعين مسألة في الدائرة المحاطة بالمثلث القائم\\

بطريقة عنصر السماء}}



\vspace{0.4cm}



{\small\color{notecolor}

\textit{Using the Method of the Celestial Element (Tian Yuan Shu)}\\

\textit{to Solve 170 Problems on the Inscribed Circle}\\

\ar{باستخدام طريقة عنصر السماء لحلّ 170 مسألة في الدائرة المحاطة}}



\vspace{1.2cm}



{\small

A foundational text of medieval Chinese algebra,\\

completed in the year 1248 of the Common Era.

}



\vspace{0.3cm}



{\small\color{notecolor}

\ar{أحد النصوص الأساسية في جبر العصور الوسطى الصينية،\\

أُنجز في سنة 1248 ميلادية}}



\vfill



{\small

\textit{Bilingual Chinese–Arabic Edition / \ar{الطبعة الثنائية الصينية–العربية}}

}



\end{titlepage}



% -------- Table of Contents --------

\tableofcontents

\newpage



% =============================================================================

\section{Introduction / \ar{مقدمة}}

\label{sec:intro}



\begin{paracol}{2}



The \textit{Ceyuan Haijing} (測圓海鏡, literally ``Sea Mirror of Circle Measurements'') is the masterpiece of Li Ye (李冶, also known as Li Zhi, 1192–1279), one of the four great mathematicians of the Song--Yuan period of China, alongside Yang Hui, Qin Jiushao, and Zhu Shijie.



Completed in 1248, this work systematically applies the \textbf{tian yuan shu} (天元術, ``method of the celestial element'')---an algebraic method for setting up and solving polynomial equations---to a total of 170 geometric problems concerning a circle inscribed in or related to a right triangle (勾股容圓).



\switchcolumn



\ar{يعتبر كتاب «مرآة البحر لقياس الدائرة» (測圓海鏡) تحفة الرياضياتي لي يه (李冶، المعروف أيضًا باسم لي جيه، 1192–1279 م)، أحد العُظماء الأربعة في الرياضيات خلال فترة سونغ–يوان في الصين، إلى جانب يانغ هُوي وتشين جيُوشاو وتشو شِجيه.}



\vspace{0.3cm}



\ar{أُنجز هذا العمل في سنة 1248 م، ويطبّق بشكل منهجي \textbf{طريقة عنصر السماء} (天元術) — وهي طريقة جبرية لوضع وحلّ المعادلات متعدّدة الحدود — على 170 مسألة هندسية تتعلق بدائرة محاطة بمثلث قائم أو مرتبطة به (勾股容圓).}



\switchcolumn*



The book is organized as follows:

\begin{itemize}[leftmargin=2em]

  \item \textbf{Volume 1} establishes the theoretical foundation: the \textit{Circle City Diagram} (圓城圖式), the \textit{General Rate Names} (總率名號), the \textit{Current Question Numbers} (今問正數), and the \textit{Identification Miscellaneous Records} (識別雜記) containing 692 geometric formulas.

  \item \textbf{Volumes 2--12} present 170 problems, solved by the tian yuan shu method. The problems involve equations of degrees ranging from 1 to 6.

\end{itemize}



The present document contains the complete text of \textbf{Volumes 1--3}, comprising all foundational material and the first suite of problems.



\switchcolumn



\ar{يُنظَّم الكتاب على النحو التالي:}

\begin{itemize}[leftmargin=2em,rightmargin=0em]

  \item \ar{\textbf{المجلد الأول}: يضع الأسس النظرية: رسمة \textit{المدينة الدائرية} (圓城圖式)، و\textit{أسماء المعدّلات العامة} (總率名號)، و\textit{أعداد المسائل الحاضرة} (今問正數)، و\textit{سجلات التمييز المتنوعة} (識別雜記) التي تحتوي على 692 صيغة هندسية.}

  \item \ar{\textbf{المجلدات 2–12}: تعرض 170 مسألة، تُحلّ بطريقة عنصر السماء، وتشمل معادلات درجتها من الدرجة الأولى إلى السادسة.}

\end{itemize}



\vspace{0.3cm}



\ar{يحتوي هذا الكتاب على النص الكامل للمجلدات \textbf{1–3}، بما فيها جميع المواد الأساسية وأول مجموعة من المسائل.}



\end{paracol}



\vspace{0.5cm}

\noindent\textcolor{rulecolor}{\rule{\textwidth}{0.4pt}}



% =============================================================================

\newpage

\section{Original Preface / \ar{المدمة الأصلية}}

\label{sec:preface}



\begin{paracol}{2}



{\color{prefacecolor}

數本難窮，吾欲以力強窮之，彼其數不惟不能得其凡，而吾之力且憊矣。然則數果不可以窮耶？既已名之數矣，則又何為而不可窮也。

}



\vspace{0.3cm}



\switchcolumn



\ar{أصل الأعداد يصعب استنفاده، ولو حاولتُ إجباره بالقوة، لما استطعتُ الحصول على جوهره، بل وستنضب قواي. فهل الأعداد لا تُستنفد حقًّا؟ بما أنها سُمّيت «أعدادًا»، فكيف لا يمكن استنفادها؟}



\switchcolumn*



{\color{prefacecolor}

故謂數為難窮斯可，謂數為不可窮斯不可。何則？彼其冥冥之中，固有昭昭者存，夫昭昭者其自然之數也。非自然之數，其自然之理也。數一出於自然，吾欲以力強窮之，使隸首復生亦末如之何也已。

}



\switchcolumn



\ar{فمن الجائز القول إنّ الأعداد يصعب استنفادها، لكن لا يجوز القول إنها لا تُستنفد. ولماذا؟ ففي وسط الظلام يوجد نورٌ باطن، وهذا النور هو العدد الطبيعي. بل هو ليس العدد الطبيعي وحسب، بل هو المبدأ الطبيعي. بما أن الأعداد تنبثق من الطبيعة، فلو حاولتُ إجبارها بالقوة، حتى لو بُعث لي شو (مخترع الحساب الأسطوري) من جديد، لعجز عن الأمر.}



\switchcolumn*



{\color{prefacecolor}

苟能推自然之理，以明自然之數，則雖遠而乾端坤倪，幽而神情鬼狀，未有不合者矣。

}



\switchcolumn



\ar{فإذا أمكن استنباط المبادئ الطبيعية لإيضاح الأعداد الطبيعية، فحتى أقاصي السماء وأطراف الأرض، وحتى أسرار الأرواح والأشباح، لن يخلوا من التوافق.}



\switchcolumn*



{\color{prefacecolor}

予自幼喜算數，恆病夫考圓之術例出於牽強，殊乖於自然，如古率、徽率、密率之不同，截弧、截矢、截背之互見，內外諸角，析剖支條，莫不各自名家，與世作法。及反覆研究，卒無以當吾心焉。

}



\switchcolumn



\ar{منذ صغري أحببتُ الحساب، وكنتُ أشكو دومًا أنّ طرق دراسة الدائرة تبدو مُكرهة وغير طبيعية — فالنسب القديمة والنسب الدقيقة والنسب المُحكَمة تختلف، والأقواس والأوتار والظهور تتباين، والزوايا الداخلية والخارجية تُحلّل إلى فروع. كلّ مدرسة وضعت اسمها وطريقتها للعالم. غير أنّ دراستها المتكرّرة لم تُشبِع فضولي.}



\switchcolumn*



{\color{prefacecolor}

老大以來，得洞淵九容之說，日夕玩繹，而向之病我者，使爆然落去而無遗余。山中多暇，客有從余求其說者，於是乎又為衍之，遂累一百七十問。

}



\switchcolumn



\ar{في الكبر، حصلتُ على نظرية «تسعة استيعابات دونغ يوان» (洞淵九容)، وقد تأمّلتُها ليل نهار، فزال ما كان يؤلمني دفعةً واحدةً بلا أثر. وإذ كانت لديّ فراغاتٌ في الجبال، طلب منّي ضيوفي شرح هذه النظرية، فوسّعتُها حتى بلغت مئة وسبعين مسألة.}



\switchcolumn*



{\color{prefacecolor}

既成編，客復目之《測圓海鏡》，蓋取夫天臨海鏡之義也。

}



\switchcolumn



\ar{ولما اكتمل التأليف، أطلق عليه أحد الضيوف اسم «مرآة البحر لقياس الدائرة»، مستوحىً من معنى «السماء التي تُشرف على مرآة البحر».}



\switchcolumn*



{\color{prefacecolor}

昔半山老人集《唐百家詩選》，自謂廢日力於此，良可惜。明道先生以上蔡謝君記誦為玩物喪志。夫文史尚矣，猶之為不足貴，況九九賤技能乎！

}



\switchcolumn



\ar{لقد جمع «شيخ نصف الجبل» مختاراتٍ من شعر مئة شاعرٍ من تانغ، معتبرًا ذلك إضاعةً للوقت. وقد زمَر مستر مينغداو بشي شيه من شانغتساي لحفظه الشعر. وإذا كان الأدب والتاريخ — رغم علوّ شأنهما — لا يُعَدّان ثمينَين، فكيف بمهارة الضرب «تسعة في تسعة» الهزيلة!}



\switchcolumn*



{\color{prefacecolor}

嗜好酸鹹，平生每痛自戒勅，竟莫能已，類有物慿之者，吾亦不知其然而然也。

}



\switchcolumn



\ar{إنّ حبّ الملح والخل، فقد حاولتُ طوال حياتي توبيخ نفسي والامتناع عنها، لكنني عجزتُ، كأنّ شيئًا ما يدفعني، ولا أدري كيف يحصل هذا.}



\switchcolumn*



{\color{prefacecolor}

故嘗私為之解曰：由技兼於事者言之，夷之禮，夔之樂，亦不免為一技。由技進乎道者言之，石之斤，扁之輪，非聖人之所與乎。

}



\switchcolumn



\ar{ولذا فسّرتُ ذلك لنفسي: من جهة تطبيق المهارات، فإنّ طقوس يي وموسيقى كويي لا تخلوا من كونها مهارة. ومن جهة الارتقاء بالمهارة إلى الطريق، فإنّ فأس بن شي وعجلة بيان — أليس هما مما يُثنى عليه الحكماء؟}



\switchcolumn*



{\color{prefacecolor}

覽吾之編，察吾苦心，其憫我者當百數，其笑我者當千數。乃若吾之所自得，則自得焉耳，寧復為人憫笑計哉。



\hfill 李冶序

}



\switchcolumn



\ar{من يطالع تأليفي ويلمس جهدي، سيجد من يشفق عليّ بالمئات، ومن يسخروا مني بالألوف. أمّا ما اكتسبته لنفسي فهو لي وحدي، وهل أبالي بمن يشفق أو يسخر؟}



\vspace{0.3cm}



\ar{\hfill تقديم لي يه}



\end{paracol}



\newpage





% =============================================================================

\section{Volume 1 / \ar{المجلد الأول} \quad 卷一}

\label{sec:vol1}



\begin{paracol}{2}



Volume 1 is the theoretical foundation of the entire work. It comprises four major sections:



\begin{enumerate}[leftmargin=2em]

  \item \textbf{Circle City Diagram} (圓城圖式) --- the master geometric configuration

  \item \textbf{General Rate Names} (總率名號) --- definitions of all 15 right triangles and their elements

  \item \textbf{Current Question Numbers} (今問正數) --- numerical values for all line segments

  \item \textbf{Identification Miscellaneous Records} (識別雜記) --- 692 geometric formulas

\end{enumerate}



\switchcolumn



\ar{المجلد الأول هو الأساس النظري للعمل بأكمله. ويتألّف من أربعة أقسام رئيسية:}



\begin{enumerate}[leftmargin=2em,rightmargin=0em]

  \item \ar{\textbf{رسمة المدينة الدائرية} (圓城圖式) — التكوين الهندسي الرئيسي}

  \item \ar{\textbf{أسماء المعدّلات العامة} (總率名號) — تعريف المثلثات القائمة الخمسة عشر وعناصرها}

  \item \ar{\textbf{أعداد المسائل الحاضرة} (今問正數) — القيم الرقمية لجميع الخطوط}

  \item \ar{\textbf{سجلات التمييز المتنوعة} (識別雜記) — 692 صيغة هندسية}

\end{enumerate}



\end{paracol}



% ------------------------------------------------------------------------------

\subsection{Circle City Diagram / \ar{رسمة المدينة الدائرية} (圓城圖式)}

\label{subsec:diagram}



\begin{paracol}{2}



The \textit{Circle City Diagram} is the master geometric configuration upon which all 170 problems are based. It consists of a large right triangle (called 勾股形, ``leg--leg shape'') with its inscribed circle, together with specific points and lines.



\vspace{0.3cm}



The scenario Li Ye constructs is:



\begin{quote}

假令有圓城一所，不知周徑，四面開門，門外縱橫各有十字大道。其西北十字道頭定為乾地，其東北十字道頭定為艮地，其東南十字道頭定為巽地，其西南十字道頭定為坤地。所有測望雜法一一設問如後。

\end{quote}



\switchcolumn



\ar{إنّ \textit{رسمة المدينة الدائرية} هي التكوين الهندسي الرئيسي الذي تستند إليه جميع المسائل الـ170. وتتألّف من مثلث قائم كبير (يُسمّى 勾股形، «شكل الضلع–الضلع») مع دائرته المحاطة، بالإضافة إلى نقاط وخطوط مُحدّدة.}



\vspace{0.3cm}



\ar{والمنطق الذي بناه لي يه هو:}



\begin{quote}

\ar{لنفترض مدينة دائرية محيطها وقطرها مجهولان، لها أبواب على الجوانب الأربعة، وخارج كلّ باب طرقات متقاطعة شمالية–جنوبية وشرقية–غربية. الطرقات المتقاطعة الشمالية–الغربية تُسمّى «تشيين» (乾)، والشمالية–الشرقية «كين» (艮)، والجنوبية–الشرقية «شون» (巽)، والجنوبية–الغربية «كون» (坤). جميع طرق القياس مفصّلة في المسائل أدناه.}

\end{quote}



\switchcolumn*



\textbf{The 22 Named Points.} The largest right triangle has vertices \textbf{Tian} (天, Heaven), \textbf{Di} (地, Earth), and \textbf{Qian} (乾). The center of the inscribed circle is called \textbf{Xin} (心, Heart). Key points include intersections of horizontal and vertical lines through the center and other constructed lines. Each point is designated by a single Chinese character:



\begin{center}

\begin{tabular}{llllllll}

\toprule

天 & 地 & 乾 & 心 & 日 & 南 & 北 & 川 \\

東 & 西 & 艮 & 坤 & 山 & 月 & 巽 & 青 \\

泉 & 朱 & 金 & 泛 & 旦 & 夕 & & \\

\bottomrule

\end{tabular}

\end{center}



\switchcolumn



\ar{\textbf{النقاط الثانية والعشرون المُسمّاة.} المثلث القائم الأكبر له رؤوس \textbf{تيان} (天، السماء)، و\textbf{دي} (地، الأرض)، و\textbf{تشيين} (乾). ويُسمّى مركز الدائرة المحاطة \textbf{شين} (心، القلب). وتشمل النقاط الرئيسية تقاطعات الخطوط الأفقية والعمودية عبر المركز وخطوط أخرى مُنشأة. كلّ نقطة تُحدّد بحرفٍ صينيّ واحد:}



\begin{center}

\begin{tabular}{llllllll}

\toprule

天 & 地 & 乾 & 心 & 日 & 南 & 北 & 川 \\

東 & 西 & 艮 & 坤 & 山 & 月 & 巽 & 青 \\

泉 & 朱 & 金 & 泛 & 旦 & 夕 & & \\

\bottomrule

\end{tabular}

\end{center}



\end{paracol}



% ------------------------------------------------------------------------------

\subsection{General Rate Names / \ar{أسماء المعدّلات العامة} (總率名號)}

\label{subsec:names}



\begin{paracol}{2}



The work studies \textbf{15 right triangles}, all having a segment of the \textit{Tiandi} (天地) line as their hypotenuse (弦, ``string''). The \textit{General Rate Names} define these triangles and their three sides.



In each triangle:

\begin{itemize}[leftmargin=2em]

  \item \textbf{弦} ($c$) --- the hypotenuse (斜邊)

  \item \textbf{股} ($b$) --- the longer vertical leg

  \item \textbf{勾} ($a$) --- the shorter horizontal leg

\end{itemize}



\switchcolumn



\ar{يدرس هذا الكتاب \textbf{15 مثلثًا قائمًا}، جميعها يتخذ قطعةً من خط \textit{تياندي} (天地) وترًا له (弦، «الوتر»). وتحدّد \textit{أسماء المعدّلات العامة} هذه المثلثات وأضلاعها الثلاثة.}



\vspace{0.3cm}



\ar{في كلّ مثلث:}

\begin{itemize}[leftmargin=2em,rightmargin=0em]

  \item \ar{\textbf{الوتر} ($c$) — الوتر (弦)}

  \item \ar{\textbf{الضلع الطويل} ($b$) — الضلع الرأسي الأطول (股)}

  \item \ar{\textbf{الضلع القصير} ($a$) — الضلع الأفقي الأقصر (勾)}

\end{itemize}



\end{paracol}



\begin{center}

\begin{tabular}{|c|l|c|c|c|}

\hline

\textbf{No.} & \textbf{Name / \ar{الاسم}} & \textbf{\ar{الوتر} 弦} & \textbf{\ar{الضلع} 股} & \textbf{\ar{القاعدة} 勾} \\

\hline

1 & \ar{تونغ (العام)} 通 & 通弦 & 通股 & 通勾 \\

2 & \ar{بيان (الجانب)} 邊 & 邊弦 & 邊股 & 邊勾 \\

3 & \ar{دي (القاعدة)} 底 & 底弦 & 底股 & 底勾 \\

4 & \ar{هوانغ غوانغ} 黃廣 & 黃廣弦 & 黃廣股 & 黃廣勾 \\

5 & \ar{هوانغ تشانغ} 黃長 & 黃長弦 & 黃長股 & 黃長勾 \\

6 & \ar{شانغ غاو (العلوي)} 上高 & 上高弦 & 上高股 & 上高勾 \\

7 & \ar{شيا غاو (السفلي)} 下高 & 下高弦 & 下高股 & 下高勾 \\

8 & \ar{شانغ بينغ} 上平 & 上平弦 & 上平股 & 上平勾 \\

9 & \ar{شيا بينغ} 下平 & 下平弦 & 下平股 & 下平勾 \\

10 & \ar{دا تشا (الكبير)} 大差 & 大差弦 & 大差股 & 大差勾 \\

11 & \ar{شياو تشا (الصغير)} 小差 & 小差弦 & 小差股 & 小差勾 \\

12 & \ar{هوانغ جي (القصوي)} 皇極 & 皇極弦 & 皇極股 & 皇極勾 \\

13 & \ar{تاي شُو (الخالص)} 太虛 & 太虛弦 & 太虛股 & 太虛勾 \\

14 & \ar{مينغ (الضوء)} 明 & 明弦 & 明股 & 明勾 \\

15 & \ar{تشوان (التركيز)} 叀 & 叀弦 & 叀股 & 叀勾 \\

\hline

\end{tabular}

\end{center}



\begin{paracol}{2}

\small Note: 上高 = 下高 and 上平 = 下平, so among 15 triangles, only 13 are distinct.

\switchcolumn

\small \ar{ملاحظة: 上高 = 下高 و 上平 = 下平، لذا من بين 15 مثلثًا، هناك 13 مثلثًا مميزًا فقط.}

\end{paracol}



% ------------------------------------------------------------------------------

\subsection{Current Question Numbers / \ar{أعداد المسائل الحاضرة} (今問正數)}

\label{subsec:numbers}



\begin{paracol}{2}



This section gives the numerical value of every line segment in the diagram. The \textbf{radius of the inscribed circle is taken as 120 steps} (步), which serves as the standard unit.



\vspace{0.3cm}



\textbf{1. 通 (Tong / General)}



\begin{center}

\begin{tabular}{ll}

弦六百八十 & 勾三百二十 \quad 股六百 \\

勾股和九百二十 & 較二百八十 \\

勾和一千 & 較三百六十 \\

股和一千二百八十 & 較八十 \\

較和九百六十 & 較四百 \\

和和一千六百 & 較二百四十 \\

\end{tabular}

\end{center}



\textit{That is:} hypotenuse $c = 680$, short leg $a = 320$, long leg $b = 600$.



\switchcolumn



\ar{يعطي هذا القسم القيمة الرقمية لكلّ قطعة خط في الرسمة. \textbf{يُؤخذ نصف قطر الدائرة المحاطة بوحدة 120 خطوة} (步)، وهي تُستخدم كوحدة قياس معيارية.}



\vspace{0.3cm}



\ar{\textbf{1. تونغ (العام)}}



\begin{center}

\begin{tabular}{ll}

\ar{الوتر 680} & \ar{القاعدة 320 \quad الضلع 600} \\

\ar{مجموع القاعدة والضلع 920} & \ar{الفرق 280} \\

\ar{مجموع القاعدة 1000} & \ar{الفرق 360} \\

\ar{مجموع الضلع 1280} & \ar{الفرق 80} \\

\ar{مجموع الفرق 960} & \ar{الفرق 400} \\

\ar{المجموع الكلي 1600} & \ar{الفرق 240} \\

\end{tabular}

\end{center}



\vspace{0.3cm}



\ar{أي: الوتر $c = 680$، القاعدة $a = 320$، الضلع $b = 600$.}



\end{paracol}



\vspace{0.3cm}



\begin{paracol}{2}

\textbf{2. 邊 (Bian / Side)}



\begin{center}

\begin{tabular}{ll}

弦五百四十四 & 勾二百五十六 \quad 股四百八十 \\

勾股和七百三十六 & 較二百二十四 \\

勾和八百 & 較二百八十八 \\

股和一千零二十四 & 較六十四 \\

較和七百六十八 & 較三百二十 \\

和和一千二百八十 & 較一百九十二 \\

\end{tabular}

\end{center}



\switchcolumn



\ar{\textbf{2. بيان (الجانب)}}



\begin{center}

\begin{tabular}{ll}

\ar{الوتر 544} & \ar{القاعدة 256 \quad الضلع 480} \\

\ar{مجموع القاعدة والضلع 736} & \ar{الفرق 224} \\

\ar{مجموع القاعدة 800} & \ar{الفرق 288} \\

\ar{مجموع الضلع 1024} & \ar{الفرق 64} \\

\ar{مجموع الفرق 768} & \ar{الفرق 320} \\

\ar{المجموع الكلي 1280} & \ar{الفرق 192} \\

\end{tabular}

\end{center}



\end{paracol}



\vspace{0.3cm}



\begin{paracol}{2}

\textbf{3. 底 (Di / Base)}



\begin{center}

\begin{tabular}{ll}

弦四百二十五 & 勾二百 \quad 股三百七十五 \\

勾股和五百七十五 & 較一百七十五 \\

勾和六百二十五 & 較二百二十五 \\

股和八百 & 較五十 \\

較和六百 & 較二百五十 \\

和和一千 & 較一百五十 \\

\end{tabular}

\end{center}



\switchcolumn



\ar{\textbf{3. دي (القاعدة)}}



\begin{center}

\begin{tabular}{ll}

\ar{الوتر 425} & \ar{القاعدة 200 \quad الضلع 375} \\

\ar{مجموع القاعدة والضلع 575} & \ar{الفرق 175} \\

\ar{مجموع القاعدة 625} & \ar{الفرق 225} \\

\ar{مجموع الضلع 800} & \ar{الفرق 50} \\

\ar{مجموع الفرق 600} & \ar{الفرق 250} \\

\ar{المجموع الكلي 1000} & \ar{الفرق 150} \\

\end{tabular}

\end{center}



\end{paracol}



\vspace{0.3cm}



\begin{paracol}{2}

\textbf{4. 黃廣 (Huang Guang)}



\begin{center}

\begin{tabular}{ll}

弦五百一十 & 勾二百四十 \quad 股四百五十 \\

\multicolumn{2}{c}{(即城徑也 --- this is the city diameter)} \\

勾股和六百九十 & 較二百一十 \\

\end{tabular}

\end{center}



\vspace{0.3cm}



The remaining 11 triangles follow the same pattern, each with 13 numerical quantities, giving a total of $15 \times 13 = 195$ numerical entries.



\switchcolumn



\ar{\textbf{4. هوانغ غوانغ}}



\begin{center}

\begin{tabular}{ll}

\ar{الوتر 510} & \ar{القاعدة 240 \quad الضلع 450} \\

\multicolumn{2}{c}{\ar{(هذا هو قطر المدينة)}} \\

\ar{مجموع القاعدة والضلع 690} & \ar{الفرق 210} \\

\end{tabular}

\end{center}



\vspace{0.3cm}



\ar{تتبع المثلثات الـ11 المتبقية النمط نفسه، ولكلّ منها 13 كمية رقمية، ما يعطي إجمالاً $15 \times 13 = 195$ مدخلاً رقميًا.}



\end{paracol}



\newpage





% ------------------------------------------------------------------------------

\subsection{Identification Miscellaneous Records / \ar{سجلات التمييز المتنوعة} (識別雜記)}

\label{subsec:identification}



\begin{paracol}{2}



This is the theoretical heart of the work. It contains \textbf{692 geometric formulas} relating the various line segments in the diagram. These formulas are purely geometric theorems, independent of any particular numerical values. The problems in all subsequent volumes are solved by using these formulas together with the tian yuan shu.



The section is divided into eight parts:



\begin{enumerate}[leftmargin=2em]

  \item \textbf{Various Names} (諸雜名目) --- general definitions and 30+ theorems

  \item \textbf{Five Sums and Five Differences} (五和五較)

  \item \textbf{Various Hypotenuses} (諸弦)

  \item \textbf{Large and Small Differences} (大小差)

  \item \textbf{Various Differences} (諸差)

  \item \textbf{Mutual Relations of Rates} (諸率互見)

  \item \textbf{Four-Place Supplement} (四位拾遺)

  \item \textbf{Supplement} (拾遺)

\end{enumerate}



\switchcolumn



\ar{هذا هو القلب النظري للعمل. ويحتوي على \textbf{692 صيغة هندسية} تربط الخطوط المختلفة في الرسمة. هذه الصيغ هي مبرهنات هندسية خالصة، مستقلة عن أي قيم رقمية مُعيّنة. وتُحلّ مسائل المجلدات التالية باستخدام هذه الصيغ معاً مع طريقة عنصر السماء.}



\vspace{0.3cm}



\ar{يُقسّم هذا القسم إلى ثمانية أجزاء:}



\begin{enumerate}[leftmargin=2em,rightmargin=0em]

  \item \ar{\textbf{الأسماء المتنوعة} (諸雜名目) — تعريفات عامة وأكثر من 30 مبرهنة}

  \item \ar{\textbf{المجاميع الخمسة والفروق الخمسة} (五和五較)}

  \item \ar{\textbf{الأوتار المختلفة} (諸弦)}

  \item \ar{\textbf{الفروق الكبيرة والصغيرة} (大小差)}

  \item \ar{\textbf{الفروق المختلفة} (諸差)}

  \item \ar{\textbf{العلاقات المتبادلة للمعدّلات} (諸率互見)}

  \item \ar{\textbf{التكميلة الرباعية} (四位拾遺)}

  \item \ar{\textbf{التكميلة} (拾遺)}

\end{enumerate}



\end{paracol}



\vspace{0.3cm}



\noindent\textbf{Sample Definitions from 諸雜名目 / \ar{نماذج تعريفات من «الأسماء المتنوعة»}:}



\vspace{0.3cm}



\begin{center}

\begin{tabular}{lll}

\toprule

\textbf{Name} & \textbf{Definition} & \textbf{\ar{التعريف}} \\

\midrule

內率 & 明勾 $\times$ 叀股 & \ar{القاعدة المضيئة $\times$ ضلع التركيز} \\

外率 & 明股 $\times$ 叀勾 & \ar{ضلع المضيء $\times$ قاعدة التركيز} \\

虛率 & 虛勾 $\times$ 虛股 & \ar{القاعدة الخالصة $\times$ الضلع الخالص} \\

角差 & 高股 $-$ 平勾 & \ar{ضلع المرتفع $-$ قاعدة المستوي} \\

次差 & (明股$-$明勾) $+$ (叀股$-$叀勾) & \ar{(ضلع المضيء$-$قاعدته) $+$ (ضلع التركيز$-$قاعدته)} \\

混同和 & 黃長股 $+$ 黃廣勾 & \ar{ضلع هوانغ تشانغ $+$ قاعدة هوانغ غوانغ} \\

傍差 & (明股$-$明勾) $-$ (叀股$-$叀勾) & \ar{(ضلع المضيء$-$قاعدته) $-$ (ضلع التركيز$-$قاعدته)} \\

\bottomrule

\end{tabular}

\end{center}



\vspace{0.3cm}



\noindent\textbf{Sample Theorems / \ar{نماذج من المبرهنات}:}



\vspace{0.3cm}



\begin{paracol}{2}



\begin{itemize}[leftmargin=2em]

  \item 凡大差小差相乘為半段徑冪。

  

  \textit{The product of the large difference and small difference equals half the square of the diameter.}

  

  \item 大差勾小差股相乘同上。

  

  \textit{The product of the large-difference short-leg and the small-difference long-leg equals the same.}

  

  \item 黃廣股黃長勾相乘為經冪。

  

  \textit{The product of the Huangguang long-leg and Huangchang short-leg equals the square of the diameter.}

  

  \item 勾股和即弦黃和。

  

  \textit{The sum of the two legs equals the sum of the hypotenuse and the diameter of the inscribed circle.}

  

  That is: $a + b = c + d$, where $d$ is the diameter.

\end{itemize}



\switchcolumn



\begin{itemize}[leftmargin=2em,rightmargin=0em]

  \item \ar{إنّ حاصل ضرب الفرق الكبير والفرق الصغير يساوي نصف مربّع القطر.}



  \item \ar{حاصل ضرب قاعدة الفرق الكبير وضلع الفرق الصغير يساوي المقدار نفسه.}



  \item \ar{حاصل ضرب ضلع هوانغ غوانغ وقاعدة هوانغ تشانغ يساوي مربّع القطر.}



  \item \ar{مجموع الضلعين القائمين يساوي مجموع الوتر وقطر الدائرة المحاطة.}



  \vspace{0.3cm}



  \ar{أي: $a + b = c + d$، حيث $d$ هو قطر الدائرة المحاطة.}

\end{itemize}



\end{paracol}



\newpage



% =============================================================================

\section{Volume 2 / \ar{المجلد الثاني} \quad 卷二}

\label{sec:vol2}



\begin{paracol}{2}



\textbf{正率一十四問} --- Fourteen Questions on Standard Rates



These fourteen problems form the ``Dongyuan Nine Containments'' (洞淵九容) sequence, derived from the earlier work of the Daoist master Li Sicong (李思聰, known as ``Master Dongyuan''). They present the nine fundamental types of circles contained in or related to a right triangle, plus five additional problems.



\switchcolumn



\ar{\textbf{أربعة عشر مسألة في المعدّلات المعيارية}}



\vspace{0.3cm}



\ar{تُشكّل هذه المسائل الأربعة عشر سلسلة «تسعة استيعابات دونغ يوان» (洞淵九容)، المستوحاة من عمل السيّد الطاووي لي سِتسونغ (李思聰، المعروف باسم «سيّد دونغ يوان»). وهي تعرض تسعة أنواع أساسية من الدوائر المحاطة بالمثلث القائم أو المرتبطة به، بالإضافة إلى خمس مسائل إضافية.}



\end{paracol}



% ------------------------------------------------------------------------------

\subsubsection*{Problem 1 / \ar{المسألة الأولى} \quad 第一問}



\begin{paracol}{2}



\noindent 或問：甲乙二人俱在乾地，乙東行三百二十步而立，甲南行六百步望見乙，問徑幾里？



\vspace{0.2cm}



\Answer{城徑二百四十步。}



\vspace{0.2cm}



\Method{此為勾股容圓也。以勾股相乘，倍之為實，並勾股冪以求弦，復加入勾股共，以為法。}



\vspace{0.2cm}



\Workings{置甲南行六百步在地，以乙東行三百二十步乘之，得一十九萬二千步，倍之得三十入萬四千步為實。以乙東行步自之，得一十萬零二千四百步為勾冪；以甲南行步自之，得三十六萬步為股冪。二冪相並，得四十六萬二千四百步為方實，以平方開之，得六百八十步，則弦也。以弦加勾股共，共得一千六百步以為法。如法而一，得二百四十步，則城徑也。合問。}



\switchcolumn



\ar{سُئل: رجلان أ وب يقفان في موقع تشيين. يسير ب نحو الشرق 320 خطوة ويقف، ويسير أ نحو الجنوب 600 خطوة ويرى ب. فكم قطر المدينة؟}



\vspace{0.2cm}



\AnswerAr{قطر المدينة 240 خطوة.}



\vspace{0.2cm}



\MethodAr{هذه مسألة الدائرة المحاطة بالمثلث القائم. اضرب الضلعين القائمين واضعف الناتج فتكون الذات، واجمع مربّعي الضلعين لاستخراج الوتر، ثم أضف مجموع الضلعين فتكون القاسية.}



\vspace{0.2cm}



\WorkingsAr{ضع مسافة سير أ نحو الجنوب (600 خطوة) واضربها بمسافة سير ب نحو الشرق (320 خطوة)، تحصل على 192000 خطوة، واضعفها فتحصل على 384000 خطوة وهي الذات. ربّع مسافة سير ب نحو الشرق فتحصل على 102400 وهي مربّع القاعدة؛ ربّع مسافة سير أ نحو الجنوب فتحصل على 360000 وهي مربّع الضلع. اجمع المربّعين فتحصل على 462400 وهي الذات المربّعة، واستخرج الجذر التربيعي فتحصل على 680 خطوة وهي الوتر. أضف الوتر إلى مجموع الضلعين، يصبح المجموع 1600 خطوة وهو القاسية. اقسم الذات على القاسية فتحصل على 240 خطوة وهي قطر المدينة. يتطابق مع المطلوب.}



\end{paracol}



\vspace{0.5cm}



% ------------------------------------------------------------------------------

\subsubsection*{Problem 2 / \ar{المسألة الثانية} \quad 第二問}



\begin{paracol}{2}



\noindent 或問：甲乙二人俱在西門，乙東行二百五十六步，甲南行四百八十步望見乙，問答同前。



\vspace{0.2cm}



\Answer{城徑二百四十步。}



\vspace{0.2cm}



\Method{此為勾上容圓也。以勾股相乘，倍之為實，並勾股冪以求弦，加入股以為法。}



\vspace{0.2cm}



\Workings{置甲南行四百八十步在地，以乙東行二百五十六步乘之，得一十二萬二千八百八十步，倍之得二十四萬五千七百六十步為實。以乙東行步自之，得六萬五千五百三十六步為勾冪；以甲南行步自之，得二十三萬零四百步為股冪。勾股冪相並，得二十九萬五千九百三十六步為方實，以平方開之，得五百四十四步為弦也。以弦加入南行步，共得一千零二十四步以為法。而一，得二百四十步則城徑。合問。}



\switchcolumn



\ar{سُئل: رجلان أ وب يقفان عند الباب الغربي. يسير ب نحو الشرق 256 خطوة، ويسير أ نحو الجنوب 480 خطوة ويرى ب. الجواب كالسابق.}



\vspace{0.2cm}



\AnswerAr{قطر المدينة 240 خطوة.}



\vspace{0.2cm}



\MethodAr{هذه مسألة الدائرة المحاطة على القاعدة. اضرب الضلعين القائمين واضعف الناتج فتكون الذات، واجمع مربّعي الضلعين لاستخراج الوتر، ثم أضف الضلع فتكون القاسية.}



\vspace{0.2cm}



\WorkingsAr{ضع مسافة سير أ نحو الجنوب (480 خطوة) واضربها بمسافة سير ب نحو الشرق (256 خطوة)، تحصل على 122880 خطوة، واضعفها فتحصل على 245760 خطوة وهي الذات. ربّع مسافة سير ب نحو الشرق فتحصل على 65536 وهي مربّع القاعدة؛ ربّع مسافة سير أ نحو الجنوب فتحصل على 230400 وهي مربّع الضلع. اجمع مربّعي الضلعين فتحصل على 295936 وهي الذات المربّعة، واستخرج الجذر التربيعي فتحصل على 544 خطوة وهي الوتر. أضف الوتر إلى مسافة السير نحو الجنوب، يصبح المجموع 1024 خطوة وهو القاسية. اقسم الذات على القاسية فتحصل على 240 خطوة وهي قطر المدينة. يتطابق مع المطلوب.}



\end{paracol}



\vspace{0.5cm}



% ------------------------------------------------------------------------------

\subsubsection*{Problem 3 / \ar{المسألة الثالثة} \quad 第三問}



\begin{paracol}{2}



\noindent 或問：甲乙二人俱在北門，乙東行二百步而止，甲南行三百七十五步望見乙，問答同前。



\vspace{0.2cm}



\Answer{城徑二百四十步。}



\vspace{0.2cm}



\Method{此為股上容圓也。以勾股相乘，倍之為實，以勾股冪求弦，加入勾以為法。}



\vspace{0.2cm}



\Workings{置甲南行三百七十五步，以乙東行二百步乘之，得七萬五千步，倍之得一十五萬步為實。以乙東行自之，得四萬步為勾冪；以甲南行自之，得一十四萬零六百二十五步為股冪。勾股冪相並，得一十八萬零六百二十五步為方實。如平方而一，得四百二十五步則弦也。加入乙東行二百步，共得六百二十五步以為法。以法除之，得二百四十步，則城徑也。合問。}



\switchcolumn



\ar{سُئل: رجلان أ وب يقفان عند الباب الشمالي. يسير ب نحو الشرق 200 خطوة ويقف، ويسير أ نحو الجنوب 375 خطوة ويرى ب. الجواب كالسابق.}



\vspace{0.2cm}



\AnswerAr{قطر المدينة 240 خطوة.}



\vspace{0.2cm}



\MethodAr{هذه مسألة الدائرة المحاطة على الضلع. اضرب الضلعين القائمين واضعف الناتج فتكون الذات، واستخرج الوتر من مربّعي الضلعين، ثم أضف القاعدة فتكون القاسية.}



\vspace{0.2cm}



\WorkingsAr{ضع مسافة سير أ نحو الجنوب (375 خطوة) واضربها بمسافة سير ب نحو الشرق (200 خطوة)، تحصل على 75000 خطوة، واضعفها فتحصل على 150000 خطوة وهي الذات. ربّع مسافة سير ب نحو الشرق فتحصل على 40000 وهي مربّع القاعدة؛ ربّع مسافة سير أ نحو الجنوب فتحصل على 140625 وهي مربّع الضلع. اجمع مربّعي الضلعين فتحصل على 180625 وهي الذات المربّعة. استخرج الجذر التربيعي فتحصل على 425 خطوة وهي الوتر. أضف مسافة سير ب نحو الشرق (200 خطوة)، يصبح المجموع 625 خطوة وهو القاسية. اقسم الذات على القاسية فتحصل على 240 خطوة وهي قطر المدينة. يتطابق مع المطلوب.}



\end{paracol}



\vspace{0.5cm}



% ------------------------------------------------------------------------------

\subsubsection*{Problem 4 / \ar{المسألة الرابعة} \quad 第四問}



\begin{paracol}{2}



\noindent 或問：甲乙二人俱在圓城中心而立，乙穿城向東行一百三十六步而止，甲穿城南行二百五十五步望見乙，問答同前。



\vspace{0.2cm}



\Answer{城徑二百四十步。}



\vspace{0.2cm}



\Method{此為勾股上容圓也。以勾股相乘，倍之為實，並勾股冪如法求弦，以為法。}



\vspace{0.2cm}



\Workings{以二行步相乘，得三萬四千六百八十步，倍之得六萬九千三百六十步為實。置乙東行自之，得一萬八千四百九十六步為勾冪；又以甲南行自之，得六萬五千零二十五步為股冪。二冪相並，得八萬三千五百二十一步為方實，以平方開之，得二百八十九步即弦也。便以為法。如法除實，得二百四十步，即圓城之徑也。合問。}



\switchcolumn



\ar{سُئل: رجلان أ وب يقفان في مركز المدينة الدائرية. يعبر ب المدينة ويسير نحو الشرق 136 خطوة ويقف، ويعبر أ المدينة ويسير نحو الجنوب 255 خطوة ويرى ب. الجواب كالسابق.}



\vspace{0.2cm}



\AnswerAr{قطر المدينة 240 خطوة.}



\vspace{0.2cm}



\MethodAr{هذه مسألة الدائرة المحاطة على القاعدة والضلع. اضرب المسافتين واضعف الناتج فتكون الذات، واجمع مربّعي الضلعين لاستخراج الوتر فتكون القاسية.}



\vspace{0.2cm}



\WorkingsAr{اضرب المسافتين فتحصل على 34680 خطوة، واضعفها فتحصل على 69360 خطوة وهي الذات. ربّع مسافة سير ب نحو الشرق فتحصل على 18496 وهي مربّع القاعدة؛ ربّع مسافة سير أ نحو الجنوب فتحصل على 65025 وهي مربّع الضلع. اجمع المربّعين فتحصل على 83521 وهي الذات المربّعة، واستخرج الجذر التربيعي فتحصل على 289 خطوة وهي الوتر. وهي القاسية. اقسم الذات على القاسية فتحصل على 240 خطوة وهي قطر المدينة الدائرية. يتطابق مع المطلوب.}



\end{paracol}



\vspace{0.5cm}



% ------------------------------------------------------------------------------

\subsubsection*{Problem 5 / \ar{المسألة الخامسة} \quad 第五問}



\begin{paracol}{2}



\noindent 或問：甲乙二人同立於乾地，乙東行一百八十步遇塔而止，甲南行三百六十步，回望其塔正居城徑之半，問答同前。



\vspace{0.2cm}



\Answer{城徑二百四十步。}



\vspace{0.2cm}



\Method{此為弦上容圓也。以勾股相乘，倍之為實，以勾股和為法。}



\vspace{0.2cm}



\Workings{以二行步相乘，得六萬四千八百步，倍之得一十二萬九千六百步為實。並二行步，得五百四十步以為法。除實，得二百四十步，即城徑也。合問。}



\switchcolumn



\ar{سُئل: رجلان أ وب يقفان في موقع تشيين. يسير ب نحو الشرق 180 خطوة حتى يصل البرج ويقف، ويسير أ نحو الجنوب 360 خطوة، وعندما ينظر إلى الوراء يرى البرج عند منتصف قطر المدينة. الجواب كالسابق.}



\vspace{0.2cm}



\AnswerAr{قطر المدينة 240 خطوة.}



\vspace{0.2cm}



\MethodAr{هذه مسألة الدائرة المحاطة على الوتر. اضرب الضلعين القائمين واضعف الناتج فتكون الذات، واجمع الضلعين القائمين فتكون القاسية.}



\vspace{0.2cm}



\WorkingsAr{اضرب المسافتين فتحصل على 64800 خطوة، واضعفها فتحصل على 129600 خطوة وهي الذات. اجمع المسافتين فتحصل على 540 خطوة وهي القاسية. اقسم الذات على القاسية فتحصل على 240 خطوة وهي قطر المدينة. يتطابق مع المطلوب.}



\end{paracol}



\vspace{0.5cm}





% ------------------------------------------------------------------------------

\subsubsection*{Problem 6 / \ar{المسألة السادسة} \quad 第六問}



\begin{paracol}{2}



\noindent 或問：甲乙二人俱在坤地，乙東行一百九十二步而止，甲南行三百六十步，望乙與城參相直，問答同前。



\vspace{0.2cm}



\Answer{城徑二百四十步。}



\vspace{0.2cm}



\Method{此為勾外容圓也。以勾股相乘，倍之為實，以弦較和為法。}



\vspace{0.2cm}



\Workings{以二行步相乘，得六萬九千一百二十步，倍之得一十三萬八千二百四十步為實。置乙東行自之，得三萬六千八百六十四步為勾冪；又置甲南行自之，得一十二萬九千六百步為股冪。二冪相並，得一十六萬六千四百六十四步為方實，以平方開之，得四百零八，即弦也。又置甲南行步內減乙東行步，餘一百六十八步，即較也。以較加弦，共得五百七十六步以為法。實如法而一，得二百四十步，為城徑也。合問。}



\vspace{0.2cm}



\Remark{此題用勾股求得弦，即可加減得弦較較為城徑。今必以勾股相乘倍積為實，求得弦較和為法，而後始得弦較較為城徑者，蓋欲因此並明勾股相乘之倍積為弦較較、弦較和相乘之積，非故為紆回也。}



\switchcolumn



\ar{سُئل: رجلان أ وب يقفان في موقع كون. يسير ب نحو الشرق 192 خطوة ويقف، ويسير أ نحو الجنوب 360 خطوة، ويرى ب والمدينة على استقامة واحدة. الجواب كالسابق.}



\vspace{0.2cm}



\AnswerAr{قطر المدينة 240 خطوة.}



\vspace{0.2cm}



\MethodAr{هذه مسألة الدائرة المحاطة خارج القاعدة. اضرب الضلعين القائمين واضعف الناتج فتكون الذات، واجمع الوتر والفرق فتكون القاسية.}



\vspace{0.2cm}



\WorkingsAr{اضرب المسافتين فتحصل على 69120 خطوة، واضعفها فتحصل على 138240 خطوة وهي الذات. ربّع مسافة سير ب نحو الشرق فتحصل على 36864 وهي مربّع القاعدة؛ ربّع مسافة سير أ نحو الجنوب فتحصل على 129600 وهي مربّع الضلع. اجمع المربّعين فتحصل على 166464 وهي الذات المربّعة، واستخرج الجذر التربيعي فتحصل على 408 وهو الوتر. ثم اطرح مسافة سير ب نحو الشرق من مسافة سير أ نحو الجنوب، يتبقّى 168 خطوة وهي الفرق. أضف الفرق إلى الوتر، يصبح المجموع 576 خطوة وهو القاسية. اقسم الذات على القاسية فتحصل على 240 خطوة وهي قطر المدينة. يتطابق مع المطلوب.}



\vspace{0.2cm}



\RemarkAr{في هذه المسألة، يمكن استخراج الوتر من الضلعين القائمين ثم إجراء الجمع والطرح للحصول على قطر المدينة. غير أنّه يجب استخدام ضعف حاصل ضرب الضلعين القائمين كذات، ومجموع الوتر والفرق كقاسية، ثم الحصول على قطر المدينة. والغرض من ذلك هو توضيح أنّ ضعف حاصل ضرب الضلعين القائمين يساوي حاصل ضرب (الوتر ناقص الفرق) و(الوتر زائد الفرق)، وليس مجرد التعقيد.}



\end{paracol}



\vspace{0.5cm}



% ------------------------------------------------------------------------------

\subsubsection*{Problem 7 / \ar{المسألة السابعة} \quad 第七問}



\begin{paracol}{2}



\noindent 或問：甲乙二人同立於艮地，甲南行一百五十步而止，乙東行八十步，望乙與城參相直，問答同前。



\vspace{0.2cm}



\Answer{城徑二百四十步。}



\vspace{0.2cm}



\Method{此為股外容圓也。以勾股相乘，倍之為實，以弦和較為法。}



\switchcolumn



\ar{سُئل: رجلان أ وب يقفان في موقع كِن. يسير أ نحو الجنوب 150 خطوة ويقف، ويسير ب نحو الشرق 80 خطوة، ويرى ب والمدينة على استقامة واحدة. الجواب كالسابق.}



\vspace{0.2cm}



\AnswerAr{قطر المدينة 240 خطوة.}



\vspace{0.2cm}



\MethodAr{هذه مسألة الدائرة المحاطة خارج الضلع. اضرب الضلعين القائمين واضعف الناتج فتكون الذات، واجعل فرق مجموع الوتر والضلع القاسية.}



\end{paracol}



\vspace{0.5cm}



% ------------------------------------------------------------------------------

\subsubsection*{Problem 8 / \ar{المسألة الثامنة} \quad 第八問}



\begin{paracol}{2}



\noindent 或問：甲乙二人同立於坤地，甲南行一百九十二步而止，乙東行七十二步，望乙與城參相直，問答同前。



\vspace{0.2cm}



\Answer{城徑二百四十步。}



\vspace{0.2cm}



\Method{此為弦外容圓也。以勾股相乘，倍之為實，以勾股較為法。}



\switchcolumn



\ar{سُئل: رجلان أ وب يقفان في موقع كون. يسير أ نحو الجنوب 192 خطوة ويقف، ويسير ب نحو الشرق 72 خطوة، ويرى ب والمدينة على استقامة واحدة. الجواب كالسابق.}



\vspace{0.2cm}



\AnswerAr{قطر المدينة 240 خطوة.}



\vspace{0.2cm}



\MethodAr{هذه مسألة الدائرة المحاطة خارج الوتر. اضرب الضلعين القائمين واضعف الناتج فتكون الذات، واجعل فرق الضلعين القاسية.}



\end{paracol}



\vspace{0.5cm}



% ------------------------------------------------------------------------------

\subsubsection*{Problem 9 / \ar{المسألة التاسعة} \quad 第九問}



\begin{paracol}{2}



\noindent 或問：甲乙二人俱在巽地，乙東行一百九十二步而止，甲南行三百六十步，望乙與城參相直，問答同前。



\vspace{0.2cm}



\Answer{城徑二百四十步。}



\vspace{0.2cm}



\Method{此為勾外容半圓也。以勾股相乘，倍之為實，以股弦和為法。}



\switchcolumn



\ar{سُئل: رجلان أ وب يقفان في موقع شون. يسير ب نحو الشرق 192 خطوة ويقف، ويسير أ نحو الجنوب 360 خطوة، ويرى ب والمدينة على استقامة واحدة. الجواب كالسابق.}



\vspace{0.2cm}



\AnswerAr{قطر المدينة 240 خطوة.}



\vspace{0.2cm}



\MethodAr{هذه مسألة نصف الدائرة المحاطة خارج القاعدة. اضرب الضلعين القائمين واضعف الناتج فتكون الذات، واجمع الضلع والوتر فتكون القاسية.}



\end{paracol}



\vspace{0.5cm}



% ------------------------------------------------------------------------------

\subsubsection*{Problem 10 / \ar{المسألة العاشرة} \quad 第十問}



\begin{paracol}{2}



\noindent 或問：甲乙二人俱在艮地，甲南行一百五十步而止，乙東行八十步，望乙與城參相直，問答同前。



\vspace{0.2cm}



\Answer{城徑二百四十步。}



\vspace{0.2cm}



\Method{此為股外容半圓也。以勾股相乘，倍之為實，以勾弦和為法。}



\switchcolumn



\ar{سُئل: رجلان أ وب يقفان في موقع كِن. يسير أ نحو الجنوب 150 خطوة ويقف، ويسير ب نحو الشرق 80 خطوة، ويرى ب والمدينة على استقامة واحدة. الجواب كالسابق.}



\vspace{0.2cm}



\AnswerAr{قطر المدينة 240 خطوة.}



\vspace{0.2cm}



\MethodAr{هذه مسألة نصف الدائرة المحاطة خارج الضلع. اضرب الضلعين القائمين واضعف الناتج فتكون الذات، واجمع القاعدة والوتر فتكون القاسية.}



\end{paracol}



\vspace{0.5cm}



\noindent\textcolor{rulecolor}{\rule{\textwidth}{0.4pt}}



\vspace{0.3cm}



\begin{paracol}{2}

\textit{The above ten problems constitute the ``Dongyuan Nine Containments'' (洞淵九容), supplemented by Problem 5 (弦上容圓). The remaining four problems of Volume 2 continue with additional standard-rate problems using the tian yuan shu method.}



\switchcolumn



\ar{تُشكّل المسائل العشرة أعلاه «تسعة استيعابات دونغ يوان» (洞淵九容)، بالإضافة إلى المسألة الخامسة (الدائرة المحاطة على الوتر). وتستمر المسائل الأربع المتبقية في المجلد الثاني بمسائل إضافية في المعدّلات المعيارية باستخدام طريقة عنصر السماء.}



\end{paracol}



\newpage





% =============================================================================

\section{Volume 3 / \ar{المجلد الثالث} \quad 卷三}

\label{sec:vol3}



\begin{paracol}{2}



\textbf{邊股一十七問} --- Seventeen Questions on the Side-Long-Leg



Volume 3 presents 17 problems in which the given data involve the \textit{side long-leg} (邊股, $b_2$), defined as the long leg of the ``Side'' triangle (邊 triangle, No.\ 2 in the table), equal to 480 steps. In all problems of this volume, the unknown to be found is the diameter of the circular city (城徑 = 240 steps).



These problems demonstrate the increasing sophistication of the tian yuan shu method, with equations ranging from quadratic to cubic degree.



\switchcolumn



\ar{\textbf{سبع عشرة مسألة في الضلع الجانبي}}



\vspace{0.3cm}



\ar{يعرض المجلد الثالث 17 مسألة تتضمّن البيانات المعطاة الضلع الطويل الجانبي (邊股، $b_2$)، المعرَّف بأنه الضلع الطويل للمثلث الجانبي (邊، رقم 2 في الجدول)، وهو يساوي 480 خطوة. وفي جميع مسائل هذا المجلد، المجهول المطلوب إيجاده هو قطر المدينة الدائرية (城徑 = 240 خطوة).}



\vspace{0.3cm}



\ar{توضّح هذه المسائل التطوّر المتزايد في طريقة عنصر السماء، حيث تتراوح المعادلات من الدرجة الثانية إلى الدرجة الثالثة.}



\end{paracol}



% ------------------------------------------------------------------------------

\subsubsection*{Problem 1 / \ar{المسألة الأولى} \quad 第一問}



\begin{paracol}{2}



\noindent 或問：乙出東門南行，不知步數而止。甲出西門南行四百八十步，望見乙，復就乙行五百一十步與乙相會，問答同前。



\vspace{0.2cm}



\Answer{城徑二百四十步。}



\vspace{0.2cm}



\Method{倍相減步，以乘二之甲南行步，為平方實，得城徑。}



\vspace{0.2cm}



\Workings{識別得二行相減，餘三十步，即乙出東門南行步也。倍相減步，得六十步，以乘二之甲南行步九百六十步，得五萬七千六百步，為平方實。如法開之，得二百四十步，即城徑也。合問。}



\switchcolumn



\ar{سُئل: يخرج ب من الباب الشرقي ويسير نحو الجنوب مسافة مجهولة ويقف. يخرج أ من الباب الغربي ويسير نحو الجنوب 480 خطوة، ويرى ب، ثم يسير نحو ب 510 خطوة حتى يلتقي به. الجواب كالسابق.}



\vspace{0.2cm}



\AnswerAr{قطر المدينة 240 خطوة.}



\vspace{0.2cm}



\MethodAr{اضعف الفرق بين المسافتين، واضربه بضعف مسافة سير أ نحو الجنوب، فتكون ذات التربيع، وتحصل على قطر المدينة.}



\vspace{0.2cm}



\WorkingsAr{بالتمييز: نطرح المسافتين، يتبقّى 30 خطوة، وهي مسافة سير ب من الباب الشرقي نحو الجنوب. نضعف الفرق فنحصل على 60 خطوة، ونضربها بضعف مسافة سير أ نحو الجنوب (960 خطوة)، فنحصل على 57600 خطوة وهي ذات التربيع. نستخرج الجذر فنحصل على 240 خطوة وهي قطر المدينة. يتطابق مع المطلوب.}



\end{paracol}



\vspace{0.5cm}



% ------------------------------------------------------------------------------

\subsubsection*{Problem 2 / \ar{المسألة الثانية} \quad 第二問}



\begin{paracol}{2}



\noindent 或問：甲出西門南行四百八十步而止，乙從艮隅東行八十步，望見甲，問答同前。



\vspace{0.2cm}



\Answer{城徑二百四十步。}



\vspace{0.2cm}



\Method{倍南行步，以東行步乘之為實，東行步為從方，一步常法，得全徑。}



\vspace{0.2cm}



\Workings{立天元一為全徑，以減於二之甲南行步，得為兩個大差也。以乙東行步乘之，得為圓徑冪，寄左。然後以天元冪與左相消，得以帶縱平方開之，得二百四十步，即城徑也。合問。}



\vspace{0.2cm}



\noindent\textbf{又法：}半之乙東行步，乘南行步為實，半之乙東行步為從，一步常法，得半徑。



\vspace{0.2cm}



\Workings{立天元一為半城徑，減甲南行步，得為大差也。以半之東行步乘之，得即半徑冪，寄左。然後以天元冪為同數，與左相消，得開帶縱平方，得一百二十步。倍之，即城徑也。合問。}



\switchcolumn



\ar{سُئل: يخرج أ من الباب الغربي ويسير نحو الجنوب 480 خطوة ويقف. يسير ب من زاوية كِن نحو الشرق 80 خطوة، ويرى أ. الجواب كالسابق.}



\vspace{0.2cm}



\AnswerAr{قطر المدينة 240 خطوة.}



\vspace{0.2cm}



\MethodAr{اضعف مسافة السير نحو الجنوب، واضربها بمسافة السير نحو الشرق فتكون الذات، واجعل مسافة السير نحو الشرق الطرف المرافق، وواحدًا القاسية الثابتة، فتحصل على القطر الكامل.}



\vspace{0.2cm}



\WorkingsAr{نضع عنصر السماء يساوي القطر الكامل، ونطرحه من ضعف مسافة سير أ نحو الجنوب، فنحصل على ضعف الفرق الكبير. نضربها بمسافة سير ب نحو الشرق فنحصل على مربّع قطر الدائرة، ونضعها على اليسار. ثم نربّع عنصر السماء ونلغِه مع الطرف الأيسر، فنستخرج الجذر التربيعي ذا الطرف فنحصل على 240 خطوة وهي قطر المدينة. يتطابق مع المطلوب.}



\vspace{0.2cm}



\ar{\textbf{طريقة ثانية:} نصف مسافة السير نحو الشرق مضروبًا بمسافة السير نحو الجنوب فتكون الذات، ونصف مسافة السير نحو الشرق هو الطرف المرافق، وواحدًا القاسية الثابتة، فتحصل على نصف القطر.}



\vspace{0.2cm}



\WorkingsAr{نضع عنصر السماء يساوي نصف قطر المدينة، ونطرحه من مسافة سير أ نحو الجنوب فنحصل على الفرق الكبير. نضربها بنصف مسافة السير نحو الشرق فنحصل على مربّع نصف القطر، ونضعها على اليسار. ثم نجعل مربّع عنصر السماء مساويًا لها، ونلغِه مع الطرف الأيسر، فنستخرج الجذر التربيعي ذا الطرف فنحصل على 120 خطوة. نضعفها فنحصل على قطر المدينة. يتطابق مع المطلوب.}



\end{paracol}



\vspace{0.5cm}



% ------------------------------------------------------------------------------

\subsubsection*{Problem 3 / \ar{المسألة الثالثة} \quad 第三問}



\begin{paracol}{2}



\noindent 或問：甲出西門南行四百八十步而止，乙從艮隅亦南行一百五十步，望見甲，問答同前。



\vspace{0.2cm}



\Answer{城徑二百四十步。}



\vspace{0.2cm}



\Method{兩行步相乘為實，南行步為從方，一為隅，得半徑。}



\vspace{0.2cm}



\Workings{立天元一為半城徑，以減乙南行步，得為半梯頭；以甲行步為梯底，以乘之，得為半徑冪，寄左。然後以天元冪與左相消，得開帶縱平方，得一百二十步。倍之，即城徑也。合問。}



\switchcolumn



\ar{سُئل: يخرج أ من الباب الغربي ويسير نحو الجنوب 480 خطوة ويقف. يسير ب من زاوية كِن أيضًا نحو الجنوب 150 خطوة، ويرى أ. الجواب كالسابق.}



\vspace{0.2cm}



\AnswerAr{قطر المدينة 240 خطوة.}



\vspace{0.2cm}



\MethodAr{اضرب المسافتين فتكون الذات، واجعل مسافة السير نحو الجنوب الطرف المرافق، والواحد الركن، فتحصل على نصف القطر.}



\vspace{0.2cm}



\WorkingsAr{نضع عنصر السماء يساوي نصف قطر المدينة، ونطرحه من مسافة سير ب نحو الجنوب فنحصل على نصف رأس السلّم. نجعل مسافة سير أ الطول السفلي للسلّم، ونضربها فنحصل على مربّع نصف القطر، ونضعها على اليسار. ثم نربّع عنصر السماء ونلغِه مع الطرف الأيسر، فنستخرج الجذر التربيعي ذا الطرف فنحصل على 120 خطوة. نضعفها فنحصل على قطر المدينة. يتطابق مع المطلوب.}



\end{paracol}



\vspace{0.5cm}



% ------------------------------------------------------------------------------

\subsubsection*{Problem 4 / \ar{المسألة الرابعة} \quad 第四問}



\begin{paracol}{2}



\noindent 或問：甲出西門南行四百八十步，乙出東門直行一十六步，望見甲，問答同前。



\vspace{0.2cm}



\Answer{城徑二百四十步。}



\vspace{0.2cm}



\Method{以四之東行步，乘南行冪為實，從空，東行為廉，一步為隅法，得全徑。}



\vspace{0.2cm}



\Workings{立天元一為圓徑，加乙東行步，得為中勾；其甲南行即中股也。置東行步為小勾，以中股乘之，得，合以中勾除，今不受除，便以為小股也。內寄中勾分母。乃復以中股乘之，得三百六十八萬六千四百，又四之，得一千四百七十四萬五千六百，為一段圓徑冪，寄中勾分母，寄左。然後以天元徑自之，又以中勾乘之，得為同數，與左相消，得以帶縱立方開之，得二百四十步，為城徑也。合問。}



\vspace{0.2cm}



\Remark{不受除者，無可除之理也。凡二數，此數於彼數有可除之理，則受除；無可除之理，則不受除也。蓋除有法有實，實可二，法不可二。此題以中勾為法，而中勾內有一元，又有十六步，其為數已二矣，又何以均分不一之數乎？故曰不受也。寄分者，姑寄其應除之數也。俟求得兩相等數，而此數內尚少一除，不除此而轉乘彼，则兩數仍相等，猶之受除者也。此所謂以乘代除也。}



\switchcolumn



\ar{سُئل: يخرج أ من الباب الغربي ويسير نحو الجنوب 480 خطوة، ويخرج ب من الباب الشرقي ويسير مباشرة 16 خطوة، ويرى أ. الجواب كالسابق.}



\vspace{0.2cm}



\AnswerAr{قطر المدينة 240 خطوة.}



\vspace{0.2cm}



\MethodAr{أربعة أضعاف مسافة السير نحو الشرق مضروبة بمربّع مسافة السير نحو الجنوب فتكون الذات، الطرف فارغ، مسافة السير نحو الشرق هي الكشك، وواحدًا قاسية الركن، فتحصل على القطر الكامل.}



\vspace{0.2cm}



\WorkingsAr{نضع عنصر السماء يساوي قطر الدائرة، ونضيف إليه مسافة سير ب نحو الشرق فنحصل على القاعدة الوسطى؛ ومسافة سير أ نحو الجنوب هي الضلع الأوسط. نضع مسافة السير نحو الشرق كقاعدة صغرى، ونضربها بالضلع الأوسط، فنحصل على... يجب تقسيمها على القاعدة الوسطى، لكنّها لا تقبل القسمة، فنعتبرها الضلع الصغير مؤقتًا، مع إبقاء القاعدة الوسطى مقامًا. ثم نضربها مرة أخرى بالضلع الأوسط فنحصل على 3686400، وأربعة أضعافها هي 14745600، وهي مربّع قطر الدائرة، مع إبقاء القاعدة الوسطى مقامًا، ونضعها على اليسار. ثم نربّع عنصر السماء (القطر) ونضربها بالقاعدة الوسطى فنحصل على المساوي، ونلغِه مع الطرف الأيسر، فنستخرج الجذر التكعبي ذا الطرف فنحصل على 240 خطوة وهي قطر المدينة. يتطابق مع المطلوب.}



\vspace{0.2cm}



\RemarkAr{القول «لا تقبل القسمة» يعني أنّ القسمة غير ممكنة. ففي حال وجود عددين، إذا كان أحدهما يقبل القسمة على الآخر فهو «يقبل القسمة»، وإذا لم يكن يقبل فهو «لا يقبل القسمة». ذلك أنّ القسمة تقتضي مقسومًا ومقسومًا عليه، والمقسوم يمكن أن يكون ثنائيًا أمّا المقسوم عليه فلا. ففي هذه المسألة، القاعدة الوسطى هي المقسوم عليها، وهي تحتوي على مجهول وعلى 16 خطوة، فهي بالفعل عدد ثنائي، فكيف نقسم عددًا غير موحّد؟ لهذا نقول إنّها لا تقبل القسمة. أمّا «الإبقاء ككسور» فهو تأجيل العدد المفترض قسمته. حتى إذا حصلنا على عددين متساويين، وكان أحدهما لا يزال يفتقر إلى قسمة واحدة، فبدل أن نقسم نضرب الطرف الآخر، فيبقيان متساويين، كأنّهما قُسما. وهذا ما يُسمّى «الاستعاضة بالضرب عن القسمة».}



\end{paracol}



\vspace{0.5cm}



% ------------------------------------------------------------------------------

\subsubsection*{Problem 5 / \ar{المسألة الخامسة} \quad 第五問}



\begin{paracol}{2}



\noindent 或問：乙出南門東行七十二步而止，甲出西門南行四百八十步，望乙與城參相直，問答同前。



\vspace{0.2cm}



\Answer{城徑二百四十步。}



\vspace{0.2cm}



\Method{以乙東行冪，乘甲南行為實；乙東行冪為從方；甲南行步內減二之東行步為益廉；一步常法，得半徑。}



\vspace{0.2cm}



\Workings{立天元一為半城徑，以減南行步，得為小股。又以天元加乙東行步，得為小勾。又以天元加南行步，得為大股。乃置大股在地，以小勾乘之，得下式，合以小股除之，今不受除，便以為大勾，內寄小股分母。又置天元半徑，以分母小股乘之，得以減大勾，得為半個梯底於上。以乙東行七十二步為半個梯頭，以乘上位，得為半徑冪，內寄小股分母，寄左。然後置天元冪，又以分母小股乘之，得為同數，與左相消，以立方開之，得一百二十步。倍之，即城徑也。合問。}



\switchcolumn



\ar{سُئل: يخرج ب من الباب الجنوبي ويسير نحو الشرق 72 خطوة ويقف. يخرج أ من الباب الغربي ويسير نحو الجنوب 480 خطوة، ويرى ب والمدينة على استقامة واحدة. الجواب كالسابق.}



\vspace{0.2cm}



\AnswerAr{قطر المدينة 240 خطوة.}



\vspace{0.2cm}



\MethodAr{مربّع مسافة سير ب نحو الشرق مضروبًا بمسافة سير أ نحو الجنوب فتكون الذات؛ مربّع مسافة سير ب نحو الشرق هو الطرف المرافق؛ مسافة سير أ نحو الجنوب ناقص ضعف مسافة السير نحو الشرق هو الكشك الموجب؛ وواحدًا القاسية الثابتة، فتحصل على نصف القطر.}



\vspace{0.2cm}



\WorkingsAr{نضع عنصر السماء يساوي نصف قطر المدينة، ونطرحه من مسافة السير نحو الجنوب فنحصل على الضلع الصغير. ونضيف عنصر السماء إلى مسافة سير ب نحو الشرق فنحصل على القاعدة الصغرى. ونضيف عنصر السماء إلى مسافة السير نحو الجنوب فنحصل على الضلع الأكبر. نضع الضلع الأكبر ونضربه بالقاعدة الصغرى فنحصل على الناتج التالي، ويجب تقسيمه على الضلع الصغير، لكنّه لا يقبل القسمة، فنعتبره القاعدة الكبرى مؤقتًا مع إبقاء الضلع الصغير مقامًا. ثم نضع نصف قطر عنصر السماء ونضربه بالمقام (الضلع الصغير)، ونطرحه من القاعدة الكبرى فنحصل على نصف قاع السلّم في الأعلى. ومسافة سير ب نحو الشرق (72 خطوة) هي نصف رأس السلّم، ونضربها بالعدد الأعلى فنحصل على مربّع نصف القطر مع إبقاء الضلع الصغير مقامًا، ونضعها على اليسار. ثم نربّع عنصر السماء ونضربها بالمقام (الضلع الصغير) فنحصل على المساوي، ونلغِه مع الطرف الأيسر، فنستخرج الجذر التكعبي فنحصل على 120 خطوة. نضعفها فنحصل على قطر المدينة. يتطابق مع المطلوب.}



\end{paracol}



\vspace{0.5cm}





% ------------------------------------------------------------------------------

\subsubsection*{Problem 6 / \ar{المسألة السادسة} \quad 第六問}



\begin{paracol}{2}



\noindent 或問：乙出南門東行，不知步數而立。甲出西門南行四百八十步，望見乙與城參相直，又就乙行四百零八步與乙相會，問答同前。



\vspace{0.2cm}



\Answer{城徑二百四十步。}



\vspace{0.2cm}



\Method{以兩行相乘為實，以兩行相減為從，一步常法，平開得半徑。}



\vspace{0.2cm}



\Workings{立天元一為半徑。以加就乙行步，得為大勾；以減甲南行步，得為大股。乃以大勾、大股相乘，得下式，合以天元加就乙行除之，今不受除，便以為小股，內寄分母。又以天元乘之，得為半徑冪，寄左。然後以天元冪與左相消，得平開，得一百二十步。倍之，即城徑也。合問。}



\switchcolumn



\ar{سُئل: يخرج ب من الباب الجنوبي ويسير نحو الشرق مسافة مجهولة ويقف. يخرج أ من الباب الغربي ويسير نحو الجنوب 480 خطوة، ويرى ب والمدينة على استقامة واحدة، ثم يسير نحو ب 408 خطوات حتى يلتقي به. الجواب كالسابق.}



\vspace{0.2cm}



\AnswerAr{قطر المدينة 240 خطوة.}



\vspace{0.2cm}



\MethodAr{اضرب المسافتين فتكون الذات، واجعل فرق المسافتين هو الطرف المرافق، وواحدًا القاسية الثابتة، فاستخرج الجذر التربيعي فتحصل على نصف القطر.}



\vspace{0.2cm}



\WorkingsAr{نضع عنصر السماء يساوي نصف القطر. نضيف إليه مسافة السير نحو ب فنحصل على القاعدة الكبرى؛ ونطرحه من مسافة سير أ نحو الجنوب فنحصل على الضلع الأكبر. نضرب القاعدة الكبرى بالضلع الأكبر فنحصل على الناتج التالي، ويجب تقسيمه على مجموع عنصر السماء ومسافة السير نحو ب، لكنّه لا يقبل القسمة، فنعتبره الضلع الصغير مؤقتًا مع إبقاء المقام. ثم نضربها بعنصر السماء فنحصل على مربّع نصف القطر، ونضعها على اليسار. ثم نربّع عنصر السماء ونلغِه مع الطرف الأيسر، فنستخرج الجذر التربيعي فنحصل على 120 خطوة. نضعفها فنحصل على قطر المدينة. يتطابق مع المطلوب.}



\end{paracol}



\vspace{0.5cm}



% ------------------------------------------------------------------------------

\subsubsection*{Problem 7 / \ar{المسألة السابعة} \quad 第七問}



\begin{paracol}{2}



\noindent 或問：甲出西門南行四百八十步而止，乙出東門南行三十步，望見甲，又斜行二百五十步與甲相會，問答同前。



\vspace{0.2cm}



\Answer{城徑二百四十步。}



\vspace{0.2cm}



\Method{兩行步相減，以乘兩行步相併為實，兩行步相減又自之為從方，兩倍兩行步相減為益廉，一步常法，得半徑。}



\vspace{0.2cm}



\Workings{立天元一為半徑。以減乙南行步，得為梯頭；以甲南行步為梯底。以梯頭乘梯底，得為半徑冪，寄左。乃以兩行步相減，餘四百五十步；又以兩行步相併，得五百一十步；以減餘乘之，得二十二萬九千五百步為實。又以減餘四百五十步自之，得二十萬零二千五百步為從方。又以減餘倍之，得九百步為益廉。以立方開之，得一百二十步。倍之，即城徑也。合問。}



\switchcolumn



\ar{سُئل: يخرج أ من الباب الغربي ويسير نحو الجنوب 480 خطوة ويقف. يخرج ب من الباب الشرقي ويسير نحو الجنوب 30 خطوة، ويرى أ، ثم يسير مائلًا 250 خطوة حتى يلتقي بأ. الجواب كالسابق.}



\vspace{0.2cm}



\AnswerAr{قطر المدينة 240 خطوة.}



\vspace{0.2cm}



\MethodAr{اطرح المسافتين، واضرب الناتج بمجموع المسافتين فتكون الذات. اطرح المسافتين وربّع الناتج فتكون الطرف المرافق. اضعف فرق المسافتين فتكون الكشك الموجب. وواحدًا القاسية الثابتة، فتحصل على نصف القطر.}



\vspace{0.2cm}



\WorkingsAr{نضع عنصر السماء يساوي نصف القطر. نطرحه من مسافة سير ب نحو الجنوب فنحصل على رأس السلّم؛ ومسافة سير أ نحو الجنوب هي قاع السلّم. نضرب رأس السلّم بقاعه فنحصل على مربّع نصف القطر، ونضعها على اليسار. نطرح المسافتين فيتبقّى 450 خطوة؛ ونجمع المسافتين فنحصل على 510 خطوات؛ نضرب الباقي بالمجموع فنحصل على 229500 خطوة وهي الذات. نربّع الباقي (450 خطوة) فنحصل على 202500 وهي الطرف المرافق. ونضعف الباقي فنحصل على 900 خطوة وهي الكشك الموجب. نستخرج الجذر التكعبي فنحصل على 120 خطوة. نضعفها فنحصل على قطر المدينة. يتطابق مع المطلوب.}



\end{paracol}



\vspace{0.5cm}



% ------------------------------------------------------------------------------

\subsubsection*{Problem 8 / \ar{المسألة الثامنة} \quad 第八問}



\begin{paracol}{2}



\noindent 或問：甲出西門南行四百八十步而止，乙出東門南行五十四步而止，望見甲，又斜行三百三十步與甲相會，問答同前。



\vspace{0.2cm}



\Answer{城徑二百四十步。}



\vspace{0.2cm}



\Method{兩行相減，以乘兩行相併為實，兩行相減自之為從方，四之兩行相減為益廉，一步常法，得半徑。}



\vspace{0.2cm}



\Workings{立天元一為半徑。以減乙南行步，得為梯頭差；以甲南行步內減天元，得為梯底差。以梯頭差乘梯底差，得為半徑冪，寄左。乃以兩行相減，餘四百二十六步；又與兩行相併，得五百三十四步；以減餘乘之，得二十二萬七千四百八十四步為實。又以減餘自之，得一十八萬一千四百七十六步為從方。又以減餘四之，得一千七百零四步為益廉。以立方開之，得一百二十步。倍之，即城徑也。合問。}



\switchcolumn



\ar{سُئل: يخرج أ من الباب الغربي ويسير نحو الجنوب 480 خطوة ويقف. يخرج ب من الباب الشرقي ويسير نحو الجنوب 54 خطوة ويقف، ويرى أ، ثم يسير مائلًا 330 خطوة حتى يلتقي بأ. الجواب كالسابق.}



\vspace{0.2cm}



\AnswerAr{قطر المدينة 240 خطوة.}



\vspace{0.2cm}



\MethodAr{اطرح المسافتين، واضرب الناتج بمجموع المسافتين فتكون الذات. اطرح المسافتين وربّع الناتج فتكون الطرف المرافق. أربعة أضعاف فرق المسافتين فتكون الكشك الموجب. وواحدًا القاسية الثابتة، فتحصل على نصف القطر.}



\vspace{0.2cm}



\WorkingsAr{نضع عنصر السماء يساوي نصف القطر. نطرحه من مسافة سير ب نحو الجنوب فنحصل على فرق رأس السلّم؛ ونطرح عنصر السماء من مسافة سير أ نحو الجنوب فنحصل على فرق قاع السلّم. نضرب فرق رأس السلّم بفرق قاعه فنحصل على مربّع نصف القطر، ونضعها على اليسار. نطرح المسافتين فيتبقّى 426 خطوة؛ ونجمع المسافتين فنحصل على 534 خطوة؛ نضرب الباقي بالمجموع فنحصل على 227484 خطوة وهي الذات. نربّع الباقي فنحصل على 181476 وهي الطرف المرافق. وأربعة أضعاف الباقي هي 1704 خطوات وهي الكشك الموجب. نستخرج الجذر التكعبي فنحصل على 120 خطوة. نضعفها فنحصل على قطر المدينة. يتطابق مع المطلوب.}



\end{paracol}



\vspace{0.5cm}



% ------------------------------------------------------------------------------

\subsubsection*{Problem 9 / \ar{المسألة التاسعة} \quad 第九問}



\begin{paracol}{2}



\noindent 或問：甲出西門南行四百八十步而止，乙出東門南行九十步而止，望見甲，又斜行三百九十步與甲相會，問答同前。



\vspace{0.2cm}



\Answer{城徑二百四十步。}



\vspace{0.2cm}



\Method{兩行相減，以乘兩行相併為實，兩行相減自之為從方，兩倍兩行相減又益四之東行為益廉，一步常法，得半徑。}



\vspace{0.2cm}



\Workings{立天元一為半徑。以減甲南行步，得為梯底差；以乙南行步內減天元，得為梯頭差。以梯頭差乘梯底差，得為半徑冪，寄左。乃以兩行相減，餘三百九十步；又以兩行相併，得五百七十步；以減餘乘之，得二十二萬二千三百步為實。又以減餘自之，得一十五萬二千一百步為從方。又以減餘倍之，得七百八十步；又益四之乙東行九十步，得三百六十步；共得一千一百四十步為益廉。以立方開之，得一百二十步。倍之，即城徑也。合問。}



\switchcolumn



\ar{سُئل: يخرج أ من الباب الغربي ويسير نحو الجنوب 480 خطوة ويقف. يخرج ب من الباب الشرقي ويسير نحو الجنوب 90 خطوة ويقف، ويرى أ، ثم يسير مائلًا 390 خطوة حتى يلتقي بأ. الجواب كالسابق.}



\vspace{0.2cm}



\AnswerAr{قطر المدينة 240 خطوة.}



\vspace{0.2cm}



\MethodAr{اطرح المسافتين، واضرب الناتج بمجموع المسافتين فتكون الذات. اطرح المسافتين وربّع الناتج فتكون الطرف المرافق. ضعف فرق المسافتين مضافًا إلى أربعة أضعاف مسافة السير نحو الشرق فتكون الكشك الموجب. وواحدًا القاسية الثابتة، فتحصل على نصف القطر.}



\vspace{0.2cm}



\WorkingsAr{نضع عنصر السماء يساوي نصف القطر. نطرحه من مسافة سير أ نحو الجنوب فنحصل على فرق قاع السلّم؛ ونطرح عنصر السماء من مسافة سير ب نحو الجنوب فنحصل على فرق رأس السلّم. نضرب فرق رأس السلّم بفرق قاعه فنحصل على مربّع نصف القطر، ونضعها على اليسار. نطرح المسافتين فيتبقّى 390 خطوة؛ ونجمع المسافتين فنحصل على 570 خطوة؛ نضرب الباقي بالمجموع فنحصل على 222300 خطوة وهي الذات. نربّع الباقي فنحصل على 152100 وهي الطرف المرافق. وضعف الباقي هو 780 خطوة؛ مضافًا إلى أربعة أضعاف مسافة سير ب نحو الشرق (90 خطوة) فتصبح 360؛ والمجموع 1140 خطوة وهي الكشك الموجب. نستخرج الجذر التكعبي فنحصل على 120 خطوة. نضعفها فنحصل على قطر المدينة. يتطابق مع المطلوب.}



\end{paracol}



\vspace{0.5cm}



% ------------------------------------------------------------------------------

\subsubsection*{Problem 10 / \ar{المسألة العاشرة} \quad 第十問}



\begin{paracol}{2}



\noindent 或問：甲出西門南行四百八十步而止，乙出東門南行一百二十步而止，望見甲，又斜行四百五十步與甲相會，問答同前。



\vspace{0.2cm}



\Answer{城徑二百四十步。}



\vspace{0.2cm}



\Method{兩行相減，以乘兩行相併為實，兩行相減自之為從方，兩倍兩行相減又益四之東行為益廉，一步常法，得半徑。}



\vspace{0.2cm}



\Workings{立天元一為半徑。以減甲南行步，得為梯底差；以乙南行步內減天元，得為梯頭差。以梯頭差乘梯底差，得為半徑冪，寄左。乃以兩行相減，餘三百六十步；又以兩行相併，得六百步；以減餘乘之，得二十一万六千步為實。又以減餘自之，得一十二萬九千六百步為從方。又以減餘倍之，得七百二十步；又益四之乙東行一百二十步，得四百八十步；共得一千二百步為益廉。以立方開之，得一百二十步。倍之，即城徑也。合問。}



\switchcolumn



\ar{سُئل: يخرج أ من الباب الغربي ويسير نحو الجنوب 480 خطوة ويقف. يخرج ب من الباب الشرقي ويسير نحو الجنوب 120 خطوة ويقف، ويرى أ، ثم يسير مائلًا 450 خطوة حتى يلتقي بأ. الجواب كالسابق.}



\vspace{0.2cm}



\AnswerAr{قطر المدينة 240 خطوة.}



\vspace{0.2cm}



\MethodAr{اطرح المسافتين، واضرب الناتج بمجموع المسافتين فتكون الذات. اطرح المسافتين وربّع الناتج فتكون الطرف المرافق. ضعف فرق المسافتين مضافًا إلى أربعة أضعاف مسافة السير نحو الشرق فتكون الكشك الموجب. وواحدًا القاسية الثابتة، فتحصل على نصف القطر.}



\vspace{0.2cm}



\WorkingsAr{نضع عنصر السماء يساوي نصف القطر. نطرحه من مسافة سير أ نحو الجنوب فنحصل على فرق قاع السلّم؛ ونطرح عنصر السماء من مسافة سير ب نحو الجنوب فنحصل على فرق رأس السلّم. نضرب فرق رأس السلّم بفرق قاعه فنحصل على مربّع نصف القطر، ونضعها على اليسار. نطرح المسافتين فيتبقّى 360 خطوة؛ ونجمع المسافتين فنحصل على 600 خطوة؛ نضرب الباقي بالمجموع فنحصل على 216000 خطوة وهي الذات. نربّع الباقي فنحصل على 129600 وهي الطرف المرافق. وضعف الباقي هو 720 خطوة؛ مضافًا إلى أربعة أضعاف مسافة سير ب نحو الشرق (120 خطوة) فتصبح 480؛ والمجموع 1200 خطوة وهي الكشك الموجب. نستخرج الجذر التكعبي فنحصل على 120 خطوة. نضعفها فنحصل على قطر المدينة. يتطابق مع المطلوب.}



\end{paracol}



\vspace{0.5cm}





% ------------------------------------------------------------------------------

\subsubsection*{Problem 11 / \ar{المسألة الحادية عشرة} \quad 第十一問}



\begin{paracol}{2}



\noindent 或問：甲出西門南行四百八十步而止，乙出東門南行一百八十步而止，望見甲，又斜行五百一十步與甲相會，問答同前。



\vspace{0.2cm}



\Answer{城徑二百四十步。}



\vspace{0.2cm}



\Method{兩行相減，以乘兩行相併為實，兩行相減自之為從方，兩倍兩行相減又益四之東行為益廉，一步常法，得半徑。}



\vspace{0.2cm}



\Workings{立天元一為半徑。以減甲南行步，得為梯底差；以乙南行步內減天元，得為梯頭差。以梯頭差乘梯底差，得為半徑冪，寄左。乃以兩行相減，餘三百步；又以兩行相併，得六百六十步；以減餘乘之，得十九萬八千步為實。又以減餘自之，得九萬步為從方。又以減餘倍之，得六百步；又益四之乙東行一百八十步，得七百二十步；共得一千三百二十步為益廉。以立方開之，得一百二十步。倍之，即城徑也。合問。}



\switchcolumn



\ar{سُئل: يخرج أ من الباب الغربي ويسير نحو الجنوب 480 خطوة ويقف. يخرج ب من الباب الشرقي ويسير نحو الجنوب 180 خطوة ويقف، ويرى أ، ثم يسير مائلًا 510 خطوات حتى يلتقي بأ. الجواب كالسابق.}



\vspace{0.2cm}



\AnswerAr{قطر المدينة 240 خطوة.}



\vspace{0.2cm}



\MethodAr{اطرح المسافتين، واضرب الناتج بمجموع المسافتين فتكون الذات. اطرح المسافتين وربّع الناتج فتكون الطرف المرافق. ضعف فرق المسافتين مضافًا إلى أربعة أضعاف مسافة السير نحو الشرق فتكون الكشك الموجب. وواحدًا القاسية الثابتة، فتحصل على نصف القطر.}



\vspace{0.2cm}



\WorkingsAr{نضع عنصر السماء يساوي نصف القطر. نطرحه من مسافة سير أ نحو الجنوب فنحصل على فرق قاع السلّم؛ ونطرح عنصر السماء من مسافة سير ب نحو الجنوب فنحصل على فرق رأس السلّم. نضرب فرق رأس السلّم بفرق قاعه فنحصل على مربّع نصف القطر، ونضعها على اليسار. نطرح المسافتين فيتبقّى 300 خطوة؛ ونجمع المسافتين فنحصل على 660 خطوة؛ نضرب الباقي بالمجموع فنحصل على 198000 خطوة وهي الذات. نربّع الباقي فنحصل على 90000 وهي الطرف المرافق. وضعف الباقي هو 600 خطوة؛ مضافًا إلى أربعة أضعاف مسافة سير ب نحو الشرق (180 خطوة) فتصبح 720؛ والمجموع 1320 خطوة وهي الكشك الموجب. نستخرج الجذر التكعبي فنحصل على 120 خطوة. نضعفها فنحصل على قطر المدينة. يتطابق مع المطلوب.}



\end{paracol}



\vspace{0.5cm}



% ------------------------------------------------------------------------------

\subsubsection*{Problem 12 / \ar{المسألة الثانية عشرة} \quad 第十二問}



\begin{paracol}{2}



\noindent 或問：甲出西門南行四百八十步而止，乙出東門南行二百一十步而止，望見甲，又斜行五百七十步與甲相會，問答同前。



\vspace{0.2cm}



\Answer{城徑二百四十步。}



\vspace{0.2cm}



\Method{兩行相減，以乘兩行相併為實，兩行相減自之為從方，兩倍兩行相減又益四之東行為益廉，一步常法，得半徑。}



\vspace{0.2cm}



\Workings{立天元一為半徑。以減甲南行步，得為梯底差；以乙南行步內減天元，得為梯頭差。以梯頭差乘梯底差，得為半徑冪，寄左。乃以兩行相減，餘二百七十步；又以兩行相併，得六百九十步；以減餘乘之，得十八萬六千三百步為實。又以減餘自之，得七萬二千九百步為從方。又以減餘倍之，得五百四十步；又益四之乙東行二百一十步，得八百四十步；共得一千三百八十步為益廉。以立方開之，得一百二十步。倍之，即城徑也。合問。}



\switchcolumn



\ar{سُئل: يخرج أ من الباب الغربي ويسير نحو الجنوب 480 خطوة ويقف. يخرج ب من الباب الشرقي ويسير نحو الجنوب 210 خطوات ويقف، ويرى أ، ثم يسير مائلًا 570 خطوة حتى يلتقي بأ. الجواب كالسابق.}



\vspace{0.2cm}



\AnswerAr{قطر المدينة 240 خطوة.}



\vspace{0.2cm}



\MethodAr{اطرح المسافتين، واضرب الناتج بمجموع المسافتين فتكون الذات. اطرح المسافتين وربّع الناتج فتكون الطرف المرافق. ضعف فرق المسافتين مضافًا إلى أربعة أضعاف مسافة السير نحو الشرق فتكون الكشك الموجب. وواحدًا القاسية الثابتة، فتحصل على نصف القطر.}



\vspace{0.2cm}



\WorkingsAr{نضع عنصر السماء يساوي نصف القطر. نطرحه من مسافة سير أ نحو الجنوب فنحصل على فرق قاع السلّم؛ ونطرح عنصر السماء من مسافة سير ب نحو الجنوب فنحصل على فرق رأس السلّم. نضرب فرق رأس السلّم بفرق قاعه فنحصل على مربّع نصف القطر، ونضعها على اليسار. نطرح المسافتين فيتبقّى 270 خطوة؛ ونجمع المسافتين فنحصل على 690 خطوة؛ نضرب الباقي بالمجموع فنحصل على 186300 خطوة وهي الذات. نربّع الباقي فنحصل على 72900 وهي الطرف المرافق. وضعف الباقي هو 540 خطوة؛ مضافًا إلى أربعة أضعاف مسافة سير ب نحو الشرق (210 خطوات) فتصبح 840؛ والمجموع 1380 خطوة وهي الكشك الموجب. نستخرج الجذر التكعبي فنحصل على 120 خطوة. نضعفها فنحصل على قطر المدينة. يتطابق مع المطلوب.}



\end{paracol}



\vspace{0.5cm}



% ------------------------------------------------------------------------------

\subsubsection*{Problem 13 / \ar{المسألة الثالثة عشرة} \quad 第十三問}



\begin{paracol}{2}



\noindent 或問：甲出西門南行四百八十步而止，乙出東門南行二百四十步而止，望見甲，又斜行六百三十步與甲相會，問答同前。



\vspace{0.2cm}



\Answer{城徑二百四十步。}



\vspace{0.2cm}



\Method{兩行相減，以乘兩行相併為實，兩行相減自之為從方，兩倍兩行相減又益四之東行為益廉，一步常法，得半徑。}



\vspace{0.2cm}



\Workings{立天元一為半徑。以減甲南行步，得為梯底差；以乙南行步內減天元，得為梯頭差。以梯頭差乘梯底差，得為半徑冪，寄左。乃以兩行相減，餘二百四十步；又以兩行相併，得七百二十步；以減餘乘之，得十七萬二千八百步為實。又以減餘自之，得五萬七千六百步為從方。又以減餘倍之，得四百八十步；又益四之乙東行二百四十步，得九百六十步；共得一千四百四十步為益廉。以立方開之，得一百二十步。倍之，即城徑也。合問。}



\switchcolumn



\ar{سُئل: يخرج أ من الباب الغربي ويسير نحو الجنوب 480 خطوة ويقف. يخرج ب من الباب الشرقي ويسير نحو الجنوب 240 خطوة ويقف، ويرى أ، ثم يسير مائلًا 630 خطوة حتى يلتقي بأ. الجواب كالسابق.}



\vspace{0.2cm}



\AnswerAr{قطر المدينة 240 خطوة.}



\vspace{0.2cm}



\MethodAr{اطرح المسافتين، واضرب الناتج بمجموع المسافتين فتكون الذات. اطرح المسافتين وربّع الناتج فتكون الطرف المرافق. ضعف فرق المسافتين مضافًا إلى أربعة أضعاف مسافة السير نحو الشرق فتكون الكشك الموجب. وواحدًا القاسية الثابتة، فتحصل على نصف القطر.}



\vspace{0.2cm}



\WorkingsAr{نضع عنصر السماء يساوي نصف القطر. نطرحه من مسافة سير أ نحو الجنوب فنحصل على فرق قاع السلّم؛ ونطرح عنصر السماء من مسافة سير ب نحو الجنوب فنحصل على فرق رأس السلّم. نضرب فرق رأس السلّم بفرق قاعه فنحصل على مربّع نصف القطر، ونضعها على اليسار. نطرح المسافتين فيتبقّى 240 خطوة؛ ونجمع المسافتين فنحصل على 720 خطوة؛ نضرب الباقي بالمجموع فنحصل على 172800 خطوة وهي الذات. نربّع الباقي فنحصل على 57600 وهي الطرف المرافق. وضعف الباقي هو 480 خطوة؛ مضافًا إلى أربعة أضعاف مسافة سير ب نحو الشرق (240 خطوة) فتصبح 960؛ والمجموع 1440 خطوة وهي الكشك الموجب. نستخرج الجذر التكعبي فنحصل على 120 خطوة. نضعفها فنحصل على قطر المدينة. يتطابق مع المطلوب.}



\end{paracol}



\vspace{0.5cm}



% ------------------------------------------------------------------------------

\subsubsection*{Problem 14 / \ar{المسألة الرابعة عشرة} \quad 第十四問}



\begin{paracol}{2}



\noindent 或問：甲出西門南行四百八十步而止，乙出東門南行二百七十步而止，望見甲，又斜行六百九十步與甲相會，問答同前。



\vspace{0.2cm}



\Answer{城徑二百四十步。}



\vspace{0.2cm}



\Method{兩行相減，以乘兩行相併為實，兩行相減自之為從方，兩倍兩行相減又益四之東行為益廉，一步常法，得半徑。}



\vspace{0.2cm}



\Workings{立天元一為半徑。以減甲南行步，得為梯底差；以乙南行步內減天元，得為梯頭差。以梯頭差乘梯底差，得為半徑冪，寄左。乃以兩行相減，餘二百一十步；又以兩行相併，得七百五十步；以減餘乘之，得十五萬七千五百步為實。又以減餘自之，得四萬四千一百步為從方。又以減餘倍之，得四百二十步；又益四之乙東行二百七十步，得一千零八十步；共得一千五百步為益廉。以立方開之，得一百二十步。倍之，即城徑也。合問。}



\switchcolumn



\ar{سُئل: يخرج أ من الباب الغربي ويسير نحو الجنوب 480 خطوة ويقف. يخرج ب من الباب الشرقي ويسير نحو الجنوب 270 خطوة ويقف، ويرى أ، ثم يسير مائلًا 690 خطوة حتى يلتقي بأ. الجواب كالسابق.}



\vspace{0.2cm}



\AnswerAr{قطر المدينة 240 خطوة.}



\vspace{0.2cm}



\MethodAr{اطرح المسافتين، واضرب الناتج بمجموع المسافتين فتكون الذات. اطرح المسافتين وربّع الناتج فتكون الطرف المرافق. ضعف فرق المسافتين مضافًا إلى أربعة أضعاف مسافة السير نحو الشرق فتكون الكشك الموجب. وواحدًا القاسية الثابتة، فتحصل على نصف القطر.}



\vspace{0.2cm}



\WorkingsAr{نضع عنصر السماء يساوي نصف القطر. نطرحه من مسافة سير أ نحو الجنوب فنحصل على فرق قاع السلّم؛ ونطرح عنصر السماء من مسافة سير ب نحو الجنوب فنحصل على فرق رأس السلّم. نضرب فرق رأس السلّم بفرق قاعه فنحصل على مربّع نصف القطر، ونضعها على اليسار. نطرح المسافتين فيتبقّى 210 خطوات؛ ونجمع المسافتين فنحصل على 750 خطوة؛ نضرب الباقي بالمجموع فنحصل على 157500 خطوة وهي الذات. نربّع الباقي فنحصل على 44100 وهي الطرف المرافق. وضعف الباقي هو 420 خطوة؛ مضافًا إلى أربعة أضعاف مسافة سير ب نحو الشرق (270 خطوة) فتصبح 1080؛ والمجموع 1500 خطوة وهي الكشك الموجب. نستخرج الجذر التكعبي فنحصل على 120 خطوة. نضعفها فنحصل على قطر المدينة. يتطابق مع المطلوب.}



\end{paracol}



\vspace{0.5cm}



% ------------------------------------------------------------------------------

\subsubsection*{Problem 15 / \ar{المسألة الخامسة عشرة} \quad 第十五問}



\begin{paracol}{2}



\noindent 或問：甲出西門南行四百八十步而止，乙出東門南行三百步而止，望見甲，又斜行七百五十步與甲相會，問答同前。



\vspace{0.2cm}



\Answer{城徑二百四十步。}



\vspace{0.2cm}



\Method{兩行相減，以乘兩行相併為實，兩行相減自之為從方，兩倍兩行相減又益四之東行為益廉，一步常法，得半徑。}



\vspace{0.2cm}



\Workings{立天元一為半徑。以減甲南行步，得為梯底差；以乙南行步內減天元，得為梯頭差。以梯頭差乘梯底差，得為半徑冪，寄左。乃以兩行相減，餘一百八十步；又以兩行相併，得七百八十步；以減餘乘之，得十四萬零四百步為實。又以減餘自之，得三萬二千四百步為從方。又以減餘倍之，得三百六十步；又益四之乙東行三百步，得一千二百步；共得一千五百六十步為益廉。以立方開之，得一百二十步。倍之，即城徑也。合問。}



\switchcolumn



\ar{سُئل: يخرج أ من الباب الغربي ويسير نحو الجنوب 480 خطوة ويقف. يخرج ب من الباب الشرقي ويسير نحو الجنوب 300 خطوة ويقف، ويرى أ، ثم يسير مائلًا 750 خطوة حتى يلتقي بأ. الجواب كالسابق.}



\vspace{0.2cm}



\AnswerAr{قطر المدينة 240 خطوة.}



\vspace{0.2cm}



\MethodAr{اطرح المسافتين، واضرب الناتج بمجموع المسافتين فتكون الذات. اطرح المسافتين وربّع الناتج فتكون الطرف المرافق. ضعف فرق المسافتين مضافًا إلى أربعة أضعاف مسافة السير نحو الشرق فتكون الكشك الموجب. وواحدًا القاسية الثابتة، فتحصل على نصف القطر.}



\vspace{0.2cm}



\WorkingsAr{نضع عنصر السماء يساوي نصف القطر. نطرحه من مسافة سير أ نحو الجنوب فنحصل على فرق قاع السلّم؛ ونطرح عنصر السماء من مسافة سير ب نحو الجنوب فنحصل على فرق رأس السلّم. نضرب فرق رأس السلّم بفرق قاعه فنحصل على مربّع نصف القطر، ونضعها على اليسار. نطرح المسافتين فيتبقّى 180 خطوة؛ ونجمع المسافتين فنحصل على 780 خطوة؛ نضرب الباقي بالمجموع فنحصل على 140400 خطوة وهي الذات. نربّع الباقي فنحصل على 32400 وهي الطرف المرافق. وضعف الباقي هو 360 خطوة؛ مضافًا إلى أربعة أضعاف مسافة سير ب نحو الشرق (300 خطوة) فتصبح 1200؛ والمجموع 1560 خطوة وهي الكشك الموجب. نستخرج الجذر التكعبي فنحصل على 120 خطوة. نضعفها فنحصل على قطر المدينة. يتطابق مع المطلوب.}



\end{paracol}



\vspace{0.5cm}



% ------------------------------------------------------------------------------

\subsubsection*{Problem 16 / \ar{المسألة السادسة عشرة} \quad 第十六問}



\begin{paracol}{2}



\noindent 或問：甲出西門南行四百八十步而止，乙出東門南行三百三十步而止，望見甲，又斜行八百一十步與甲相會，問答同前。



\vspace{0.2cm}



\Answer{城徑二百四十步。}



\vspace{0.2cm}



\Method{兩行相減，以乘兩行相併為實，兩行相減自之為從方，兩倍兩行相減又益四之東行為益廉，一步常法，得半徑。}



\vspace{0.2cm}



\Workings{立天元一為半徑。以減甲南行步，得為梯底差；以乙南行步內減天元，得為梯頭差。以梯頭差乘梯底差，得為半徑冪，寄左。乃以兩行相減，餘一百五十步；又以兩行相併，得八百一十步；以減餘乘之，得十二万一千五百步為實。又以減餘自之，得二萬二千五百步為從方。又以減餘倍之，得三百步；又益四之乙東行三百三十步，得一千三百二十步；共得一千六百二十步為益廉。以立方開之，得一百二十步。倍之，即城徑也。合問。}



\switchcolumn



\ar{سُئل: يخرج أ من الباب الغربي ويسير نحو الجنوب 480 خطوة ويقف. يخرج ب من الباب الشرقي ويسير نحو الجنوب 330 خطوة ويقف، ويرى أ، ثم يسير مائلًا 810 خطوات حتى يلتقي بأ. الجواب كالسابق.}



\vspace{0.2cm}



\AnswerAr{قطر المدينة 240 خطوة.}



\vspace{0.2cm}



\MethodAr{اطرح المسافتين، واضرب الناتج بمجموع المسافتين فتكون الذات. اطرح المسافتين وربّع الناتج فتكون الطرف المرافق. ضعف فرق المسافتين مضافًا إلى أربعة أضعاف مسافة السير نحو الشرق فتكون الكشك الموجب. وواحدًا القاسية الثابتة، فتحصل على نصف القطر.}



\vspace{0.2cm}



\WorkingsAr{نضع عنصر السماء يساوي نصف القطر. نطرحه من مسافة سير أ نحو الجنوب فنحصل على فرق قاع السلّم؛ ونطرح عنصر السماء من مسافة سير ب نحو الجنوب فنحصل على فرق رأس السلّم. نضرب فرق رأس السلّم بفرق قاعه فنحصل على مربّع نصف القطر، ونضعها على اليسار. نطرح المسافتين فيتبقّى 150 خطوة؛ ونجمع المسافتين فنحصل على 810 خطوات؛ نضرب الباقي بالمجموع فنحصل على 121500 خطوة وهي الذات. نربّع الباقي فنحصل على 22500 وهي الطرف المرافق. وضعف الباقي هو 300 خطوة؛ مضافًا إلى أربعة أضعاف مسافة سير ب نحو الشرق (330 خطوة) فتصبح 1320؛ والمجموع 1620 خطوة وهي الكشك الموجب. نستخرج الجذر التكعبي فنحصل على 120 خطوة. نضعفها فنحصل على قطر المدينة. يتطابق مع المطلوب.}



\end{paracol}



\vspace{0.5cm}



% ------------------------------------------------------------------------------

\subsubsection*{Problem 17 / \ar{المسألة السابعة عشرة} \quad 第十七問}



\begin{paracol}{2}



\noindent 或問：甲出西門南行四百八十步而止，乙出東門南行三百六十步而止，望見甲，又斜行八百七十步與甲相會，問答同前。



\vspace{0.2cm}



\Answer{城徑二百四十步。}



\vspace{0.2cm}



\Method{兩行相減，以乘兩行相併為實，兩行相減自之為從方，兩倍兩行相減又益四之東行為益廉，一步常法，得半徑。}



\vspace{0.2cm}



\Workings{立天元一為半徑。以減甲南行步，得為梯底差；以乙南行步內減天元，得為梯頭差。以梯頭差乘梯底差，得為半徑冪，寄左。乃以兩行相減，餘一百二十步；又以兩行相併，得八百四十步；以減餘乘之，得十萬零八百步為實。又以減餘自之，得一万四千四百步為從方。又以減餘倍之，得二百四十步；又益四之乙東行三百六十步，得一千四百四十步；共得一千六百八十步為益廉。以立方開之，得一百二十步。倍之，即城徑也。合問。}



\switchcolumn



\ar{سُئل: يخرج أ من الباب الغربي ويسير نحو الجنوب 480 خطوة ويقف. يخرج ب من الباب الشرقي ويسير نحو الجنوب 360 خطوة ويقف، ويرى أ، ثم يسير مائلًا 870 خطوة حتى يلتقي بأ. الجواب كالسابق.}



\vspace{0.2cm}



\AnswerAr{قطر المدينة 240 خطوة.}



\vspace{0.2cm}



\MethodAr{اطرح المسافتين، واضرب الناتج بمجموع المسافتين فتكون الذات. اطرح المسافتين وربّع الناتج فتكون الطرف المرافق. ضعف فرق المسافتين مضافًا إلى أربعة أضعاف مسافة السير نحو الشرق فتكون الكشك الموجب. وواحدًا القاسية الثابتة، فتحصل على نصف القطر.}



\vspace{0.2cm}



\WorkingsAr{نضع عنصر السماء يساوي نصف القطر. نطرحه من مسافة سير أ نحو الجنوب فنحصل على فرق قاع السلّم؛ ونطرح عنصر السماء من مسافة سير ب نحو الجنوب فنحصل على فرق رأس السلّم. نضرب فرق رأس السلّم بفرق قاعه فنحصل على مربّع نصف القطر، ونضعها على اليسار. نطرح المسافتين فيتبقّى 120 خطوة؛ ونجمع المسافتين فنحصل على 840 خطوة؛ نضرب الباقي بالمجموع فنحصل على 100800 خطوة وهي الذات. نربّع الباقي فنحصل على 14400 وهي الطرف المرافق. وضعف الباقي هو 240 خطوة؛ مضافًا إلى أربعة أضعاف مسافة سير ب نحو الشرق (360 خطوة) فتصبح 1440؛ والمجموع 1680 خطوة وهي الكشك الموجب. نستخرج الجذر التكعبي فنحصل على 120 خطوة. نضعفها فنحصل على قطر المدينة. يتطابق مع المطلوب.}



\end{paracol}



\newpage





% =============================================================================

\section{The Tian Yuan Shu Method / \ar{طريقة عنصر السماء 天元術}}

\label{sec:tianyuan}



\begin{paracol}{2}



The \textit{tian yuan shu} (天元術, ``method of the celestial element'') is the algebraic technique employed throughout the \textit{Ceyuan Haijing}. It is the earliest systematic method in Chinese mathematics for setting up and solving polynomial equations with one unknown.



\vspace{0.3cm}



The name 天元 (tian yuan, ``celestial element'') refers to the unknown quantity in an equation. The method is equivalent to what we now call polynomial algebra in one variable.



\switchcolumn



\ar{إنّ \textit{طريقة عنصر السماء} (天元術، «طريقة عنصر السماء») هي التقنية الجبرية المستخدمة في جميع أنحاء كتاب «مرآة البحر لقياس الدائرة». وهي أول طريقة منهجية في الرياضيات الصينية لوضع وحلّ المعادلات متعدّدة الحدود بمجهول واحد.}



\vspace{0.3cm}



\ar{يشير اسم 天元 (تيان يوان، «عنصر السماء») إلى الكمّية المجهولة في المعادلة. وتُعادل هذه الطريقة ما نسمّيه الآن بالجبر متعدّد الحدود بمتغيّر واحد.}



\end{paracol}



% ------------------------------------------------------------------------------

\subsection{Methodology / \ar{المنهجية}}



\begin{paracol}{2}



The procedure follows these steps:



\begin{enumerate}[leftmargin=2em]

  \item \textbf{立天元一為\textit{X}} (``Set the celestial element unity to be $X$'') --- This declares the unknown quantity. In modern notation, it is equivalent to ``Let $x$ = \textit{X}.''

  

  \item Express all given quantities and relationships in terms of the celestial element $x$.

  

  \item Construct two equal expressions (often representing the same geometric quantity computed by different methods), one designated as the ``left'' (左) expression and stored temporarily, the other as the ``equal number'' (同數).

  

  \item \textbf{相消} (``mutual cancellation'') --- Set the two expressions equal and simplify to obtain the polynomial equation in standard form.

  

  \item Solve the equation by root-extraction methods (開方): square-root extraction (開平方), cube-root extraction (開立方), or higher-degree root extraction.

\end{enumerate}



\switchcolumn



\ar{تتّبع العملية الخطوات التالية:}



\begin{enumerate}[leftmargin=2em,rightmargin=0em]

  \item \ar{\textbf{نضع عنصر السماء يساوي \textit{س}} (立天元一為\textit{X}) — هذا يعلن عن الكمّية المجهولة. في الرمز الحديث، يُعادل «لتكن $x$ = \textit{X}».}



  \item \ar{عبّر عن جميع الكمّيات المعطاة والعلاقات بدلالة عنصر السماء $x$.}



  \item \ar{ابنِ تعبيرين متساويين (غالبًا ما يمثّلان الكمّية الهندسية نفسها المحسوبة بطريقتين مختلفتين)، أحدهما يُسمّى التعبير «الأيسر» (左) ويُحفظ مؤقتًا، والآخر يُسمّى «العدد المساوي» (同數).}



  \item \ar{\textbf{الإلغاء المتبادل} (相消) — اجعل التعبيرين متساويين وبسّط للحصول على المعادلة متعدّدة الحدود بالصيغة القياسية.}



  \item \ar{حلّ المعادلة بطرق استخراج الجذر (開方): استخراج الجذر التربيعي (開平方)، أو استخراج الجذر التكعبي (開立方)، أو استخراج جذر الدرجة العليا.}

\end{enumerate}



\end{paracol}



% ------------------------------------------------------------------------------

\subsection{Polynomial Notation / \ar{الرمز متعدّد الحدود}}



\begin{paracol}{2}



Li Ye uses a positional notation on a counting-rod board:

\begin{itemize}[leftmargin=2em]

  \item Coefficients are arranged vertically, with the constant term (實) at the top.

  \item Below it: the coefficient of $x$ (方, ``square''), then $x^2$ (廉, ``edge''), then $x^3$ (隅, ``corner'').

  \item For higher degrees: first 廉, second 廉, etc.

  \item Negative coefficients are indicated by a diagonal stroke through the last digit.

\end{itemize}



\switchcolumn



\ar{يستخدم لي يه رمزًا موضعيًا على لوحة أعواد الحساب:}

\begin{itemize}[leftmargin=2em,rightmargin=0em]

  \item \ar{تُرتّب المعاملات عموديًا، مع وضع الحدّ الثابت (實) في الأعلى.}

  \item \ar{تحته: معامل $x$ (方، «المربّع»)، ثم $x^2$ (廉، «الكشك»)، ثم $x^3$ (隅، «الركن»).}

  \item \ar{للدرجات العليا: الكشك الأول، الكشك الثاني، إلخ.}

  \item \ar{تُشير المعاملات السالبة بخط مائل عبر الرقم الأخير.}

\end{itemize}



\end{paracol}



\vspace{0.5cm}



% ------------------------------------------------------------------------------

\subsection{Example: Quadratic Equation / \ar{مثال: معادلة تربيعية}}



\begin{paracol}{2}



The first problem leads to the equation equivalent to:

\[

384000 = 1600 \cdot d

\]

where $d$ is the diameter. But more complex problems yield true polynomial equations. For instance, Problem 4 of Volume 3 produces a \textbf{cubic equation}:

\[

(4 \cdot 16)(480)^2 = (480 + 16) \cdot d^2 - d^3

\]

in modern form, which Li Ye solves by ``opening the cube with an edge'' (帶縱立方).



\switchcolumn



\ar{تؤدّي المسألة الأولى إلى المعادلة المكافئة:}

\[

384000 = 1600 \cdot d

\]

\ar{حيث $d$ هو القطر. لكنّ المسائل الأكثر تعقيدًا تُنتج معادلات متعدّدة الحدود حقيقية. فعلى سبيل المثال، تُنتج المسألة الرابعة من المجلد الثالث \textbf{معادلة تكعيبية}:}

\[

(4 \cdot 16)(480)^2 = (480 + 16) \cdot d^2 - d^3

\]

\ar{بالصيغة الحديثة، والتي يحلّها لي يه بـ«استخراج الجذر التكعبي ذي الطرف» (帶縱立方).}



\end{paracol}



\vspace{0.5cm}



% ------------------------------------------------------------------------------

\subsection{Historical Significance / \ar{الأهمية التاريخية}}



\begin{paracol}{2}



The \textit{Ceyuan Haijing} is the earliest surviving work that systematically employs the tian yuan shu. The equations solved include:



\begin{center}

\begin{tabular}{lr}

\toprule

\textbf{Degree} & \textbf{Number of Equations} \\

\midrule

Linear (degree 1) & 31 \\

Quadratic (degree 2) & 106 \\

Cubic (degree 3) & 24 \\

Quartic (degree 4) & 20 \\

Sextic (degree 6) & 1 \\

\midrule

\textbf{Total} & \textbf{182} \\

\bottomrule

\end{tabular}

\end{center}



\switchcolumn



\ar{إنّ كتاب «مرآة البحر لقياس الدائرة» هو أقدم عمل باقٍ يستخدم طريقة عنصر السماء بشكل منهجي. وتشمل المعادلات المحلولة:}



\begin{center}

\begin{tabular}{lr}

\toprule

\textbf{\ar{الدرجة}} & \textbf{\ar{عدد المعادلات}} \\

\midrule

\ar{خطية (درجة 1)} & 31 \\

\ar{تربيعية (درجة 2)} & 106 \\

\ar{تكعيبية (درجة 3)} & 24 \\

\ar{رباعية (درجة 4)} & 20 \\

\ar{سداسية (درجة 6)} & 1 \\

\midrule

\textbf{\ar{المجموع}} & \textbf{182} \\

\bottomrule

\end{tabular}

\end{center}



\end{paracol}



\vspace{0.5cm}



% ------------------------------------------------------------------------------

\subsection{Key Terms Glossary / \ar{قائمة المصطلحات الرئيسية}}



\vspace{0.3cm}



\begin{center}

\begin{tabular}{lll}

\toprule

\textbf{Chinese} & \textbf{Arabic} & \textbf{English} \\

\midrule

天元 & \ar{عنصر السماء} (unsur al-sama') & Celestial Element / Unknown \\

天元術 & \ar{طريقة عنصر السماء} (tariqat unsur al-sama') & Method of the Celestial Element \\

立天元一為 & \ar{نضع عنصر السماء يساوي} & Let the unknown equal \\

勾股 & \ar{الضلعان القائمان} (al-dil'an al-qaiman) & Two legs of right triangle \\

勾股容圓 & \ar{الدائرة المحاطة بالمثلث القائم} & Circle inscribed in right triangle \\

弦 & \ar{الوتر} (al-watar) & Hypotenuse \\

股 & \ar{الضلع الطويل} (al-dhil' al-tawil) & Longer vertical leg \\

勾 & \ar{القاعدة} (al-qaida) & Shorter horizontal leg \\

相消 & \ar{الإلغاء المتبادل} (al-ilgha' al-mutabadil) & Mutual cancellation \\

開方 & \ar{استخراج الجذر} (istikhraj al-jadhr) & Root extraction \\

開平方 & \ar{استخراج الجذر التربيعي} & Square root extraction \\

開立方 & \ar{استخراج الجذر التكعبي} & Cube root extraction \\

實 & \ar{الذات} (al-dhat) & Constant term / Dividend \\

方 & \ar{المربّع} (al-murabba') & Coefficient of $x$ \\

廉 & \ar{الكشك} (al-kashk) & Coefficient of $x^2$ \\

隅 & \ar{الركن} (al-rukn) & Coefficient of $x^3$ \\

從 & \ar{الطرف المرافق} & Linear coefficient (in certain forms) \\

帶縱 & \ar{ذو الطرف} & With linear term \\

城徑 & \ar{قطر المدينة} & City diameter \\

\bottomrule

\end{tabular}

\end{center}



\vspace{0.5cm}



\noindent\textcolor{rulecolor}{\rule{\textwidth}{0.4pt}}



\vspace{0.5cm}



\begin{paracol}{2}

\textit{Note:} This bilingual edition presents the complete text of Volumes 1--3 of the \textit{Ceyuan Haijing}. Volume 1 contains all definitions, numerical data, and 692 geometric formulas. Volumes 2 and 3 present the first 31 of 170 problems, solved by the tian yuan shu method. The remaining nine volumes (卷四--卷十二) continue with 139 further problems of increasing complexity, culminating in a sixth-degree equation in Volume 7.



\switchcolumn



\ar{ملاحظة: تعرض هذه الطبعة الثنائية النص الكامل للمجلدات 1–3 من «مرآة البحر لقياس الدائرة». يحتوي المجلد الأول على جميع التعريفات والبيانات الرقمية و692 صيغة هندسية. وتعرض المجلدات الثاني والثالث أول 31 مسألة من أصل 170، تُحلّ بطريقة عنصر السماء. ويستمر المجلدات التسعة المتبقية (卷四–卷十二) بـ139 مسألة إضافية ذات تعقيد متزايد، تبلغ ذروتها بمعادلة من الدرجة السادسة في المجلد السابع.}



\end{paracol}



\vspace{1cm}



\begin{center}

{\color{rulecolor}\rule{0.4\textwidth}{0.4pt}}



\vspace{0.5cm}



{\small\color{sectioncolor}

\textit{End of Volumes 1--3 / \ar{نهاية المجلدات 1–3}}

}



\vspace{0.3cm}



{\footnotesize\color{notecolor}

\textit{測圓海鏡卷第一至卷第三終}

}

\end{center}



\end{document}







% ============================================================

% Source: translations/ar/chinese/qin-jiushao-shuxue-jiuzhang-fascicle1_arabic.tex

% ============================================================



%% ========================================================================

%% Qin Jiushao's Shuxue Jiuzhang (数学九章) — Fascicle 1: Dayan Category

%% CHINESE → ARABIC BILINGUAL TRANSLATION EDITION

%% ========================================================================

\documentclass[11pt,a4paper]{article}



%% -------- Core Packages --------

\usepackage{fontspec}

\usepackage{amsmath,amssymb,amsthm}

\usepackage{geometry}

\usepackage{graphicx}

\usepackage{xcolor}

\usepackage{fancyhdr}

\usepackage{array,booktabs,longtable,tabularx}

\usepackage{hyperref}

\usepackage{etoolbox}

\usepackage{xeCJK}



%% -------- Page Geometry --------

\geometry{

  paperwidth=210mm,paperheight=297mm,

  left=15mm,right=15mm,top=22mm,bottom=22mm,

  headheight=14pt,footskip=24pt

}



%% -------- Font Configuration --------

%% Arabic font (for use inside \beginR...\endR blocks)

\newfontfamily\arabicfont[Script=Arabic]{Arial}

\newfontfamily\arabicfontbf[Script=Arabic]{Arial}

\TeXXeTstate=1

\newcommand{\ar}[1]{{\arabicfont\beginR #1\endR}}

\newcommand{\arbf}[1]{{\arabicfontbf\beginR #1\endR}}

\newcommand{\artxt}[1]{\arabicfont\beginR #1\endR}



%% Bullet command that uses Latin font

\newcommand{\arbullet}{{\fontspec{Times New Roman}\textbullet}}



%% CJK font for Chinese

\setCJKmainfont{Microsoft YaHei}

\newCJKfontfamily\chinesebf{Microsoft YaHei}

\newCJKfontfamily\chineselg[Scale=1.15]{Microsoft YaHei}

\newfontfamily\chfont{Microsoft YaHei}

\newfontfamily\chfontbf{Microsoft YaHei}



%% Latin font

\setmainfont{Times New Roman}

\setsansfont{Arial}



%% Chinese text line breaking

\tolerance=3000

\emergencystretch=3em

\hbadness=10000



%% -------- Colors --------

\definecolor{titlecolor}{RGB}{45,55,72}

\definecolor{sectioncolor}{RGB}{55,65,81}

\definecolor{notecolor}{RGB}{120,53,15}

\definecolor{rulecolor}{RGB}{180,160,140}

\definecolor{prefacecolor}{RGB}{80,60,40}

\definecolor{colrule}{RGB}{160,140,120}



%% -------- Parallel Text Layout --------

\newlength{\cnwidth}

\newlength{\arwidth}

\setlength{\cnwidth}{0.44\textwidth}

\setlength{\arwidth}{0.44\textwidth}



%% Parallel text: Chinese left, Arabic right

\newcommand{\paralleltext}[2]{%

  \par\medskip\noindent

  \begin{minipage}[t]{\cnwidth}%

    \small #1%

  \end{minipage}%

  \hfill\textcolor{colrule}{\vrule width 0.6pt}\hfill

  \begin{minipage}[t]{\arwidth}%

    \small\arabicfont\beginR #2\endR%

  \end{minipage}%

  \par\medskip

}



%% Section header in parallel

\newcommand{\parallelsection}[2]{%

  \par\bigskip\noindent

  \begin{minipage}[t]{\cnwidth}%

    \Large\bfseries\color{sectioncolor} #1%

  \end{minipage}%

  \hfill\textcolor{colrule}{\vrule width 0.8pt}\hfill

  \begin{minipage}[t]{\arwidth}%

    \Large\arabicfontbf\color{sectioncolor}\beginR #2\endR%

  \end{minipage}%

  \par\medskip\noindent\textcolor{colrule}{\rule{\textwidth}{0.6pt}}\par\medskip

}



\newcommand{\parallelsubsection}[2]{%

  \par\medskip\noindent

  \begin{minipage}[t]{\cnwidth}%

    \large\bfseries\color{sectioncolor} #1%

  \end{minipage}%

  \hfill\textcolor{colrule}{\vrule width 0.6pt}\hfill

  \begin{minipage}[t]{\arwidth}%

    \large\arabicfontbf\color{sectioncolor}\beginR #2\endR%

  \end{minipage}%

  \par\smallskip

}



%% Labels

\newcommand{\cnlabel}[1]{\textbf{\color{notecolor}#1}}

\newcommand{\arlabel}[1]{\arabicfontbf\color{notecolor}\beginR #1\endR}



%% Chinese environments

\newenvironment{cnquote}{\begin{quote}\color{prefacecolor}}{\end{quote}}



%% Header/Footer

\pagestyle{fancy}

\fancyhf{}

\fancyhead[L]{\small\chfont 数学九章 · 大衍类}

\fancyhead[R]{\small\arabicfont\beginR الطريقة الكبرى\endR}

\fancyhead[C]{\small\thepage}

\renewcommand{\headrulewidth}{0.4pt}

\renewcommand{\footrulewidth}{0pt}



\hypersetup{

  colorlinks=true,linkcolor=sectioncolor,citecolor=sectioncolor,urlcolor=sectioncolor,

  pdfencoding=auto,

  pdftitle={Qin Jiushao - Shuxue Jiuzhang Fascicle 1 - Chinese-Arabic Edition},

  pdfauthor={Bilingual Translation Edition}

}



%% ========================================================================

\begin{document}



%% ========================================================================

%% TITLE PAGE

%% ========================================================================

\begin{titlepage}

\thispagestyle{empty}

\centering

\vspace*{1.5cm}



{\fontsize{24}{30}\selectfont\chineselg

\textcolor{titlecolor}{数书九章}}



\vspace{0.4cm}



{\fontsize{18}{24}\selectfont\arabicfont

\textcolor{titlecolor}{\beginR كتاب المسائل التسع في الحساب\endR}}



\vspace{0.8cm}



{\Large\textbf{Fascicle I / الفصل الأول}}



\vspace{0.3cm}



{\large

\textcolor{sectioncolor}{The Dayan Category (\chfont 大衍类)}\\[0.3em]

\textcolor{sectioncolor}{\arabicfont\beginR باب الطريقة الكبرى\endR}

}



\vspace{1.2cm}



{\large

\textbf{Qin Jiushao} / \chfont 秦九韶 / \arabicfont\beginR تشين جيوشاو\endR\\[0.3em]

c.\ 1202--1261 CE / \arabicfont\beginR حوالي ٦٠٢--٦٦١ هـ\endR\\[0.5em]

\textit{Presented in the 7th year of Chunyou ({\chfont 淳祐七年}), 1247 CE}\\[0.2em]

{\arabicfont\beginR وُلِّف في السنة السابعة من حقبة تشونيو، ٦٤٥ هـ / ١٢٤٧ م\endR}

}



\vspace{1cm}



{\normalsize\textit{

One of the greatest masterpieces of Chinese mathematics,\\

containing the general Chinese Remainder Theorem (Dayan method).}}



\vspace{0.3cm}



{\arabicfont\beginR 

من أعظم مصنفات الرياضيات الصينية، يضم نظرية الباقي الصينية العامة (الطريقة الكبرى).

\endR}



\vspace{1cm}



\textcolor{colrule}{\rule{0.6\textwidth}{0.8pt}}



\vspace{0.5cm}



{\normalsize\textbf{\textsc{Chinese--Arabic Bilingual Edition}}\\[0.2em]

\arabicfont\beginR الطبعة الثنائية اللغة: الصينية--العربية\endR}



\vspace{0.5cm}



{\small

\textit{Left column}: Original Classical Chinese\\

\textit{Right column}: Arabic Translation\\[0.2em]

\arabicfont\beginR العمود الأيسر: النص الصيني الكلاسيكي الأصلي --- العمود الأيمن: الترجمة العربية\endR

}



\vfill



{\small Modern LaTeX edition\ with XeTeX\\

Microsoft YaHei / Arial}



\end{titlepage}



%% ========================================================================

%% TABLE OF CONTENTS

%% ========================================================================

\newpage

\tableofcontents

\newpage



%% ========================================================================

%% INTRODUCTION

%% ========================================================================

\parallelsection{Introduction / 引言}{المقدمة}



\paralleltext{%

The \textit{Shuxue Jiuzhang} ({\chfont 数书九章}), Mathematical Treatise in Nine Sections) is the magnum opus of the Song dynasty mathematician Qin Jiushao ({\chfont 秦九韶}), c.~1202--1261). Completed in 1247 CE during the Chunyou era, it stands as one of the greatest achievements in the history of Chinese mathematics and indeed of world mathematics. The work is organized into nine categories ({\chfont 类})), each containing nine problems ({\chfont 问})), for a total of eighty-one problems.%

}{%

يُعد \textit{كتاب المسائل التسع في الحساب} (数学九章) تحفةَ عصرِ سلالة سونغ الرياضية التي ألَّفها عالم الرياضيات تشين جيوشاو (秦九韶، حوالي ١٢٠٢--١٢٦١ م). أُتمّ في عام ١٢٤٧ م في حقبة تشونيو، وهو من أعظم الإنجازات في تاريخ الرياضيات الصينية وفي تاريخ الرياضيات العالمية بشكل عام. يُنظَّم العمل في تسعة أبواب (类)، كل باب يتضمن تسع مسائل (问)، بمجموع واحد وثمانين مسألة.

}



\paralleltext{%

The nine categories are: 1.~\textbf{Dayan} ({\chfont 大衍})) --- indeterminate analysis; 2.~\textbf{Tianshi} ({\chfont 天时})) --- astronomical phenomena; 3.~\textbf{Tianyu} ({\chfont 田域})) --- land measurement; 4.~\textbf{Cewang} ({\chfont 测望})) --- surveying; 5.~\textbf{Fuyi} ({\chfont 赋役})) --- taxation; 6.~\textbf{Qian'gu} ({\chfont 钱谷})) --- commerce; 7.~\textbf{Yingjian} ({\chfont 营建})) --- construction; 8.~\textbf{Junlv} ({\chfont 军旅})) --- military affairs; 9.~\textbf{Shiyi} ({\chfont 市易})) --- market transactions.%

}{%

الأبواب التسعة هي: ١.~\textbf{الطريقة الكبرى} (大衍) --- التحليل اللامحدود. ٢.~\textbf{الظواهر الفلكية} (天时) --- الحسابات الفلكية. ٣.~\textbf{قياس الأراضي} (田域) --- مساحة الأراضي. ٤.~\textbf{القياس والمراقبة} (测望) --- المساحة والرصد. ٥.~\textbf{الضرائب والأشغال} (赋役) --- الضرائب والخدمات. ٦.~\textbf{النقود والحبوب} (钱谷) --- التجارة. ٧.~\textbf{البناء والهندسة} (营建) --- البناء والهندسة. ٨.~\textbf{الشؤون العسكرية} (军旅) --- الأمور العسكرية. ٩.~\textbf{تجارة السوق} (市易) --- معاملات السوق.

}



\paralleltext{%

Each problem follows a rigorous four-part structure:

\textbullet\ \textbf{Wen} ({\chfont 问})) --- the problem statement;

\textbullet\ \textbf{Da} ({\chfont 答})) --- the numerical answer;

\textbullet\ \textbf{Shu} ({\chfont 术})) --- the general method and algorithm;

\textbullet\ \textbf{Cao} ({\chfont 草})) --- the detailed working.%

}{%

كل مسألة تتبع هيكلاً رباعياً صارماً:

--- \textbf{المسألة} (问) صياغة المسألة.

--- \textbf{الجواب} (答) الإجابة الرقمية.

--- \textbf{الطريقة} (术) الطريقة العامة والخوارزمية.

--- \textbf{الحل المفصل} (草) العمليات الحسابية المفصلة.

}



\parallelsubsection{The Dayan Category (大衍类)}{باب الطريقة الكبرى (大衍类)}



\paralleltext{%

Fascicle I covers the first category, \textit{Dayan} ({\chfont 大衍}), ``Great Derivation''), devoted to the general solution of systems of linear congruences---what is now known worldwide as the \textbf{Chinese Remainder Theorem}. Qin Jiushao's \textit{Dayan zongshu shu} ({\chfont 大衍总数术}), ``Dayan General Method'') provides a systematic algorithm, and his \textit{Dayan qiu yi shu} ({\chfont 大衍求一术}), ``Dayan Method for Finding Unity'') gives the procedure for finding modular multiplicative inverses. The nine problems apply this method to divination, calendar calculation, engineering, finance, and commerce.%

}{%

يغطي الفصل الأول البابَ الأول، \textit{الطريقة الكبرى} (大衍، ``الاشتقاق العظيم'')، المكرس للحل العام لأنظمة التطابقات الخطية --- المعروفة عالمياً باسم \textbf{نظرية الباقي الصينية}. تُوفِّر طريقة تشين جيوشاو \textit{الطريقة العامة الكبرى} (大衍总数术) خوارزمية منهجية، وتُعطي \textit{طريقة استخراج الوحدة} (大衍求一术) الإجراء لإيجاد المعكوسات الضربية المعيارية. تطبق المسائل التسع هذه الطريقة على الكهانة والتقويم والهندسة والمالية والتجارة.

}



\newpage



%% ========================================================================

%% PREFACE

%% ========================================================================

\parallelsection{Preface (序)}{التمهيد (序)}



\paralleltext{%

\begin{cnquote}

周教六艺，数实成之。学士大夫所从事尚矣。其用本太虚生一，而周流无穷。大则可以通神明、顺性命，小则可以经世务、类万物，讵容以浅近窥哉？

\end{cnquote}%

}{%

لقد أكملت علومُ الرياضياتِ المبادئَ الستَّ التي علَّمها زوا [ملك سلالة جو]. إنَّها مَوضِعُ اهتمام العُلَماء والكُبَّراء منذ زمنٍ بعيدٍ جداً. أصلُها من الفراغ الأعظم الذي ولَّد الوحدة، فتدور دورةً لا نهاية لها. في كُبراها، يُمكنها أن تتصل بالروحانيَّات وتُطيع الطبيعة؛ وفي صُغراها، يُمكنها أن تدير شؤون العالم وتصنَّف الأشياء. كيف يُمكن اختزالها بنظرةٍ سطحية؟

}



\paralleltext{%

\begin{cnquote}

若昔推策以迎日，定律而和气。髀距濬川，土圭度晷。天地之大，囿焉而不能外，况其间总总者乎？爰自河图洛书，闿发幽秘；八卦九畴，错综精微。极而至于大衍、皇极之用，而人事之变无不该，鬼神之情莫能隐矣。

\end{cnquote}%

}{%

فمنذ القدم، حُسِبت الدرجات الفلكية لاستقبال الشمس، ورُتِّبت القوانين لتهيئة الانسجام. مِسطرة الكَعْبُوس لحفر الأنهار، ومِيزان التراب لقياس الظلال. عظمةُ السماء والأرض محصورةٌ فيها ولا يمكن الخروجُ عنها، فكيف بما هو أصغر؟ منذ خرائط النهر والكتاب المقدَّس، استُكشفت الأسرار الخفية؛ والرموز الثمانية والتسع فئات، تعقيدٌ دقيق. تصل إلى استخدامات الطريقة الكبرى والقمة العظمى، بحيث لا تخلو منها أيُّ تغيُّرات في شؤون البشر، ولا يمكن إخفاء طبيعة الأرواح.

}



\paralleltext{%

\begin{cnquote}

呜呼！乐有制氏，仅记铿锵之声，而谓与天地同和者止于是，可乎？今数术之书尚三十余家。天象历度谓之缀术，太乙壬甲谓之三式，皆曰内算，言其秘也。九章所载及周官九数，系于方圆者为专术，皆曰外算，对内而言也。其用相通，不可歧二。独大衍法不载九章，未有能推之者。历家演法颇用之，以为方程者，误也。

\end{cnquote}%

}{%

يا للعجب! إنَّ لموسيقى آلتها التي تسجل الأصوات الرنَّانة فحسب، أيكون هذا هو كلُّ ما يتوافق مع السماء والأرض؟ كلا! اليوم، كتبُ الرياضيات لا تزال تزيد عن ثلاثين مؤلِّفاً. حسابات الفلك والتقويم تُسمَّى فُنون الجمع، والطريقة العظمى والرموز الثلاثة تُسمَّى الصيغ الثلاث، وجميعها تُسمَّى ``الحساب الداخلي'' لسرِّيتها. ما ورد في الجُزْء التاسع وتسع قواعد زوا المرتبطة بالدائرة والمربع هي فنون خاصة تُسمَّى ``الحساب الخارجي''، مقابل الداخلي. استخداماتها متصلة ولا يمكن فصلها. إلاَّ أنَّ طريقة الطريقة الكبرى وحدها لم تُدرَج في الجُزْء التاسع، ولم يكن أحدٌ قادراً على استنباطها. استخدمها حاسبو التقويم كثيراً، لكنهم ظنُّوها معادلات — وهذا خطأ.

}



\paralleltext{%

\begin{cnquote}

九韶愚陋，不闲于艺。然早岁侍亲中都，因得访习于太史，又尝从隐君子受数学。时际兵难，历岁遥塞，不自意全于矢石之间。更险离忧，荏苒十禩，心槁气落。信知夫物莫不有数也，乃肆意其间，旁诹方能，探索杳眇，粗若有得焉。

\end{cnquote}%

}{%

أما أنا، جيوشاو، فجاهلٌ ضئيلٌ غير ماهرٍ في الفنون. لكن في صباي خدمتُ أهلي في العاصمة، فتمكَّنتُ من التعلُّم لدى كاتب التاريخ، وقد تلقَّيتُ الرياضيات من رجلٍ حكيمٍ متصوِّف. جاء وقتُ الحروب، وامتُحنتُ السنون الطويلة، ولم أتوقَّع النجاة بين السهام والحجارة. بعد المخاطر والفراق والهموم، عشر سنواتٍ مرتْ، وذبلَ قلبي وضعُفَتْ نفسي. آمنتُ أنَّ كلَّ شيءٍ له عددٌ، فانغمستُ فيها بكلِّيَّة، واستقيتُ الفنون من كلِّ جانب، واستكشفتُ الغامض الدقيق، حتى نلتُ من الفهم شيئاً.

}



\paralleltext{%

\begin{cnquote}

所谓通神明、顺性命，固肤末于见。若其小者，窃尝设为问答，以拟于用。积多而惜其弃，因取八十一题，厘为九类，立术具草，间以图发之。

\end{cnquote}%

}{%

أمَّا ما يُقال عن الاتصال بالروحانيَّات وإطاعة الطبيعة، فلا أزال بعيداً عنه. أمَّا الصغير منها، فقد جرَّبتُ أن أصيِّغَ أسئلةً وأجوبةً لأُحاكيَ بها التطبيقات. كثُرتْ عندي، فكرهتُ إهمالها، فاختَرْتُ إحدى وثمانين مسألة، رتَّبتُها في تسعة أبواب، وضعتُ لها الطرائق وشرحتُ الحلول، وأحياناً أوضحتُ بالأشكال.

}



\paralleltext{%

\begin{cnquote}

时淳祐七年九月，鲁郡秦九韶叙。

\end{cnquote}%

}{%

كتبَ هذا في شهر أيلول/سبتمبر من السنة السابعة لحقبة تشونيو، ١٢٤٧ م. --- تشين جيوشاو من مقاطعة لو.

}



\newpage



%% ========================================================================

%% THE DAYAN GENERAL METHOD

%% ========================================================================

\parallelsection{The Dayan General Method (大衍总数术)}{الطريقة العامة الكبرى (大衍总数术)}



\parallelsubsection{Editor's Commentary on the Dayan Method}{تعليق المحرِّر على الطريقة الكبرى}



\paralleltext{%

\begin{cnquote}

按大衍术，以各分数之奇零求各分数之总数。大而天行，小而物数，皆可御之。其法有求元、求定、求衍、求奇、求乘、求用之目。大约以数之奇偶为根，而以诸数相度之尽不尽为用。



有求彼此不能度尽之诸数者，元数、定数是也。有求诸数皆能度尽之一数者，衍母数是也。有求诸数皆能度尽而一数不能度尽之数者，各衍数是也。其不尽之数，即奇数也。有求二数相度余一之数者，乘数是也。有求二数相度余一而诸数又能度尽之数者，用数是也。

\end{cnquote}%

}{%

\begin{quote}

أمَّا طريقة الطريقة الكبرى، فهي تسعى للوصول إلى المجاميع الكلية من الكسور الفردية، كُبرى كالحركات السماوية أو صُغرى كأعداد الأشياء، وكلُّها يُمكن التغلُّب عليها. طرائقُها تتضمَّن: طلب الأصليَّات، وطلب المثبتات، وطلب المشتقَّات، وطلب الفرديَّات، وطلب المعاملات، وطلب المستعملات. في الجملة، تعتمدُ على الفرديَّة والزوجيَّة للأعداد كأصل، وعلى إمكانية أو عدم إمكانية القسمة المتبادلة للأعداد كأداة.



طلبُ الأعداد التي لا يمكن قسمتها بشكلٍ متبادلٍ --- هي الأعداد الأصليَّة والمثبتة. طلبُ العدد الواحد الذي يمكن لجميع الأعداد قسمته --- هو عدد الجذر المشترك. طلبُ الأعداد التي يمكن لجميع الأعداد قسمتها بينما لا يمكن لعددٍ واحدٍ قسمتها --- هي الأعداد المشتقَّة. والباقي من القسمة هو العدد الفردي. طلبُ العدد الذي إذا قُسِمَ على عددَين يبقَ واحداً --- هو معامل الضرب. والعدد الذي إذا قُسِمَ على عددَين يبقَ واحداً ويمكن لجميع الأعداد قسمته --- هو العدد المستعمل.

\end{quote}

}



\parallelsubsection{The Four Types of Numbers}{الأنواع الأربعة من الأعداد}



\paralleltext{%

\textbf{大衍总数术曰：}置诸问数，类名有四。\\

\textbf{1. 元数} --- 谓尾位见单零者。本门揲蓍、酒息、斛粜、砌砖、失米之类是也。\\

\textbf{2. 收数} --- 谓尾位见分厘者。假令冬至三百六十五日二十五刻，欲与甲子六十日为一会而求积日之类。\\

\textbf{3. 通数} --- 谓诸数各有分子分母者。本门问一会积年是也。\\

\textbf{4. 复数} --- 谓尾位见十或百及千以上者。本门筑堤并急足之类是也。%

}{%

\textbf{تقول الطريقة العامة الكبرى:} ضع جميع أعداد المسألة، ولها أربعة أسماء.\\

\textbf{١. الأعداد الأصليَّة} (元数) --- هي التي تظهر في نهايتها أصفاراً مفردة. مثل: العدد في الكهانة بالسعف، وفوائد الخمر، وبيع الحبوب بالكيل، وبناء الجدران بالطوب، وفقدان الأرز في هذا الباب.\\

\textbf{٢. الأعداد المجمَّعة} (收数) --- هي التي تظهر في نهايتها كسوراً. كأن تكون الشتاء ثلاثمئة وخمسة وستين يوماً وخمساً وعشرين دقيقة، تريد أن تجعلها مع دورة الستين يوماً اجتماعاً واحداً فتطلب الأيام المتراكمة.\\

\textbf{٣. الأعداد المشتركة} (通数) --- هي التي لها بسط ومقام. مثل مسألة اجتماع السنوات المتراكمة في هذا الباب.\\

\textbf{٤. الأعداد المركَّبة} (复数) --- هي التي تظهر في نهايتها عشرات أو مئات أو آلاف. مثل مسألة بناء السدود والرسل السريعين في هذا الباب.

}



\parallelsubsection{Finding the Definitive Numbers (求定数)}{طلب الأعداد المثبتة (求定数)}



\paralleltext{%

\textbf{元数者：}先以两两连环求等，约奇弗约偶。或约得五而彼有十，乃约偶弗约奇。或元数俱偶，约毕可存一位见偶。或皆约而犹有类数存，姑置之，俟与其他约遍而后乃与姑置者求等约之。或诸数皆不可尽类，则以诸元数命曰定数，以复数格入之。



\textbf{收数者：}乃命尾位分厘作单零，以进所问之数，定位讫，用元数格入之。或如意立数为母，收进分厘以从所问，用通数格入之。



\textbf{通数者：}置问数通分纳子，互乘之，皆曰通数。求总等，不约一位，约众位，得各原法数，用元数格入之。



\textbf{复数者：}问数尾位见十以上者，以诸数求总等，存一位约众位，始得元数。两两连环求等，约奇弗约偶，复乘偶；或约偶弗约奇，弗乘奇。%

}{%

\textbf{أمَّا الأعداد الأصليَّة:} ابتدئْ بِطلب المقاس المشترك بين كلِّ زوجَين من الأعداد، فاقسم الفرديَّ دون الزوجيَّ. أو إذا قُسِمَ على خمسة والآخر يقبل العشرة، فاقسم الزوجيَّ دون الفرديَّ. أو إذا كانت جميع الأعداد الأصليَّة زوجيَّة، فبعد القسمة يمكن الإبقاء على واحدةٍ زوجيَّة. أو إذا قُسِمَت جميعُها وبقيَ عددٌ من نوعٍ واحد، فاتركه مؤقَّتاً، ثم بعد القسمة الشاملة مع غيره اطلب المقاس المشترك واقسمْهُ. أو إذا لم يكن جميعُ الأعداد من نوعٍ واحد، فسمِّ جميع الأعداد الأصليَّة أعداداً مثبتة، وادخلها في طريقة الأعداد المركَّبة.



\textbf{أمَّا الأعداد المجمَّعة:} فاسمِّ الكسور في نهايتها أصفاراً مفردة، فتقدَّم بالأعداد المطلوبة، وبعد التحديد ادخلها في طريقة الأعداد الأصليَّة. أو ضعْ عدداً كما تشاء قاسماً، فاجمع الكسور لتتبع المطلوب، وادخلها في طريقة الأعداد المشتركة.



\textbf{أمَّا الأعداد المشتركة:} ضعْ أعداد المسألة واضرب البسط في المقام، فجميعُها تُسمَّى أعداداً مشتركة. اطلب المقاس المشترك الأعم، فلا تقسمْ واحداً واقسم الباقي، فتحصل على أعداد قانونها الأصلي، وادخلها في طريقة الأعداد الأصليَّة.



\textbf{أمَّا الأعداد المركَّبة:} إذا ظهرتْ في نهاية أعداد المسألة عشراتٌ فما فوق، فاطلب المقاس المشترك الأعم لجميع الأعداد، وأبقِ واحداً واقسم الباقي، فتتحصَّل على الأعداد الأصليَّة. ثم اطلب المقاس المشترك بين كلِّ زوجَين، فاقسم الفرديَّ دون الزوجيَّ، ثم اضرب الزوجيَّ؛ أو اقسم الزوجيَّ دون الفرديَّ، ولا تضرب الفرديَّ.

}



\parallelsubsection{Computing the Derivative Numbers}{حساب الأعداد المشتقَّة}



\paralleltext{%

\textbf{求定数}勿使两位见偶，勿使见一太多。见一多则借用繁，不欲借则任得一。



\textbf{以定相乘为衍母}，以各定约衍母，得各衍数。或列各定数于右行，各立天元一为子于左行，以母互乘子，亦得衍数。



Then each derivative number is reduced modulo its corresponding definitive number; the remainder is called the \textbf{odd number} ({\chfont 奇数})). Using the Dayan Method for Finding Unity ({\chfont 大衍求一术})), one finds the \textbf{multiplier} ({\chfont 乘率})) for each. Multiplying each multiplier by its derivative number gives the \textbf{use number} ({\chfont 用数})). The final answer is obtained by multiplying each remainder by its use number, summing, and reducing modulo the derivative mother ({\chfont 衍母})).%

}{%

\textbf{في طلب الأعداد المثبتة:} لا تدعْ رقمَين يظهران زوجيَّين، ولا تدعِ الواحدَ يكثر. إذا كثُر الواحد كثُرتِ الاستعارات، فإنْ لم تردْ الاستعارة فدعْهُ.



\textbf{اضربْ المثبتات بعضها في بعض فتنال الجذر المشترك} (衍母)، واقسمْ الجذر المشترك على كلِّ مثبتٍ فتنال الأعداد المشتقَّة (衍数). أو اصطفِّ المثبتات في الصفِّ الأيمن، وضع الواحد الأصليَّة ابنَ كلٍّ منها في الصفِّ الأيسر، فاضرب الأمَّهات متبادلاً في الأبناء، فتنال أيضاً الأعداد المشتقَّة.



ثم يُختزل كلُّ عددٍ مشتقٍّ على معيار عدده المثبت؛ والباقي يُسمَّى \textbf{العدد الفردي} (奇数). وبطريقة استخراج الوحدة (大衍求一术)، يُوجد \textbf{معامل الضرب} (乘率) لكلٍّ منها. ضربُ كلِّ معاملٍ في عدده المشتقِّ يُنتج \textbf{العدد المستعمل} (用数). تُستخرَج الإجابة النهائية بضرب كلِّ باقٍ في عدده المستعمل، وتجميعها، والاختزال على الجذر المشترك (衍母).

}



\parallelsubsection{The Dayan Method for Finding Unity (大衍求一术)}{طريقة استخراج الوحدة (大衍求一术)}



\paralleltext{%

\begin{cnquote}

大衍求一术云：置奇右上，定居右下。立天元一于左上。先以右行上下两位，以少除多，所得商数，乃递互乘归左行。须使右上末后奇一而止。乃验左上所得，以为乘率。或奇数已见单一者，便为乘率。

\end{cnquote}%

}{%

\begin{quote}

تقول طريقة استخراج الوحدة: ضعِ الفرديَّ في أعلى اليمين، والمثبَّتَ في أسفل اليمين. وضع الواحد الأصليَّة في أعلى اليسار. ابتدئْ بالصفِّ الأيمن من أعلى وأسفل، فاقسم الأصغر على الأكبر، فناتج القسمة يُضرب تبادليّاً في الصفِّ الأيسر. ويجب أن يصبح الفرديُّ في الأعلى في النهاية واحداً. ثم افحص ما حُصِلَ في الأعلى الأيسر، فذاك هو معامل الضرب. أو إذا كان الفرديُّ قد ظهر مفرداً واحداً، فهو معامل الضرب.

\end{quote}

}



\paralleltext{%

\textit{In modern terms:} Given coprime integers $a$ (the ``odd number'') and $m$ (the ``definitive number''), find $k$ such that:

\[k \cdot a \equiv 1 \pmod{m}\]

This is the modular multiplicative inverse of $a$ modulo $m$, computed by the Euclidean algorithm with back-substitution.%

}{%

\textit{بالمصطلحات الحديثة:} مُعطى عددان أوليَّان فيما بينهما $a$ (``العدد الفردي'') و $m$ (``العدد المثبَّت'')، أوجِد $k$ بحيث:

\[k \cdot a \equiv 1 \pmod{m}\]

وهذا هو المعكوس الضربي المعياري لـ $a$ على المعيار $m$، المُحتَسب بخوارزمية إقليدس مع الاستعاضة العكسية.

}



\newpage



%% ========================================================================

%% PROBLEM 1: DIVINATION BY YARROW STALKS

%% ========================================================================

\parallelsection{Problem 1: Shigua Fawei (蓍卦发微)}{المسألة الأولى: كشف خفايا القرعة بالسعف (蓍卦发微)}



\paralleltext{%

\cnlabel{问曰} 《易》曰：大衍之数五十，其用四十有九。又曰：分而为二以象两，挂一以象三，揲之以四以象四时。三变而成爻，十有八变而成卦。欲知所衍之术及其数各几何？%

}{%

\arlabel{المسألة:} يقول \textit{كتاب التغييرات}: عدد الطريقة الكبرى خمسون، وما يُستعمل منها تسعة وأربعون. ويقول أيضاً: اقسِمْها إلى اثنين لترمز إلى الاثنين [السماء والأرض]، وعلِّق واحداً لترمز إلى الثلاثة [السماء والأرض والإنسان]، واستقرِها بأربعة لترمز إلى الفصول الأربعة. ثلاثُ تحويلاتٍ تُنتج الخطَّ، وثماني عشرةَ تحويلةً تُنتج الرسمَ الكَمِّيَّ. نريد أن نعرف طريقة الاشتقاق وأعدادَها.

}



\paralleltext{%

\textit{The Book of Changes says: ``The number of the Great Derivation is fifty; of these, forty-nine are used.'' It also says: ``Divide into two to symbolize the Two [heaven and earth]; hang one to symbolize the Three; count off by fours to symbolize the Four Seasons.'' Three changes produce a line; eighteen changes produce a hexagram.}%

}{%

\textit{يقول كتاب التغييرات: ``عدد الطريقة الكبرى خمسون، ويُستعمل منها تسعة وأربعون.'' ويقول أيضاً: ``اقسِمْها إلى اثنين ترمز إلى الاثنين، وعلِّق واحداً يرمز إلى الثلاثة، واستقرِها بأربعة ترمز إلى الفصول الأربعة.'' ثلاثُ تغييراتٍ تُنتج الخط، وثماني عشرة تغييراً تُنتج الرسم الكَمِّي.}

}



\paralleltext{%

\cnlabel{答曰} 衍母十二，衍法三。



一元衍数二十四，二元衍数一十二，三元衍数八，四元衍数六。以上四位衍数，计五十一。



一揲用数一十二，二揲用数二十四，三揲用数四，四揲用数九。以上四位用数，计四十九。%

}{%

\arlabel{الجواب:} الجذر المشترك (衍母) اثنا عشر، وطريقة الاشتقاق ثلاثة.



العدد المشتقُّ الأوَّل أربعة وعشرون، والثاني اثنا عشر، والثالث ثمانية، والرابع ستة. مجموع هذه الأعداد المشتقَّة الأربعة واحد وخمسون.



عدد الاستعمال الأوَّل اثنا عشر، والثاني أربعة وعشرون، والثالث أربعة، والرابع تسعة. مجموع هذه الأعداد المستعملة الأربعة تسعة وأربعون.

}



\paralleltext{%

\cnlabel{术曰} 置诸元数，两两连环求等，约奇弗约偶。遍约毕，乃变元数皆曰定母，列右行。各立天元一为子，列左行。以诸定母互乘左行之子，各得名曰衍数。次以各定母满去衍数，各余名曰奇数。以奇数与定母用大衍术求一，得各乘率。以乘衍数，各得用数。验次所揲余几何，以其余数乘诸用数，并之，名曰总数。满衍母去之，不满为所求数。%

}{%

\arlabel{الطريقة:} ضع جميع الأعداد الأصليَّة، واطلب المقاس المشترك بين كلِّ زوجَين، فاقسم الفرديَّ دون الزوجيَّ. بعد القسمة الشاملة، حوِّل جميع الأعداد الأصليَّة إلى أمَّهات مثبتة، واصطفَّها في الصفِّ الأيمن. وضع لكلٍّ منها الواحد الأصليَّة ابناً، واصطفَّها في الصفِّ الأيسر. اضرب الأمَّهات المثبتة متبادلاً في أبناء الصفِّ الأيسر، فتُسمَّى كلُّها أعداداً مشتقَّة. ثم اختزل كلَّ عددٍ مشتقٍّ بكامل قسمته على أمِّه المثبتة، فالباقي من كلٍّ يُسمَّى عدداً فرديّاً. بالعدد الفردي والأمِّ المثبَّتة استعملْ الطريقة الكبرى لاستخراج الوحدة، فتحصل على معاملات الضرب. اضربها في الأعداد المشتقَّة، فتنال أعداداً مستعملة. افحص ما يبقى من الاستقراء، فاضرب تلك البواقي في أعداد الاستعمال، واجمعْها، فذلك يُسمَّى المجموع الكليَّ. اختزلْهُ بالجذر المشترك، فما لا يقبل الاختزال هو المطلوب.

}



\paralleltext{%

\cnlabel{草曰} 置一、二、三、四列右行，立天元一列左行。



以右行一、二、三、四互乘左行：一乘一得一，二乘一得二，三乘二得六，四乘六得二十四。各得衍数。



次以定母满法去衍数：以十二除二十四得二余零，一元衍数为二十四；以十二除一十二得一余零，二元衍数为一十二；以八除八得一余零，三元衍数为八；以六除六得一余零，四元衍数为六。



各满定母去之，得奇数。以大衍求一入之，得乘数：得二十三为乘率；得五十一为乘率；得七为乘率。



以乘率各乘衍数，得用数：一揲用数一十二；二揲用数二十四；三揲用数四；四揲用数九。以上四位用数，计四十九。%

}{%

\arlabel{الحل المفصل:} ضع واحداً واثنين وثلاثة وأربعة في الصفِّ الأيمن، والواحد الأصليَّة في الصفِّ الأيسر.



اضرب الصفَّ الأيمن (١، ٢، ٣، ٤) متبادلاً في الصفِّ الأيسر: واحدٌ في واحدٍ يساوي واحداً، اثنان في واحدٍ يساوي اثنين، ثلاثة في اثنين تساوي ستة، أربعة في ستة تساوي أربعة وعشرين. فتُسمَّى هذه الأعداد المشتقَّة.



ثم اختزل الأعداد المشتقَّة بكامل قسمتها على الأمهات المثبتة: اثنا عشر تقسيم أربعة وعشرين يساوي اثنين ولا باقي، فالعدد المشتقُّ الأوَّل أربعة وعشرون؛ اثنا عشر تقسيم اثنا عشر يساوي واحداً ولا باقي، فالثاني اثنا عشر؛ ثمانية تقسيم ثمانية يساوي واحداً ولا باقي، فالثالث ثمانية؛ ستة تقسيم ستة يساوي واحداً ولا باقي، فالرابع ستة.



اختزلْ كلَّ واحدٍ بكامل قسمته على أمه المثبتة، فتنال الأعداد الفرديَّة. ثم ادخلْ في استخراج الوحدة، فتنال معاملات الضرب: ثلاثة وعشرون معاملاً، وواحد وخمسون معاملاً، وسبعة معاملات.



اضربْ كلَّ معاملٍ في عدده المشتقَّ، فتنال أعداد الاستعمال: الاستقراء الأوَّل اثنا عشر، والثاني أربعة وعشرون، والثالث أربعة، والرابع تسعة. مجموع هذه الأعداد المستعملة الأربعة تسعة وأربعون.

}



\paralleltext{%

\textit{This problem applies the Dayan method to the classical yarrow-stalk divination procedure of the Book of Changes (I Ching), interpreting the numerological operations in terms of congruence arithmetic.}%

}{%

\textit{تطبِّق هذه المسألة الطريقة الكبرى على طقس الكهانة الكلاسيكي بسعف الكَفْ---من كتاب التغييرات (الإي-تشينغ)، فتفسِّر العمليات الرقمية بمصطلحات حساب التطابقات.}

}



\newpage



%% ========================================================================

%% PROBLEM 2: ANCIENT CALENDAR ACCUMULATION

%% ========================================================================

\parallelsection{Problem 2: Guli Huiji (古历会积)}{المسألة الثانية: جمع الحسابات الفلكية القديمة (古历会积)}



\paralleltext{%

\cnlabel{问曰} 问古历会积，欲求积年。以历法积元为问，历过无考，但据近世演纪之法，推而上之。%

}{%

\arlabel{المسألة:} تتعلَّق هذه المسألة بحساب السنوات المتراكمة (积年) في الأنظمة التقويميَّة القديمة. مُعطى بعض الدورات الفلكية، يجب إيجاد عددٍ يُحقِّق شروط تطابقٍ متعددةٍ في آنٍ واحد.

}



\paralleltext{%

\cnlabel{答曰} 答曰：积年若干。%

}{%

\arlabel{الجواب:} الإجابة: عدد السنوات المتراكمة هو العدد المطلوب.

}



\paralleltext{%

\cnlabel{术曰} 术曰：以大衍求之。置元历相关诸周期，连环求等，约奇弗约偶，得定母。诸定相乘为衍母。以各定约衍母得衍数。各满定母去之，余为奇数。大衍求一入之，得乘率。以乘衍数为用数。次置各余数乘用数，并之为总数。满衍母去之，不满为所求积年。%

}{%

\arlabel{الطريقة:} أوجِدْها بالطريقة الكبرى. ضع جميع الدورات المتعلِّقة بالتقويم الأصلي، فاطلب المقاس المشترك بين كلِّ زوجَين، واقسم الفرديَّ دون الزوجيَّ، فتنال الأمَّهات المثبتة. اضربْ المثبتات بعضها في بعض فتنال الجذر المشترك. اقسمْ الجذر المشترك على كلِّ مثبتةٍ فتنال الأعداد المشتقَّة. اختزلْ كلَّ عددٍ مشتقٍّ بأمه المثبتة، فالباقي هو العدد الفردي. ادخلْ في استخراج الوحدة، فتنال معاملات الضرب. اضربها في الأعداد المشتقَّة فتنال أعداد الاستعمال. ثم ضع كلَّ باقٍ في استقراء واضربه في عدد استعماله، فاجمعْها فتنال المجموع الكلي. اختزلْهُ بالجذر المشترك، فما لا يقبل الاختزال هو السنوات المتراكمة المطلوبة.

}



\paralleltext{%

\textit{This problem applies the Dayan method to astronomical calendar calculations, finding a common epoch that aligns multiple celestial cycles (such as the tropical year, synodic month, and sexagenary cycle).}%

}{%

\textit{تطبِّق هذه المسألة الطريقة الكبرى على الحسابات التقويميَّة الفلكية، فتجدهراً مشتركاً يُوافق بين دورات سماويَّة متعددة (مثل السنة المدارية، والشهر القمري الاقتراني، ودورة الستين). وكانت هذه من أبرز تطبيقات الطريقة الكبرى في الفلك الصيني التقليدي.}

}



\newpage



%% ========================================================================

%% PROBLEM 3: ESTIMATING EARTHWORKS

%% ========================================================================

\parallelsection{Problem 3: Tuiji Tugong (推计土功)}{المسألة الثالثة: تقدير أعمال التراب (推计土功)}



\paralleltext{%

\cnlabel{问曰} 四县（甲、乙、丙、丁）共同筑堤，根据各县的物力分配筑堤任务。\\

--- 甲县物力：138,600贯\\

--- 乙县物力：146,300贯\\

--- 丙县物力：192,500贯\\

--- 丁县物力：184,800贯\\

筑堤标准：每770贯物力分配一名劳动力，每人每天筑堤60平方尺。\\

进度情况：甲县和乙县已完成任务，丙县剩余51丈，丁县剩余18丈。%

}{%

\arlabel{المسألة:} أربع مقاطعات (أ، ب، ج، د) تبني سداً مشتركاً، وتوزَّع مهام البناء بحسب الموارد المادية لكل مقاطعة:\\

--- موارد مقاطعة أ: ١٣٨٬٦٠٠ كوان\\

--- موارد مقاطعة ب: ١٤٦٬٣٠٠ كوان\\

--- موارد مقاطعة ج: ١٩٢٬٥٠٠ كوان\\

--- موارد مقاطعة د: ١٨٤٬٨٠٠ كوان\\

معيار البناء: لكل ٧٧٠ كوان من الموارد عامل واحد، ولكل عامل ٦٠ قدماً مربَّعةً يومياً من بناء السد.\\

حالة التقدُّم: أكملتْ مقاطعتا أ وب مهامَهما، وتبقَّى لمقاطعة ج ٥١ زانغ، ولمقاطعة د ١٨ زانغ.

}



\paralleltext{%

\cnlabel{答曰} 堤的总长度为19里235步5尺。\\

各县所筑长度：\\

--- 甲县：1,260丈\\

--- 乙县：1,768步5尺6寸\\

--- 丙县：4里328步5尺6寸\\

--- 丁县：与前三县相应%

}{%

\arlabel{الجواب:} الطول الكلي للسد ١٩ لي و٢٣٥ بو و٥ تشي.\\

الأطوال المبنية من كل مقاطعة:\\

--- مقاطعة أ: ١٢٦٠ زانغ\\

--- مقاطعة ب: ١٧٦٨ بو و٥ تشي و٦ سون\\

--- مقاطعة ج: ٤ لي و٣٢٨ بو و٥ تشي و٦ سون\\

--- مقاطعة د: بالتناسب مع المقاطعات الثلاث الأولى

}



\paralleltext{%

\cnlabel{术曰} 术曰：以大衍求之。置各县物力，连环求等，约奇弗约偶，得定母。诸定相乘为衍母。以各定约衍母得衍数。各满定母去之，余为奇数。大衍求一入之，得乘率。以乘衍数为用数。次置各县余数乘用数，并之为总数。满衍母去之，不满为所求。%

}{%

\arlabel{الطريقة:} أوجِدْها بالطريقة الكبرى. ضعْ موارد كل مقاطعة، فاطلب المقاس المشترك بين كلِّ زوجَين، واقسم الفرديَّ دون الزوجيَّ، فتنال الأمَّهات المثبتة. اضربْ المثبتات بعضها في بعض فتنال الجذر المشترك. اقسمْ الجذر المشترك على كل مثبتة فتنال الأعداد المشتقَّة. اختزلْ كلَّ عددٍ مشتقٍّ بأمه المثبتة، فالباقي هو العدد الفردي. ادخلْ في استخراج الوحدة، فتنال معاملات الضرب. اضربها في الأعداد المشتقَّة فتنال أعداد الاستعمال. ثم ضع باقي كل مقاطعة واضربه في عدد استعماله، فاجمعْها فتنال المجموع الكلي. اختزلْهُ بالجذر المشترك، فما لا يقبل الاختزال هو المطلوب.

}



\paralleltext{%

\cnlabel{草曰} 计算各县每日筑堤长度：\\

--- 甲县：54丈 / 乙县：57丈 / 丙县：75丈 / 丁县：72丈\\

使用大衍求一术：\\

--- 衍母：102,600\\

--- 甲县衍数：3,800 / 乙县衍数：5,400 / 丙县衍数：4,140 / 丁县衍数：12,825\\

乘率：甲23 / 乙5 / 丙19 / 丁1\\

用数：甲17,400 / 乙27,000 / 丙77,976 / 丁12,825\\

验证：堤的总长度19里235步5尺。%

}{%

\arlabel{الحل المفصل:} احسب الأطوال اليوميَّة لبناء السد من كل مقاطعة:\\

--- مقاطعة أ: ٥٤ زانغ / مقاطعة ب: ٥٧ زانغ / مقاطعة ج: ٧٥ زانغ / مقاطعة د: ٧٢ زانغ\\

بطريقة استخراج الوحدة:\\

--- الجذر المشترك: ١٠٢٬٦٠٠\\

--- العدد المشتقُّ لمقاطعة أ: ٣٬٨٠٠ / لمقاطعة ب: ٥٬٤٠٠ / لمقاطعة ج: ٤٬١٤٠ / لمقاطعة د: ١٢٬٨٢٥\\

معاملات الضرب: أ ٢٣ / ب ٥ / ج ١٩ / د ١\\

أعداد الاستعمال: أ ١٧٬٤٠٠ / ب ٢٧٬٠٠٠ / ج ٧٧٬٩٧٦ / د ١٢٬٨٢٥\\

التحقق: الطول الكلي للسد ١٩ لي و٢٣٥ بو و٥ تشي.

}



\newpage



%% ========================================================================

%% PROBLEM 4: CALCULATING TREASURY FUNDS

%% ========================================================================

\parallelsection{Problem 4: Tuiku E Qian (推库额钱)}{المسألة الرابعة: حساب أموال الخزينة (推库额钱)}



\paralleltext{%

\cnlabel{问曰} 问有外邑七库，日纳息足钱适等，递年成贯整纳。近缘见钱希少，听各库照当处市陌准解旧会。其甲库有零钱一十文，丁、庚二库各零四文，戊库零六文，余库无零钱。甲库所在市陌一十二文，递减一文至庚库而止。欲求诸库日息原纳足钱、展省及今纳旧会，并大小月分各几何？%

}{%

\arlabel{المسألة:} سُئل عن سبع خزائن في بلدات خارجية، كلها تدفع الفوائد اليوميَّة بالنقود الكاملة بالتساوي، وتدفع سنوياً بالكوان الكاملة. وقد قلَّ النقد مؤخَّراً، فسُمح لكل خزينة بتحويل قيمتها وفق سعر السوق المحلِّي إلى النقود القديمة. الخزينة أ لديها باقٍ عشرة وين، والخزينتان د وج لكل منهما أربعة وين، والخزينة هـ لديها ستة وين، والباقي لا بواقي. سعر السوق عند الخزينة أ اثنا عشر وين، ينقص واحداً حتى الخزينة ج. نريد معرفة: النقود الكاملة الأصلية اليومية، والقيمة المحوَّلة، والنقود القديمة الحالية، وحساب الأشهر الكبيرة والصغيرة.

}



\paralleltext{%

\cnlabel{答曰} 诸库纳日息足钱二十贯九百五十文。展省三十五贯文。\\

各库旧会：\\

--- 甲库日息旧会：二百二十四贯五百一十文\\

--- 乙库日息旧会：二百四十五贯文\\

--- 丙库日息旧会：二百六十九贯五百文\\

--- 庚库日息旧会：四百四十九贯一百四文%

}{%

\arlabel{الجواب:} النقود الكاملة التي تدفعها الخزائن يومياً: عشرون كوان وتسعمئة وخمسون ويناً. القيمة المحوَّلة: خمسة وثلاثون كوان.\\

النقود القديمة لكل خزينة:\\

--- خزينة أ يومياً: مئتان وأربعة وعشرون كوان وخمسمئة وعشرة وين\\

--- خزينة ب يومياً: مئتان وخمسة وأربعون كوان\\

--- خزينة ج يومياً: مئتان وتسعة وستون كوان وخمسمئة وين\\

--- خزينة ج (السابعة) يومياً: أربعمائة وتسعة وأربعون كوان ومئة وأربعة وين

}



\paralleltext{%

\cnlabel{术曰} 术曰：以大衍求之。置甲库市陌，以递减数减之，各得诸库原陌。连环求等，约奇弗约偶，得定母。诸定相乘为衍母。以定约衍母得衍数。各满定母去之，余为奇数。大衍求一入之，得乘率。以乘衍数为用数。次置有零文库零钱数，乘本用数，并为总数。满衍母去之，不满为诸库日息足钱。%

}{%

\arlabel{الطريقة:} أوجِدْها بالطريقة الكبرى. ضعْ سعر السوق للخزينة أ، فاطرح منه مقدار النقصان، فتنال أسعار السوق الأصلية لجميع الخزائن. اطلب المقاس المشترك بين كلِّ زوجَين، واقسم الفرديَّ دون الزوجيَّ، فتنال الأمَّهات المثبتة. اضربْ المثبتات بعضها في بعض فتنال الجذر المشترك. اقسمْ الجذر المشترك على المثبتات فتنال الأعداد المشتقَّة. اختزلْ كلَّ عددٍ مشتقٍّ بأمه المثبتة، فالباقي هو العدد الفردي. ادخلْ في استخراج الوحدة، فتنال معاملات الضرب. اضربها في الأعداد المشتقَّة فتنال أعداد الاستعمال. ثم ضعْ بواقي الخزائن التي لديها بواقي، واضربها في أعداد استعمالها، فاجمعْها فتنال المجموع الكلي. اختزلْهُ بالجذر المشترك، فما لا يقبل الاختزال هو النقود الكاملة اليومية لجميع الخزائن.

}



\paralleltext{%

\cnlabel{草曰} 置甲库市陌一十二，递减一得一十一为乙库陌，十为丙库陌，九为丁库陌，八为戊库陌，七为己库陌，六为庚库陌。得诸库原陌。\\

连环求等约讫：甲得一 / 乙得十一 / 丙得五 / 丁得九 / 戊得八 / 己得七 / 庚得一。各为定母。\\

先以诸定相乘，得27,720为衍母。\\

次以诸定互乘诸子：甲27,720 / 乙2,520 / 丙5,544 / 丁3,080 / 戊3,465 / 己3,960 / 庚27,720。各为衍数。%

}{%

\arlabel{الحل المفصل:} ضع سعر السوق للخزينة أ اثني عشر، فانقص واحداً فتنال أحد عشر لخزينة ب، وعشرة لخزينة ج، وتسعة لخزينة د، وثمانية لخزينة هـ، وسبعة لخزينة و، وستة لخزينة ج. فتنال أسعار السوق الأصلية.\\

بعد طلب المقاس المشترك والقسمة: أ ينال واحداً / ب أحد عشر / ج خمسة / د تسعة / هـ ثمانية / وسبعة / ج واحداً. فكلٌّ منها أمٌّ مثبتة.\\

أولاً اضرب جميع المثبتات بعضها في بعض، فتنال ٢٧٬٧٢٠ جذراً مشتركاً.\\

ثم اضرب المثبتات متبادلاً في الأبناء: أ ٢٧٬٧٢٠ / ب ٢٬٥٢٠ / ج ٥٬٥٤٤ / د ٣٬٠٨٠ / هـ ٣٬٤٦٥ / و ٣٬٩٦٠ / ج ٢٧٬٧٢٠. فكلٌّ منها عددٌ مشتقٌّ.

}



\newpage



%% ========================================================================

%% PROBLEM 5: GRAIN DISTRIBUTION

%% ========================================================================

\parallelsection{Problem 5: Fentiao Tuiyuan (分粜推原)}{المسألة الخامسة: تتبُّع أصل توزيع الحبوب (分粜推原)}



\paralleltext{%

\cnlabel{问曰} 有兄弟三人，各将收获之米送往不同的地方：\\

--- 甲将米送往官场，余3斗2升\\

--- 乙将米送往安吉乡，余7斗\\

--- 丙将米送往平江，余3斗\\

官斛每石八斗三升，安吉斛每石一石一斗，平江斛每石一石三斗五升。欲求三人共收获米数以及各自分米数几何？%

}{%

\arlabel{المسألة:} ثلاثة إخوة، كلٌّ منهم يرسل محصولَه من الأرز إلى مكان مختلف:\\

--- أ يرسل الأرز إلى الديوان الحكومي، يبقى ٣ دول و٢ لتر\\

--- ب يرسل الأرز إلى ضاحية أنجي، يبقى ٧ دول\\

--- ج يرسل الأرز إلى بينغجيانغ، يبقى ٣ دول\\

الكيل الحكومي ٨ دول و٣ لترات للشِّع، وكيل أنجي شِّع واحد ودول واحد، وكيل بينغجيانغ شِّع واحد و٣ دول و٥ لترات. نريد معرفة إجمالي محصول الأرز الذي جنوه الثلاثة، وكمية كلٍّ منهم.

}



\paralleltext{%

\cnlabel{答曰} 三人共收获米738石。\\

各分米数：\\

--- 甲：296石，余3斗2升\\

--- 乙：223石，余7斗\\

--- 丙：182石，余3斗%

}{%

\arlabel{الجواب:} مجموع ما جنوه الثلاثة من الأرز ٧٣٨ شِّع.\\

الكمية الموزَّعة لكلٍّ منهم:\\

--- أ: ٢٩٦ شِّع، يبقى ٣ دول و٢ لتر\\

--- ب: ٢٢٣ شِّع، يبقى ٧ دول\\

--- ج: ١٨٢ شِّع، يبقى ٣ دول

}



\paralleltext{%

\cnlabel{术曰} 术曰：以大衍求之。置各斛率（单位：升）：官斛83升，安吉斛110升，平江斛135升。连环求等，约奇弗约偶，得定母。诸定相乘为衍母。以各定约衍母得衍数。各满定母去之，余为奇数。大衍求一入之，得乘率。以乘衍数为用数。次置各余数乘用数，并之为总数。满衍母去之，不满为所求率数。%

}{%

\arlabel{الطريقة:} أوجِدْها بالطريقة الكبرى. ضع معدَّلات الكيل (وحدة: لتر): الكيل الحكومي ٨٣ لتراً، وكيل أنجي ١١٠ لتراً، وكيل بينغجيانغ ١٣٥ لتراً. اطلب المقاس المشترك بين كلِّ زوجَين، واقسم الفرديَّ دون الزوجيَّ، فتنال الأمَّهات المثبتة. اضربْ المثبتات بعضها في بعض فتنال الجذر المشترك. اقسمْ الجذر المشترك على كل مثبتة فتنال الأعداد المشتقَّة. اختزلْ كلَّ عددٍ مشتقٍّ بأمه المثبتة، فالباقي هو العدد الفردي. ادخلْ في استخراج الوحدة، فتنال معاملات الضرب. اضربها في الأعداد المشتقَّة فتنال أعداد الاستعمال. ثم ضع كلَّ باقٍ واضربه في عدد استعماله، فاجمعْها فتنال المجموع الكلي. اختزلْهُ بالجذر المشترك، فما لا يقبل الاختزال هو المطلوب.

}



\paralleltext{%

\cnlabel{草曰} 置官斛83升、安吉斛110升、平江斛135升，连环求等。求总等：得总等1，不约。连环求等，得定母：246,510。\\

各斛衍数：官斛2,970 / 安吉2,241 / 平江1,826\\

以大衍求一术，得乘率：官斛51 / 安吉7 / 平江23\\

以乘率乘衍数得用数：官斛151,470 / 安吉15,687 / 平江41,998\\

展算得：分米246石，以兄弟三人因之，得738石为共米。%

}{%

\arlabel{الحل المفصل:} ضع الكيل الحكومي ٨٣ لتراً وكيل أنجي ١١٠ لترات وكيل بينغجيانغ ١٣٥ لتراً، فاطلب المقاس المشترك بين كلِّ زوجَين. اطلب المقاس المشترك الأعم: ينال واحداً، فلا تقسمْ. بطلب المقاس المشترك بين كلِّ زوجَين تنال الأمَّهات المثبتة: ٢٤٦٬٥١٠.\\

الأعداد المشتقَّة لكل كيل: الحكومي ٢٬٩٧٠ / أنجي ٢٬٢٤١ / بينغجيانغ ١٬٨٢٦\\

بطريقة استخراج الوحدة، تنال معاملات الضرب: الحكومي ٥١ / أنجي ٧ / بينغجيانغ ٢٣\\

اضربْ معاملات الضرب في الأعداد المشتقَّة فتنال أعداد الاستعمال: الحكومي ١٥١٬٤٧٠ / أنجي ١٥٬٦٨٧ / بينغجيانغ ٤١٬٩٩٨\\

بالتوسيع: توزيع الأرز ٢٤٦ شِّع، مضروباً في ثلاثة إخوة، فينال ٧٣٨ شِّعاً إجمالاً.

}



\newpage



%% ========================================================================

%% PROBLEM 6: JOURNEY DISTANCE CALCULATION

%% ========================================================================

\parallelsection{Problem 6: Chengxing Jidi (程行计地)}{المسألة السادسة: حساب مسافة الرحلة (程行计地)}



\paralleltext{%

\cnlabel{问曰} 军师派遣三名急足（甲、乙、丙）分别前往都下报捷。\\

--- 甲每天行300里\\

--- 乙每天行240里\\

--- 丙每天行180里\\

甲提前数日申未到，乙后数日未正到，丙今日辰末到。欲知从军前到都下的距离以及三人各自行走的天数几何？%

}{%

\arlabel{المسألة:} أرسل القائد ثلاثة رسل سريعين (أ، ب، ج) إلى العاصمة للإبلاغ عن النصر.\\

--- أ يسير ٣٠٠ لي يومياً\\

--- ب يسير ٢٤٠ لي يومياً\\

--- ج يسير ١٨٠ لي يومياً\\

وصل أ قبل موعده بعدة أيَّام في ساعة الشين، ووصل ب بعد موعده بعدة أيَّام في ساعة الواي، وج وصل اليوم في نهاية ساعة التشين. نريد معرفة المسافة من المعسكر إلى العاصمة، وعدد الأيَّام التي سارها كلٌّ منهم.

}



\paralleltext{%

\cnlabel{答曰} 从军前到都下的距离为3,300里。\\

--- 甲行11天\\

--- 乙行13天4时\\

--- 丙行18天2时%

}{%

\arlabel{الجواب:} المسافة من المعسكر إلى العاصمة ٣٬٣٠٠ لي.\\

--- أ سار ١١ يوماً\\

--- ب سار ١٣ يوماً و٤ ساعات\\

--- ج سار ١٨ يوماً وساعتين

}



\paralleltext{%

\cnlabel{术曰} 术曰：以大衍求之。置三人日行里数，连环求等，约奇弗约偶，得定母。诸定相乘为衍母。以各定约衍母得衍数。各满定母去之，余为奇数。大衍求一入之，得乘率。以乘衍数为用数。次置各余数乘用数，并之为总数。满衍母去之，不满为所求距离。%

}{%

\arlabel{الطريقة:} أوجِدْها بالطريقة الكبرى. ضعْ ما يسيره كلُّ رجلٍ يومياً من اللي، فاطلب المقاس المشترك بين كلِّ زوجَين، واقسم الفرديَّ دون الزوجيَّ، فتنال الأمَّهات المثبتة. اضربْ المثبتات بعضها في بعض فتنال الجذر المشترك. اقسمْ الجذر المشترك على كل مثبتة فتنال الأعداد المشتقَّة. اختزلْ كلَّ عددٍ مشتقٍّ بأمه المثبتة، فالباقي هو العدد الفردي. ادخلْ في استخراج الوحدة، فتنال معاملات الضرب. اضربها في الأعداد المشتقَّة فتنال أعداد الاستعمال. ثم ضع كلَّ باقٍ واضربه في عدد استعماله، فاجمعْها فتنال المجموع الكلي. اختزلْهُ بالجذر المشترك، فما لا يقبل الاختزال هو المسافة المطلوبة.

}



\paralleltext{%

\cnlabel{草曰} 置甲300里、乙240里、丙180里，求总等，连环求等。计算得乘率：甲4 / 乙1 / 丙7。通过乘率乘以各自的用数，得到总用数，再通过总用数除以衍母，得到最终结果。\\

验证：\\

--- 3,300 $\div$ 300 = 11天\\

--- 3,300 $\div$ 240 = 13天余120里 = 13天4时（1天=240里，120里=半天=4时）\\

--- 3,300 $\div$ 180 = 18天余60里 = 18天2时。合问。%

}{%

\arlabel{الحل المفصل:} ضع ٣٠٠ لي لـ أ، و٢٤٠ لي لـ ب، و١٨٠ لي لـ ج، فاطلب المقاس المشترك الأعم والمقاس المشترك بين كلِّ زوجَين. يُحتَسب فتنال معاملات الضرب: أ ٤ / ب ١ / ج ٧. بضرب معاملات الضرب في أعداد استعمالها تنال أعداد الاستعمال الكلية، ثم بقسمة المجموع الكلي على الجذر المشترك يُحصَل على النتيجة النهائية.\\

التحقق:\\

--- ٣٬٣٠٠ $\div$ ٣٠٠ = ١١ يوماً\\

--- ٣٬٣٠٠ $\div$ ٢٤٠ = ١٣ يوماً و١٢٠ لي باقية = ١٣ يوماً و٤ ساعات (١ يوم = ٢٤٠ لي، ١٢٠ لي = نصف يوم = ٤ ساعات)\\

--- ٣٬٣٠٠ $\div$ ١٨٠ = ١٨ يوماً و٦٠ لي باقية = ١٨ يوماً وساعتين. هذا يُحقِّق المطلوب.

}



\newpage



%% ========================================================================

%% PROBLEM 7: CATCHING UP ON A JOURNEY

%% ========================================================================

\parallelsection{Problem 7: Chengxing Xiangji (程行相及)}{المسألة السابعة: اللحاق في الرحلة (程行相及)}



\paralleltext{%

\cnlabel{问曰} 问：甲、乙、丙三人行程相关。甲先行若干日，乙后行，丙又后行。各以不同速度行进。求某处三人相遇之时间及各地距离。%

}{%

\arlabel{المسألة:} ثلاثة مسافرون (أ، ب، ج) يسيرون على مسارٍ واحدٍ. أ يسبق ببضعة أيَّام، ثم ب يسير، ثم ج يسير بعده. كلٌّ يسير بسرعة مختلفة. نريد معرفة: الوقت والمكان الذي يلتقون فيه.

}



\paralleltext{%

\textit{This problem involves multiple travelers with different speeds and starting times, requiring the use of the Dayan method to solve a system of congruences for finding meeting times and distances.}%

}{%

\textit{تتضمَّن هذه المسألة مسافرين متعدِّدين بسرعات مختلفة وأوقات انطلاق مختلفة، وتتطلَّب استخدام الطريقة الكبرى لحل نظام تطابقات لإيجاد أوقات اللقاء والمسافات.}

}



\paralleltext{%

\cnlabel{答曰} 答曰：所求距离及相关数值若干。%

}{%

\arlabel{الجواب:} الإجابة: المسافة المطلوبة والقيم ذات الصلة هي العدد المطلوب.

}



\paralleltext{%

\cnlabel{术曰} 术曰：以大衍求之。置各人行程速度及时间差，连环求等，约奇弗约偶，得定母。诸定相乘为衍母。以各定约衍母得衍数。各满定母去之，余为奇数。大衍求一入之，得乘率。以乘衍数为用数。次置各余数乘用数，并之为总数。满衍母去之，不满为所求。%

}{%

\arlabel{الطريقة:} أوجِدْها بالطريقة الكبرى. ضعْ سرعة كلِّ مسافر وفرق الوقت، فاطلب المقاس المشترك بين كلِّ زوجَين، واقسم الفرديَّ دون الزوجيَّ، فتنال الأمَّهات المثبتة. اضربْ المثبتات بعضها في بعض فتنال الجذر المشترك. اقسمْ الجذر المشترك على كل مثبتة فتنال الأعداد المشتقَّة. اختزلْ كلَّ عددٍ مشتقٍّ بأمه المثبتة، فالباقي هو العدد الفردي. ادخلْ في استخراج الوحدة، فتنال معاملات الضرب. اضربها في الأعداد المشتقَّة فتنال أعداد الاستعمال. ثم ضع كلَّ باقٍ واضربه في عدد استعماله، فاجمعْها فتنال المجموع الكلي. اختزلْهُ بالجذر المشترك، فما لا يقبل الاختزال هو المطلوب.

}



\newpage



%% ========================================================================

%% PROBLEM 8: FINDING DIMENSIONS FROM VOLUME

%% ========================================================================

\parallelsection{Problem 8: Jichi Xunyuan (积尺寻源)}{المسألة الثامنة: البحث عن المصدر من الحجم (积尺寻源)}



\paralleltext{%

\cnlabel{问曰} 问欲砌基一段，见管大小方砖六门城砖四色。令匠取便，或平或侧，只用一色砖砌，须要适足。匠以砖量地计料，称用：\\

--- 大方：广多六寸，深少六寸\\

--- 小方：广多二寸，深少三寸\\

--- 城砖：长广多三寸，深少一寸\\

--- 六门砖：长广多三寸，深多一寸\\

皆不合匝，未免修破砖料裨补。其四色砖尺寸：\\

--- 大方：方一尺三寸\\

--- 小方：方一尺一寸\\

--- 城砖：长一尺二寸，阔六寸，厚二寸五分\\

--- 六门砖：长一尺，阔五寸，厚二寸\\

欲知基深广几何？%

}{%

\arlabel{المسألة:} يُريد بناء أساسٍ جزئيٍّ، ويملك لبنات مربَّعة كبيرة وصغيرة، ولبنات سورية، وأربعة أنواع من لبنات البوَّابات. يأمر النجار بالمرونة في استخدامها، إمَّا مسطَّحة أو جانبيَّة، بنوعٍ واحدٍ فقط، ويجب أن تكون كافيةً تماماً. يقيس النجار الأرض بالطوب ويحسب المواد:\\

--- المربَّع الكبير: العرض يزيد ٦ سون، العمق يقل ٦ سون\\

--- المربَّع الصغير: العرض يزيد ٢ سون، العمق يقل ٣ سون\\

--- اللبنة السوريَّة: الطول والعرض يزيدان ٣ سون، العمق يقل سوناً واحداً\\

--- لبنة البوَّابات: الطول والعرض يزيدان ٣ سون، العمق يزيد سوناً واحداً\\

كلُّها غير كافية، فيضطر لإصلاح وكسر اللبنات لتعويض النقص. أبعاد اللبنات الأربعة:\\

--- المربَّع الكبير: ضلع ١ تشي و٣ سون\\

--- المربَّع الصغير: ضلع ١ تشي وسون\\

--- اللبنة السوريَّة: طول ١ تشي و٢ سون، وعرض ٦ سون، وسمك ٢ سون ونصف\\

--- لبنة البوَّابات: طول ١ تشي، وعرض ٥ سون، وسمك ٢ سون\\

نريد معرفة عمق وعرض الأساس.

}



\paralleltext{%

\cnlabel{答曰} 答曰：深三丈七尺一寸，广一丈二尺三寸。%

}{%

\arlabel{الجواب:} العمق ٣ جوانغ و٧ تشي وسون واحد، والعرض ١ جوانغ و٢ تشي و٣ سون.

}



\paralleltext{%

\textit{This is a classic problem of finding a common multiple with given remainders. The dimensions of the foundation must be such that when measured with each type of brick (in different orientations), there are specific remainders. The Dayan method is used to find such dimensions.}%

}{%

\textit{هذه مسألة كلاسيكية في إيجاد المضاعف المشترك ببواقٍ معطاة. يجب أن تكون أبعاد الأساس بحيث عند قياسها بكل نوع من الطوب (في اتجاهات مختلفة)، تكون هناك بواقٍ محدَّدة. وتُستخدَم الطريقة الكبرى لإيجاد هذه الأبعاد.}

}



\paralleltext{%

\cnlabel{术曰} 术曰：以大衍求之。置四色砖尺寸及所余尺寸，连环求等，约奇弗约偶，得定母。诸定相乘为衍母。以各定约衍母得衍数。各满定母去之，余为奇数。大衍求一入之，得乘率。以乘衍数为用数。次置各余数乘用数，并之为总数。满衍母去之，不满为基之深广。%

}{%

\arlabel{الطريقة:} أوجِدْها بالطريقة الكبرى. ضعْ أبعاد اللبنات الأربعة والأبعاد الباقية، فاطلب المقاس المشترك بين كلِّ زوجَين، واقسم الفرديَّ دون الزوجيَّ، فتنال الأمَّهات المثبتة. اضربْ المثبتات بعضها في بعض فتنال الجذر المشترك. اقسمْ الجذر المشترك على كل مثبتة فتنال الأعداد المشتقَّة. اختزلْ كلَّ عددٍ مشتقٍّ بأمه المثبتة، فالباقي هو العدد الفردي. ادخلْ في استخراج الوحدة، فتنال معاملات الضرب. اضربها في الأعداد المشتقَّة فتنال أعداد الاستعمال. ثم ضع كلَّ باقٍ واضربه في عدد استعماله، فاجمعْها فتنال المجموع الكلي. اختزلْهُ بالجذر المشترك، فما لا يقبل الاختزال هو عمق وعرض الأساس.

}



\newpage



%% ========================================================================

%% PROBLEM 9: CALCULATING RICE REMAINDERS

%% ========================================================================

\parallelsection{Problem 9: Yumi Tuishu (余米推数)}{المسألة التاسعة: حساب بقايا الأرز (余米推数)}



\paralleltext{%

\cnlabel{问曰} 问：有米若干，以不同容量的容器量之，各余若干。欲求米之总数。%

}{%

\arlabel{المسألة:} سُئل عن كمية من الأرز، تُقاس بأوعية مختلفة السعة، ولكل منها باقٍ معيَّن. نريد معرفة العدد الكلي من الأرز.

}



\paralleltext{%

\textit{This problem involves finding the total number of units of rice when measured with containers of different capacities, leaving different remainders each time. It is the classic ``unknown quantity'' problem that leads directly to the Chinese Remainder Theorem.}%

}{%

\textit{تتعلَّق هذه المسألة بإيجاد العدد الكلي لوحدات الأرز عند قياسها بأوعية مختلفة السعة، تاركةً بواقٍ مختلفة في كل مرة. وهي مسألة ``الكمية المجهولة'' الكلاسيكية التي تقود مباشرة إلى نظرية الباقي الصينية.}

}



\paralleltext{%

\cnlabel{答曰} 答曰：所求米总数若干。%

}{%

\arlabel{الجواب:} الإجابة: العدد الكلي المطلوب من الأرز هو العدد المطلوب.

}



\paralleltext{%

\cnlabel{术曰} 术曰：以大衍求之。置各容器容量为问数，连环求等，约奇弗约偶，得定母。诸定相乘为衍母。以各定约衍母得衍数。各满定母去之，余为奇数。大衍求一入之，得乘率。以乘衍数为用数。次置各余数乘用数，并之为总数。满衍母去之，不满为所求米数。%

}{%

\arlabel{الطريقة:} أوجِدْها بالطريقة الكبرى. ضعْ سعة كلِّ وعاءٍ كعددٍ للمسألة، فاطلب المقاس المشترك بين كلِّ زوجَين، واقسم الفرديَّ دون الزوجيَّ، فتنال الأمَّهات المثبتة. اضربْ المثبتات بعضها في بعض فتنال الجذر المشترك. اقسمْ الجذر المشترك على كل مثبتة فتنال الأعداد المشتقَّة. اختزلْ كلَّ عددٍ مشتقٍّ بأمه المثبتة، فالباقي هو العدد الفردي. ادخلْ في استخراج الوحدة، فتنال معاملات الضرب. اضربها في الأعداد المشتقَّة فتنال أعداد الاستعمال. ثم ضع كلَّ باقٍ واضربه في عدد استعماله، فاجمعْها فتنال المجموع الكلي. اختزلْهُ بالجذر المشترك، فما لا يقبل الاختزال هو عدد الأرز المطلوب.

}



\paralleltext{%

\textit{This problem is structurally identical to the famous ``problem of the unknown quantity'' from the \textit{Sunzi Suanjing} ({\chfont 孙子算经})), which asks: ``We have things of which we do not know the number; if we count them by threes, the remainder is 2; if by fives, the remainder is 3; if by sevens, the remainder is 2. How many things are there?'' Qin Jiushao's general method solves all problems of this type systematically.}%

}{%

\textit{هذه المسألة متماثلة البنية مع مسألة ``الكمية المجهولة'' الشهيرة من \textit{كتاب سون-تسي في الحساب} (孙子算经)، والتي تسأل: ``لدينا أشياء لا نعرف عددها؛ إذا عددناها بثلاثات يبقى اثنان، وإذا بخمسات يبقى ثلاثة، وإذا بسبعات يبقى اثنان. كم عددها؟'' الطريقة العامة لتشين جيوشاو تحل جميع مسائل هذا النوع منهجياً.}

}



\newpage



%% ========================================================================

%% MATHEMATICAL ANALYSIS

%% ========================================================================

\parallelsection{Mathematical Analysis of the Dayan Method}{التحليل الرياضي للطريقة الكبرى}



\parallelsubsection{The Chinese Remainder Theorem}{نظرية الباقي الصينية}



\paralleltext{%

The Dayan method solves systems of simultaneous linear congruences. Given a set of pairwise coprime moduli $m_1, m_2, \ldots, m_n$ and remainders $r_1, r_2, \ldots, r_n$, find $N$ such that:

\begin{align*}

N &\equiv r_1 \pmod{m_1} \\

N &\equiv r_2 \pmod{m_2} \\

&\vdots \\

N &\equiv r_n \pmod{m_n}

\end{align*}%

}{%

تحلُّ الطريقة الكبرى أنظمة التطابقات الخطيَّة المتزامنة. مُعطى مجموعة من المعايير أوليَّة فيما بينها $m_1, m_2, \ldots, m_n$ وبواقٍ $r_1, r_2, \ldots, r_n$, أوجِد $N$ بحيث:

\begin{align*}

N &\equiv r_1 \pmod{m_1} \\

N &\equiv r_2 \pmod{m_2} \\

&\vdots \\

N &\equiv r_n \pmod{m_n}

\end{align*}

}



\parallelsubsection{Qin Jiushao's Algorithm}{خوارزمية تشين جيوشاو}



\paralleltext{%

\textbf{Step 1: Definitive Numbers (求定数).} Given the original numbers (元数), reduce them pairwise using the greatest common divisor, keeping them coprime. The rule ``约奇弗约偶'' (reduce the odd one, not the even one) ensures the product will have unique factorization.%

}{%

\textbf{الخطوة الأولى: الأعداد المثبتة (求定数).} مُعطى الأعداد الأصليَّة (元数)، اختزلها بأزواج باستخدام المقسوم المشترك الأعظم، مع الحفاظ على أوليَّتها فيما بينها. قاعدة ``اقسم الفرديَّ دون الزوجيَّ'' (约奇弗约偶) تضمن أن يكون الجذر ذا تحليلٍ وحيد.

}



\paralleltext{%

\textbf{Step 2: Derivative Mother (求衍母).} Compute $M = m_1 \times m_2 \times \cdots \times m_n$.%

}{%

\textbf{الخطوة الثانية: الجذر المشترك (求衍母).} احسب $M = m_1 \times m_2 \times \cdots \times m_n$.

}



\paralleltext{%

\textbf{Step 3: Derivative Numbers (求衍数).} Compute $M_i = M / m_i$ for each $i$.%

}{%

\textbf{الخطوة الثالثة: الأعداد المشتقَّة (求衍数).} احسب $M_i = M / m_i$ لكل $i$.

}



\paralleltext{%

\textbf{Step 4: Odd Numbers (求奇数).} Compute $g_i = M_i \bmod m_i$.%

}{%

\textbf{الخطوة الرابعة: الأعداد الفرديَّة (求奇数).} احسب $g_i = M_i \bmod m_i$.

}



\paralleltext{%

\textbf{Step 5: Finding Unity (大衍求一术).} For each $i$, find $k_i$ such that:

\[k_i \cdot g_i \equiv 1 \pmod{m_i}\]

This is the modular multiplicative inverse.%

}{%

\textbf{الخطوة الخامسة: استخراج الوحدة (大衍求一术).} لكل $i$, أوجِد $k_i$ بحيث:

\[k_i \cdot g_i \equiv 1 \pmod{m_i}\]

وهذا هو المعكوس الضربي المعياري.

}



\paralleltext{%

\textbf{Step 6: Use Numbers (求用数).} Compute $u_i = k_i \cdot M_i$.%

}{%

\textbf{الخطوة السادسة: الأعداد المستعملة (求用数).} احسب $u_i = k_i \cdot M_i$.

}



\paralleltext{%

\textbf{Step 7: Final Computation.} The solution is:

\[N = \sum_{i=1}^{n} r_i \cdot u_i \pmod{M}\]%

}{%

\textbf{الخطوة السابعة: الحساب النهائي.} الحل هو:

\[N = \sum_{i=1}^{n} r_i \cdot u_i \pmod{M}\]

}



\parallelsubsection{Historical Significance}{الأهمية التاريخية}



\paralleltext{%

Qin Jiushao's method predated similar European work by several centuries. The algorithm for finding modular inverses (大衍求一术) is essentially the extended Euclidean algorithm. The method was later applied by Guo Shoujing (郭守敬) to astronomy, by Li Ye (李冶) to geometry, and was transmitted to Europe where it became known as the ``Chinese Remainder Theorem.''%

}{%

سبقت طريقة تشين جيوشاو الأعمال الأوروبية المماثلة بقرون عديدة. خوارزمية إيجاد المعكوسات المعيارية (大衍求一术) هي في جوهرها خوارزمية إقليدس الموسَّعة. طُبِّقت الطريقة لاحقاً من قبل غو شوجينغ (郭守敬) في الفلك، ومن قبل لي يه (李冶) في الهندسة، ثم نُقلت إلى أوروبا حيث عُرفت باسم ``نظرية الباقي الصينية.''

}



\paralleltext{%

Gauss gave the first complete treatment in Europe in his \textit{Disquisitiones Arithmeticae} (1801), calling it the ``problem of finding a number that, when divided by given numbers, leaves given remainders.'' The theorem is now recognized as one of the fundamental results in number theory. It has applications in cryptography (RSA algorithm), error-correcting codes, computer algebra systems, and many areas of modern mathematics.%

}{%

قدَّم غاوس أول معالجة كاملة في أوروبا في كتابه \textit{أبحاث الحساب} (١٨٠١)، سمَّاها ``مسألة إيجاد عدد يُقسم على أعداد معطاة فتترك بواقٍ معطاة.'' تُعرف النظرية اليوم كأحد النتائج الأساسية في نظرية الأعداد. ولها تطبيقات في التشفير (خوارزمية آر-إس-إيه)، ورموز تصحيح الأخطاء، وأنظمة جبر الحاسوب، ومجالات عديدة من الرياضيات الحديثة.

}



\newpage



%% ========================================================================

%% GLOSSARY OF MATHEMATICAL TERMS

%% ========================================================================

\parallelsection{Glossary of Mathematical Terms}{مسرد المصطلحات الرياضية}



\begin{center}

\renewcommand{\arraystretch}{1.3}

\newcommand{\arentry}[1]{{\arabicfont\beginR #1\endR}}

\begin{tabular}{>{\chfont}l >{\itshape}l l}

\toprule

\textbf{中文} & \textbf{English} & \textbf{العربية} \\

\midrule

大衍 & Great Derivation & \arentry{الطريقة الكبرى} \\

大衍求一术 & Method for Finding Unity & \arentry{طريقة استخراج الوحدة} \\

大衍总数术 & General Method & \arentry{الطريقة العامة الكبرى} \\

元数 & Original Numbers & \arentry{الأعداد الأصليَّة} \\

定数 & Definitive Numbers & \arentry{الأعداد المثبتة} \\

定母 & Definitive Mother & \arentry{الأُمّ المثبتة} \\

衍数 & Derivative Numbers & \arentry{الأعداد المشتقَّة} \\

衍母 & Derivative Mother & \arentry{الجذر المشترك} \\

奇数 & Odd Numbers & \arentry{الأعداد الفرديَّة / البواقي} \\

乘率 & Multipliers & \arentry{معاملات الضرب} \\

用数 & Use Numbers & \arentry{الأعداد المستعملة} \\

总数 & Total Number & \arentry{المجموع الكلي} \\

问 & Question & \arentry{المسألة} \\

答 & Answer & \arentry{الجواب} \\

术 & Method & \arentry{الطريقة} \\

草 & Working & \arentry{الحل المفصل} \\

连环求等 & Chain GCD & \arentry{طلب المقاس المشترك المتسلسل} \\

约奇弗约偶 & Reduce odd, not even & \arentry{اقسم الفرديَّ دون الزوجيَّ} \\

满去 & Reduce modulo & \arentry{الاختزال المعياري} \\

余 & Remainder & \arentry{الباقي} \\

天元一 & Unity (1) & \arentry{الوحدة (١)} \\

\bottomrule

\end{tabular}

\end{center}



\newpage



%% ========================================================================

%% ABOUT THIS EDITION

%% ========================================================================

\parallelsection{About This Edition}{عن هذه الطبعة}



\paralleltext{%

This typeset edition presents Fascicle I of Qin Jiushao's \textit{Shuxue Jiuzhang} (Mathematical Treatise in Nine Sections), covering the Dayan category (大衍类) with its nine problems on the Chinese Remainder Theorem and related indeterminate analysis.



The \textbf{Chinese--Arabic bilingual format} presents the original Classical Chinese text on the left and the Arabic translation on the right, following the tradition of parallel translation editions used for mathematical and scientific texts since the medieval period. This format allows scholars of both traditions to compare the original formulations with their Arabic renditions directly.



The source text is based on the Siku Quanshu (四库全书) edition as preserved in the Chinese Textual Tradition, with cross-references to the Yijia Tang (宜稼堂) edition. The editorial commentary (按) included in the Siku edition has been preserved where available.%

}{%

تقدِّم هذه الطبعة المُحكَّمة الفصلَ الأول من \textit{كتاب المسائل التسع في الحساب} لتشين جيوشاو، ويغطي باب الطريقة الكبرى (大衍类) بتسع مسائل حول نظرية الباقي الصينية والتحليل اللامحدود ذي الصلة.



تعرض \textbf{الطبعة الثنائية الصينية--العربية} النص الصيني الكلاسيكي الأصلي في العمود الأيسر والترجمة العربية في العمود الأيمن، على غرار طبعات الترجمة المتوازية المُستخدَمة للنصوص الرياضية والعلمية منذ العصور الوسطى. يتيح هذا التنسيق لعلماء التراثين مقارنة الصياغات الأصلية بترجماتها العربية مباشرةً.



يعتمد النص الأصلي على طبعة سيكو تشياشو (四库全书) كما وُرِّثت في التراث النصي الصيني، مع إحالات صليبية إلى طبعة ييجيا تانغ (宜稼堂). التعليق التحريري (按) المدرج في طبعة سيكو قد حُفِظَ حيثما وُجِدَ.

}



\parallelsubsection{Textual Structure}{بنية النص}



\paralleltext{%

\begin{itemize}

  \item The traditional structure of 问 (question), 答 (answer), 术 (method), and 草 (working) has been preserved

  \item Arabic mathematical terminology follows the tradition of medieval Arabic mathematical texts, with modern standardizations where appropriate

\end{itemize}%

}{%

\begin{itemize}

  \item الهيكل التقليدي لـ 问 (المسألة)، و答 (الجواب)، و术 (الطريقة)، و草 (الحل المفصل) قد حُفِظَ

  \item المصطلحات الرياضية العربية تتبع تراث النصوص الرياضية العربية العصور الوسطى، مع التقييسات الحديثة حيثما كان ذلك ملائماً

\end{itemize}

}



\parallelsubsection{Further Reading}{قراءات إضافية}



\paralleltext{%

\begin{itemize}

  \item Libbrecht, U.\ \textit{Chinese Mathematics in the Thirteenth Century}. MIT Press, 1973.

  \item Qian, Baocong.\ \textit{Zhongguo Shuxue Shi} (History of Chinese Mathematics). Science Press, 1964.

  \item Martzloff, J.-C.\ \textit{A History of Chinese Mathematics}. Springer, 1997.

  \item Lam Lay Yong and Ang Tian Se.\ \textit{Fleeting Footsteps}. World Scientific, 2004.

\end{itemize}%

}{%

\begin{itemize}

  \item ليبريخت، أ. \textit{الرياضيات الصينية في القرن الثالث عشر}. مطبعة معهد ماساتشوستس للتقنية، ١٩٧٣.

  \item تشيان، باو-تسونغ. \textit{تاريخ الرياضيات الصينية}. مطبعة العلوم، ١٩٦٤.

  \item مارتسلوف، ج.-س. \textit{تاريخ الرياضيات الصينية}. سبرنجر، ١٩٩٧.

  \item لام لاي يونغ وآنج تيان سي. \textit{خطوات عابرة: تتبُّع مفهوم الحساب والجبر في الصين القديمة}. وورلد ساينتيفيك، ٢٠٠٤.

\end{itemize}

}



\vspace{2em}

\begin{center}

\textcolor{colrule}{\rule{0.5\textwidth}{0.8pt}}\\[1em]

{\large\scshape Chinese--Arabic Bilingual Edition}\\[0.3em]

{\large\arabicfont\beginR الطبعة الثنائية الصينية--العربية\endR}\\[1em]

\textit{Modern LaTeX edition\ --- \the\year}

\end{center}



\end{document}







% ============================================================

% Source: translations/ar/chinese/sunzi-suanjing_arabic.tex

% ============================================================



%% ========================================================================

%% Sunzi Suanjing (孙子算经) — Chinese → Arabic Bilingual Edition

%% 孫子算經 — الطبعة الثنائية الصينية-العربية

%% Parallel text: Classical Chinese | Arabic Translation

%% ========================================================================

\documentclass[11pt,a4paper]{article}



%% -------- Core Packages --------

\usepackage{fontspec}

\usepackage{polyglossia}

\usepackage{amsmath,amssymb,amsthm}

\usepackage{geometry}

\usepackage{graphicx}

\usepackage{xcolor}

\usepackage{titlesec}

\usepackage{fancyhdr}

\usepackage{enumitem}

\usepackage{array,booktabs,longtable,multirow}

\usepackage{hyperref}

\usepackage{tocloft}

\usepackage{indentfirst}

\usepackage{etoolbox}

\usepackage{framed}

\usepackage{xeCJK}



%% -------- Page Geometry --------

\geometry{

  paperwidth=210mm,paperheight=297mm,

  left=18mm,right=18mm,top=25mm,bottom=25mm,

  headheight=14pt,footskip=30pt

}



%% -------- Language Setup --------

\setmainlanguage{english}

\setotherlanguage{arabic}



%% -------- Font Configuration --------

\setmainfont{Times New Roman}

\setsansfont{Arial}

\setmonofont{Courier New}



%% -------- Arabic Font --------

\newfontfamily\arabicfont[Script=Arabic]{Arial}

\TeXXeTstate=1

\newcommand{\ar}[1]{\arabicfont\beginR #1\endR}



%% -------- CJK Font --------

\setCJKmainfont{Microsoft YaHei}

\setCJKsansfont{Microsoft YaHei}

\setCJKmonofont{Microsoft YaHei}



%% -------- Colors --------

\definecolor{titlecolor}{RGB}{45,55,72}

\definecolor{sectioncolor}{RGB}{55,65,81}

\definecolor{notecolor}{RGB}{120,53,15}

\definecolor{rulecolor}{RGB}{180,160,140}

\definecolor{chinesebg}{RGB}{255,252,245}

\definecolor{problembg}{RGB}{250,248,244}

\definecolor{answerbg}{RGB}{245,250,245}

\definecolor{methodbg}{RGB}{245,248,252}

\definecolor{crtgold}{RGB}{139,90,0}

\definecolor{arabicbg}{RGB}{252,255,250}



%% -------- Section Formatting --------

\titleformat{\section}

  {\normalfont\LARGE\bfseries\color{titlecolor}}

  {\thesection}{1em}{}

\titleformat{\subsection}

  {\normalfont\large\bfseries\color{sectioncolor}}

  {\thesubsection}{1em}{}

\titleformat{\subsubsection}

  {\normalfont\normalsize\bfseries\color{sectioncolor}}

  {\thesubsubsection}{1em}{}



%% -------- Header/Footer --------

\pagestyle{fancy}

\fancyhf{}

\fancyhead[L]{\small\textit{Sunzi Suanjing — الطبعة الصينية-العربية}}

\fancyhead[R]{\small\thepage}

\renewcommand{\headrulewidth}{0.4pt}

\renewcommand{\footrulewidth}{0pt}



%% -------- Custom Commands --------

\newcommand{\longline}{\noindent\rule{\textwidth}{0.4pt}}

\newcommand{\shortline}{\noindent\rule{0.5\textwidth}{0.4pt}\par}

\newcommand{\zh}[1]{{\CJKfamily{}#1}}



%% Bilingual line environment

\newcommand{\biline}[2]{%

  \noindent\begin{minipage}[t]{0.46\textwidth}\small#1\end{minipage}%

  \hfill\vrule\hfill%

  \begin{minipage}[t]{0.46\textwidth}\small\ar{#2}\end{minipage}%

  \par\vspace{6pt}%

}



%% Bilingual display environment (for longer passages)

\newenvironment{bilingual}

  {\par\noindent\ignorespaces}

  {\par}

\newcommand{\cnline}[1]{\begin{minipage}[t]{0.46\textwidth}\small#1\end{minipage}}

\newcommand{\arline}[1]{\begin{minipage}[t]{0.46\textwidth}\small\ar{#1}\end{minipage}}

\newcommand{\bipassage}[2]{%

  \noindent\cnline{#1}\hfill\vrule\hfill\arline{#2}\par\vspace{8pt}%

}



%% Custom framed boxes (replacing tcolorbox)

\newcommand{\custombox}[4]{%

  \par\smallskip\noindent%

  \colorbox{#1}{%

    \fcolorbox{#2}{#1}{%

      \begin{minipage}{0.97\textwidth}#4\end{minipage}%

    }%

  }\par\smallskip%

}



\newcommand{\biproblem}[2]{%

  \par\smallskip\noindent%

  \fcolorbox{rulecolor}{problembg}{%

    \begin{minipage}{0.97\textwidth}%

    \small

    \begin{minipage}[t]{0.46\textwidth}\textbf{\zh{問}}\\#1\end{minipage}%

    \hfill\vrule\hfill%

    \begin{minipage}[t]{0.46\textwidth}\ar{\textbf{المسألة:}\\#2}\end{minipage}%

    \end{minipage}%

  }\par\smallskip%

}

\newcommand{\bianswer}[2]{%

  \par\smallskip\noindent%

  \fcolorbox{rulecolor!70!green}{answerbg}{%

    \begin{minipage}{0.97\textwidth}%

    \small

    \begin{minipage}[t]{0.46\textwidth}\textbf{\zh{答}}\\#1\end{minipage}%

    \hfill\vrule\hfill%

    \begin{minipage}[t]{0.46\textwidth}\ar{\textbf{الجواب:}\\#2}\end{minipage}%

    \end{minipage}%

  }\par\smallskip%

}

\newcommand{\bimethod}[2]{%

  \par\smallskip\noindent%

  \fcolorbox{rulecolor!70!blue}{methodbg}{%

    \begin{minipage}{0.97\textwidth}%

    \small

    \begin{minipage}[t]{0.46\textwidth}\textbf{\zh{術}}\\#1\end{minipage}%

    \hfill\vrule\hfill%

    \begin{minipage}[t]{0.46\textwidth}\ar{\textbf{الطريقة:}\\#2}\end{minipage}%

    \end{minipage}%

  }\par\smallskip%

}



%% Preface shaded environment

\newenvironment{prefacebox}{%

  \definecolor{shadecolor}{RGB}{255,252,245}%

  \begin{shaded}\small\begin{minipage}{0.95\textwidth}%

}{%

  \end{minipage}\end{shaded}%

}



%% Gold CRT box

\newenvironment{crtbox}{%

  \par\medskip\noindent%

  \fcolorbox{crtgold}{crtgold!8}{%

    \begin{minipage}{0.96\textwidth}%

    \medskip

}{%

    \medskip\end{minipage}%

  }\par\medskip%

}



%% Song box

\newenvironment{songbox}{%

  \par\medskip\noindent%

  \fcolorbox{rulecolor}{chinesebg}{%

    \begin{minipage}{0.9\textwidth}\centering\large%

}{%

    \end{minipage}%

  }\par\medskip%

}



%% -------- Hyperref Setup --------

\hypersetup{

  colorlinks=true,

  linkcolor=sectioncolor,

  citecolor=sectioncolor,

  urlcolor=sectioncolor,

  pdfencoding=auto,

  pdftitle={Sunzi Suanjing — Chinese-Arabic Bilingual Edition},

  pdfauthor={Sunzi (c. 400 CE) — Arabic translation edition}

}



%% -------- TOC Depth --------

\setcounter{tocdepth}{2}

\setcounter{secnumdepth}{2}



%% ========================================================================

\begin{document}



%% ========================================================================

%% TITLE PAGE

%% ========================================================================

\begin{titlepage}

\begin{center}

\vspace*{1.5cm}



{\fontsize{13}{15}\selectfont\textsc{Classical Chinese Mathematics} — \ar{رياضيات الصين الكلاسيكية}}\\[0.8cm]



\longline\\[1cm]



{\fontsize{32}{38}\selectfont\bfseries\color{titlecolor}

\zh{孫子算經}}\\[0.6cm]



{\fontsize{20}{24}\selectfont\itshape Sunzi Suanjing}\\[0.3cm]



{\fontsize{26}{30}\selectfont\color{crtgold}\ar{كِتَابُ سُونْزِ فِي الْحِسَاب}}\\[0.6cm]



{\fontsize{14}{18}\selectfont Master Sun's Mathematical Manual}\\[0.3cm]



{\fontsize{14}{18}\selectfont\ar{كتاب سيدنا سُونْز في علم الحساب}}\\[1.2cm]



\longline\\[1.2cm]



\fcolorbox{rulecolor}{chinesebg}{%

\begin{minipage}{0.85\textwidth}

\centering

\small

\textit{One of the Ten Computational Canons (\zh{算經十書}) — \ar{من كتب الرياضيات العشرة}}}\\[4pt]

\textit{Contains the earliest known description of the}\\[2pt]

\textbf{\large Chinese Remainder Theorem}\\[3pt]

{\large\ar{مُسَلَّمَةُ البَاقِي الصِّينِيَّة}}

\end{minipage}%

}



\vfill



{\normalsize\ar{الطبعة الثنائية: الصينية والعربية} — Bilingual Chinese/Arabic Edition}\\[0.4cm]

{\small\today}



\end{center}

\end{titlepage}



%% ========================================================================

%% TABLE OF CONTENTS

%% ========================================================================

\tableofcontents

\newpage



%% ========================================================================

%% INTRODUCTION

%% ========================================================================

\section{Introduction — \ar{مقدمة}}



\bipassage{%

\textbf{The Sunzi Suanjing} (\zh{孫子算經}, ``Master Sun's Mathematical Manual'') is one of the earliest surviving Chinese mathematical texts, written during the 3rd to 5th centuries CE. It was listed as one of the \textit{Ten Computational Canons} (\zh{算經十書}) during the Tang dynasty. The treatise consists of three volumes (\zh{卷}):

\begin{itemize}[leftmargin=1em,topsep=0pt,itemsep=1pt]

  \item \textbf{Volume I} (\zh{卷上}): Units, counting rods, multiplication

  \item \textbf{Volume II} (\zh{卷中}): Fractions, geometry, practical problems

  \item \textbf{Volume III} (\zh{卷下}): 36 problems including the CRT

\end{itemize}%

}{%

\textbf{كتاب سونز في الحساب} (孫子算經، «كتاب سيدنا سُونْز في علم الحساب») هو أحد أقدم النصوص الرياضية الصينية التي وصلت إلينا، وقد أُلِّفَ في القرن الثالث إلى الخامس الميلادي. وُسِّمَ الكتاب إلى ثلاثة أجزاء (卷):

\begin{itemize}[leftmargin=1em,topsep=0pt,itemsep=1pt]

  \item \textbf{الجزء الأول} (卷上): وحدات القياس، أعواد الحساب، جدول الضرب

  \item \textbf{الجزء الثاني} (卷中): الكسور، الهندسة، مسائل عملية

  \item \textbf{الجزء الثالث} (卷下): ست وثلاثون مسألة تشمل مسألة الباقي الصينية

\end{itemize}%

}



\bipassage{%

The text's most celebrated contribution is \textbf{Problem 26 of Volume III}:

\begin{quote}\small

``There are certain things whose number is unknown. If we count them by threes, we have two left over; by fives, we have three left over; and by sevens, two are left over. How many things are there?''

\end{quote}

This problem represents the earliest known statement of the Chinese Remainder Theorem. The solution method gives:

\[N = 70r_1 + 21r_2 + 15r_3 - 105k\]

where $r_1, r_2, r_3$ are remainders upon division by 3, 5, and 7.

}{%

أهم إسهام في هذا الكتاب هو \textbf{المسألة السادسة والعشرون من الجزء الثالث}:

\begin{quote}\small\ar{توجد أشياء لا نعرف عددها. إذا عددناها على ثلاثات بقيت اثنتان، وعلى خمسات بقيت ثلاث، وعلى سبعات بقيت اثنتان. كم عددها؟}

\end{quote}

هذه المسألة تمثل أقدم صياغة معروفة لمسألة الباقي الصينية. طريقة الحل:

\[N = 70r_1 + 21r_2 + 15r_3 - 105k\]

حيث $r_1, r_2, r_3$ هي بقايا القسمة على ٣ و٥ و٧.

}



\newpage



%% ========================================================================

%% PREFACE

%% ========================================================================

\section*{Preface — \ar{المقدمة} (\zh{序})}

\addcontentsline{toc}{section}{Preface (\zh{序}) — \ar{المقدمة}}



\begin{prefacebox}

\large

\zh{孫子曰：夫算者，天地之經緯，群生之元首；五常之本末，陰陽之父母；星辰之建號，三光之表裹；五行之準平，四時之終始；萬物之祖宗，六藝之綱紀。}

\end{prefacebox}



\vspace{0.3cm}



\bipassage{%

\textbf{Sun Zi said:} Calculation is the warp and weft of heaven and earth, the origin of all living beings; the root and branch of the Five Constants, the father and mother of yin and yang; the establishment of titles for stars and constellations, the exterior and interior of the Three Lights; the standard balance of the Five Elements, the beginning and end of the Four Seasons; the ancestor of all things, the framework of the Six Arts.

}{%

\textbf{قال سيدنا سونز:} علم الحساب هو نسيج السماء والأرض، وأصل جميع الكائنات؛ أساس الخمس ثوابت وأب وأم الين واليانغ؛ ترتيب النجوم والأبراج، وباطن وظاهر الأجرام الثلاثة؛ موازين العناصر الخمسة، وبداية ونهاية الفصول الأربعة؛ أصل الأشياء كلها، ومنهاج الفنون الستة.

}



\bipassage{%

Investigate the gathering and dispersing of all categories, examine the descending and ascending of the Two Qi; track the cycles of cold and heat, pace the differences between near and far; observe the subtle foundations of the heavenly Dao, examine the lengths and widths of geography; gather signs of spirits, test omens of success and failure; exhaust the principles of virtue, investigate the nature of life. Establish rules and compasses, standardize squares and circles, observe laws and regulations, measure with precision. Enduring through countless generations without decay, its influence extends to all directions without limit. Those who face toward it are rich with abundance, while those who turn away are poor and destitute. Those whose minds open understand even as children, while those whose intentions remain closed find mastery difficult even with white hair.

}{%

فَحَصُ اجتماع وتفرق الفئات كلها، وراقب صعود وهبوط القِوَتَيْن؛ تتبع دورات البرد والحر، وقِسْ الاختلاف بين القريب والبعيد؛ راقب أسس الطريق السماوي الدقيقة، وفحص أطوال وعرض الجغرافيا؛ اجمع علامات الأرواح، واختبر فُؤُول النجاح والفشل؛ استنفد مبادئ الفضيلة، وفحص طبيعة الحياة. اضع القواعد والدوائر، وسوِّ المربعات والدوائر، واعمل القوانين واللوائح، وقِسْ بدقة. يبقى عبر أجيال لا تنحصر بلا فناء، ونفوذه يمتد إلى الجهات كلها بلا حدود. من يواجهه يكون غنيًا بالوفرة، ومن يتولَّ عنه يكون فقيرًا معدمًا. من تفتح عقله يفهم ولو كان صغيرًا، ومن بقي بصيرته مغلقة يجد الصعوبة حتى ولو شاب شعره.

}



\bipassage{%

Therefore, one who wishes to study must necessarily measure his own ability and assess himself carefully, setting his mind on specialization. If so, how could there be any failure?

}{%

فمن أراد أن يتعلم فلا بد أن يقيس قدرته ويُقَوِّم نفسه بعناية، ويوجه عقله نحو التخصص. إن فعل ذلك، فكيف يمكن أن يفشل؟

}



\newpage



%% ========================================================================

%% VOLUME I

%% ========================================================================

\section{Volume I (\zh{卷上}) — \ar{الجزء الأول: وحدات القياس وأعواد الحساب}}



\subsection*{1. The Origin of Length Measurements — \ar{أصل وحدات الطول}}



\biproblem{%

\zh{度之所起，起於忽。欲知其忽，蠶所生，吐絲為忽。十忽為一秒，十秒為一毫，十毫為一釐，十釐為一分，十分為一寸，十寸為一尺，十尺為一丈，十丈為一引。}%

}{%

أصل القياس يبدأ بالوحدة \textit{هُو} (忽). لتعرف الهُو: فإن دودة القز تنتجه من غزل الحرير. عشر هُو تساوي \textit{مِيَاو} (秒)، وعشر مِيَاو تساوي \textit{هَاو} (毫)، وعشر هَاو تساوي \textit{لِي} (釐)، وعشر لِي تساوي \textit{فَن} (分)، وعشر فَن تساوي \textit{تْسُون} (寸)، وعشر تْسُون تساوي \textit{تْشِ} (尺)، وعشر تْشِ تساوي \textit{جَانْغ} (丈)، وعشر جَانْغ تساوي \textit{يِن} (引).

}



\subsection*{3. Large Numbers — \ar{الأعداد الكبيرة}}



\biproblem{%

\zh{凡大數之法，萬萬曰億，萬萬億曰兆，萬萬兆曰京，萬萬京曰陔，萬萬陔曰秭，萬萬秭曰壤，萬萬壤曰溝，萬萬溝曰澗，萬萬澗日正，萬萬正曰載。}%

}{%

قاعدة الأعداد الكبيرة: عشرة آلاف مضروبة في عشرة آلاف تُسَمَّى \textit{يِ} (億)، وعشرة آلاف \textit{يِ} تُسَمَّى \textit{جَاوْ} (兆)، وعشرة آلاف \textit{جَاوْ} تُسَمَّى \textit{جِينْغ} (京)، وعشرة آلاف \textit{جِينْغ} تُسَمَّى \textit{غاي} (陔)، وعشرة آلاف \textit{غاي} تُسَمَّى \textit{زِ} (秭)، وهكذا إلى \textit{رَانْغ} و\textit{قُو} و\textit{جْيَانْ} و\textit{جَنْغ} و\textit{زاي} (載).

}



\subsection*{6. Counting Rod Notation — \ar{طريقة أعواد الحساب}}



\biproblem{%

\zh{凡筭之法，先識其位，一從十橫，百立千僵，千十相望，萬百相當。}%

}{%

طريقة الحساب تبدأ بمعرفة المواضع: الواحد عمودي، والعشرة أفقية، والمئة واقفة، والألف مستلقية. الألوف والعشرات تقابل بعضها، والعشرات والمئات تتطابق.

}



\bipassage{%

This describes the traditional Chinese counting rod numeral system: Units (1--9) are represented by \textbf{vertical} rods; Tens (10--90) by \textbf{horizontal} rods; Hundreds (100--900) by \textbf{vertical} rods; Thousands (1000--9000) by \textbf{horizontal} rods. The pattern alternates for each place value.

}{%

هذا يصف نظام أعواد الحساب الصينية التقليدية: الوحدات (١-٩) تُمثَّل بأعواد \textbf{عمودية}؛ العشرات (١٠-٩٠) بأعواد \textbf{أفقية}؛ المئات (١٠٠-٩٠٠) بأعواد \textbf{عمودية}؛ الألوف (١٠٠٠-٩٠٠٠) بأعواد \textbf{أفقية}. النمط يتناوب بين العمودي والأفقي مع كل منزلة.

}



\subsection*{7. Method of Multiplication — \ar{طريقة الضرب}}



\biproblem{%

\zh{凡乘之法，重置其位。上下相觀，上位有十步至十，有百步至百。以上命下，所得之數列於中位。言十即過，不滿自如。上位乘訖者先去之。下位乘訖者則俱退之。六不積，五不隻。上下相乘，至盡則已。}%

}{%

طريقة الضرب: رتِّب المواضع. انظر من الأعلى والأسفل: إن كان في الأعلى عشرة فتقدم إلى العشرات، وإن كان مئة فتقدم إلى المئات. اضرب الأعلى في الأسفل، والنتيجة تُوضع في الوسط. إن بلغت عشرة فاحمل، وإن لم تبلغ فاتركها. الموضع الأعلى إذا انتهى من الضرب يُزال أولاً. الموضع الأسفل إذا انتهى من الضرب تتراجع كلها. الستة لا تُجمع، والخمسة لا تكون وحيدة. اضرب الأعلى والأسفل حتى الانتهاء.

}



\subsection*{8. Method of Division — \ar{طريقة القسمة}}



\biproblem{%

\zh{凡除之法，與乘正異。乘得在中央，除得在上方。假令六為法，百為實。以六除百，當進之二等，令在正百下，以六除一，則法多而實少，不可除，故當退就十位。}%

}{%

طريقة القسمة تختلف عن الضرب تمامًا. في الضرب الناتج في الوسط، وفي القسمة الناتج في الأعلى. لنفرض أن الستة هي القاسم والمئة هي المقسوم. قسمة المئة على الستة: تقدم منزلتين وضع تحت المئة، والقسمة على الستة إنما القاسم أكبر من المقسوم فلا يمكن القسمة، فيجب الرجوع إلى منزلة العشرات.

}



\subsection*{16. Example: $81 \times 81$ — \ar{مثال: ضرب ٨١ في ٨١}}



\biproblem{%

\zh{九九八十一，自相乘，得幾何？}%

}{%

ثمانية وتسعون في ثمانية وثمانون، نضربهما في بعضهما، فكم يكون؟

}

\bianswer{%

\zh{答曰：六千五百六十一。}%

}{%

الجواب: ستة آلاف وخمسمئة وواحد وستون.

}

\bimethod{%

\zh{術曰：重置其位，以上八呼下八，八八六十四，即下六千四百於中位。以上八呼下一，一八如八，即於中位下八十。退下位一等，收上位八十。以上位一呼下八，一八如八，即於中位下八十。以上一呼下一，一一如一，即於中位下一。上下位俱收，中位即得六千五百六十一。}%

}{%

الطريقة: رتِّب المواضع، انادِ ثمانية من الأعلى بثمانية من الأسفل: ثمانية في ثمانية أربعة وستون، فضع أربعة آلاف وستمئة في الوسط. انادِ ثمانية من الأعلى بواحد من الأسفل: واحد في ثمانية بثمانية، فضع ثمانين في الوسط. أنزل الموضع السفلي منزلة واحدة، واحفظ ثمانين في الأعلى. انادِ واحدًا من الأعلى بثمانية من الأسفل: واحد في ثمانية بثمانية، فضع ثمانين في الوسط. انادِ واحدًا من الأعلى بواحد من الأسفل: واحد في واحد بواحد، فضع واحدًا في الوسط. اجمع المواضع العليا والسفلى، فيصبح في الوسط ستة آلاف وخمسمئة وواحد وستون.

}



\newpage



%% ========================================================================

%% VOLUME II

%% ========================================================================

\section{Volume II (\zh{卷中}) — \ar{الجزء الثاني: الكسور والمساحات}}



\subsection*{Problem 1. Reducing Fractions — \ar{المسألة الأولى: تبسيط الكسور}}



\biproblem{%

\zh{今有一十八分之一十二。問約之得幾何？}%

}{%

لدينا اثنا عشر من ثمانية عشر. إذا بسَّطناها فكم تكون؟

}

\bianswer{%

\zh{答曰：三分之二。}%

}{%

الجواب: ثلثان.

}

\bimethod{%

\zh{術曰：置十八分在下，一十二分在上。副置二位，以少減多，等數得六，為法。約之，即得。}%

}{%

الطريقة: ضع ثمانية عشر في الأسفل واثني عشر في الأعلى. ضع الرقمين معًا، اطرح الأصغر من الأكبر؛ العدد المشترك هو ستة وهو القاسم. اقسم عليه فتحصل على النتيجة.

}



\subsection*{Problem 9. Brick Calculation — \ar{المسألة التاسعة: حساب الطوب}}



\biproblem{%

\zh{今有屋基南北三丈，東西六丈，欲以磚砌之。凡積二尺，用磚五枚。問計幾何？}%

}{%

لدينا أساس بيت طوله من الشمال إلى الجنوب ثلاثة \textit{جانْغ}، وعرضه من الشرق إلى الغرب ستة \textit{جانْغ}، ونريد بناءه بالطوب. كل قدمين مكعبين يحتاجان خمس طوبات. كم تحتاج؟

}

\bianswer{%

\zh{答曰：四千五百枚。}%

}{%

الجواب: أربعة آلاف وخمسمئة طوبة.

}

\bimethod{%

\zh{術曰：置東西六丈，以南北三丈乘之，得一千八百尺。以五乘之，得九千尺。以二除之，即得。}%

}{%

الطريقة: ضع عرض الشرق والغرب ستة \textit{جانْغ}، اضربه في طول الشمال والجنوب ثلاثة \textit{جانْغ}، فتحصل على ألف وثمان مئة قدم. اضربها في خمسة فتكون تسعة آلاف قدم. اقسمها على اثنين فتحصل على النتيجة.

}



\subsection*{Problem 13. Area of a Circular Field — \ar{المسألة الثالثة عشر: مساحة الحقل الدائري}}



\biproblem{%

\zh{今有圓田周三百步，徑一百步。問得田幾何？}%

}{%

لدينا حقل دائري محيطه ثلاثمئة خطوة، وقطره مئة خطوة. كم تبلغ مساحة الحقل؟

}

\bianswer{%

\zh{答曰：三十一畝奇六十步。}%

}{%

الجواب: واحد وثلاثون \textit{مُو} وستون خطوة زائدة.

}

\bimethod{%

\zh{術曰：先置周三百步，半之，得一百五十步。又置徑一百步，半之，得五十步。相乘，得七千五百步。以畝法二百四十步除之，即得。}%

}{%

الطريقة: ضع المحيط ثلاثمئة خطوة، نصفها فتصبح مئة وخمسون. ثم ضع القطر مئة خطوة، نصفها فتصبح خمسين. اضربهما فتصبح سبعة آلاف وخمسمئة. اقسم على قانون المُو وهو مئتان وأربعون خطوة، فتحصل على النتيجة.

}



\subsection*{Problem 23. Canal Construction — \ar{المسألة الثالثة والعشرون: حفر قناة}}



\biproblem{%

\zh{今有穿渠，長二十九里一百四步，上廣一丈二尺六寸，下廣八尺，深一丈八尺。秋程人功三百尺。問須功幾何？}%

}{%

لدينا قناة لحفرها، طولها تسعة وعشرون لِي ومئة وأربع خطوات، عرضها العلوي \textit{جانْغ} واحد وقدمان وستة \textit{تْسُون}، وعرضها السفلي ثمانية أقدام، وعمقها \textit{جانْغ} واحد وثمانية أقدام. كمية عمل الشخص في الخريف ثلاثمئة قدم. كم تحتاج من العمال؟

}

\bianswer{%

\zh{答曰：三萬二千六百四十五人，不盡六十九尺六寸。}%

}{%

الجواب: اثنان وثلاثون ألفًا وستمئة وخمسة وأربعون شخصًا، وبقية تسعة وستون قدمًا وستة \textit{تْسُون}.

}



\newpage



%% ========================================================================

%% VOLUME III

%% ========================================================================

\section{Volume III (\zh{卷下}) — \ar{الجزء الثالث: مسائل عملية}}



\subsection*{Problem 5. Go Board — \ar{المسألة الخامسة: لوحة اللو}}



\biproblem{%

\zh{今有碁局方一十九道。問用碁幾何？}%

}{%

لدينا لوحة لو مربعة ذات تسعة عشر خطًا في كل جانب. كم حجرًا يُستخدم؟

}

\bianswer{%

\zh{答曰：三百六十一。}%

}{%

الجواب: ثلاثمئة وواحد وستون.

}

\bimethod{%

\zh{術曰：置一十九道，自相乘之，即得。}%

}{%

الطريقة: ضع تسعة عشر خطًا، واضربها في نفسها، فتحصل على النتيجة.

}



\bipassage{%

This gives $19^2 = 361$ intersection points on a Go board.

}{%

وهذا يعطي $19^2 = 361$ نقطة تقاطع على لوحة اللو.

}



\subsection*{Problem 15. People and Carts — \ar{المسألة الخامسة عشر: الناس والعربات}}



\biproblem{%

\zh{今有三人共車，二車空；二人共車，九人步。問人與車各幾何？}%

}{%

إذا استقلَّ ثلاثة أشخاص كل عربة، فتبقى عربتان فارغتين. وإذا استقلَّ اثنان في كل عربة، فتسعة يمشون. كم عدد الناس والعربات؟

}

\bianswer{%

\zh{答曰：一十五車。三十九人。}%

}{%

الجواب: خمس عشرة عربة. تسعة وثلاثون شخصًا.

}

\bimethod{%

\zh{術曰：置二車，以三乘之，得六。加步者九人，得車一十五。欲知人者，以二乘車，加九人，即得。}%

}{%

الطريقة: ضع العربتين، واضربهما في ثلاثة فتصبح ستة. أضف المشاة التسعة، فتصبح العربات خمس عشرة. ولمعرفة الناس: اضرب العربات في اثنين، وأضف التسعة، فتحصل على النتيجة.

}



\subsection*{Problem 17. Guests at a Feast — \ar{المسألة السابعة عشر: الضيوف في الوليمة}}



\biproblem{%

\zh{今有婦人河上蕩桮。津吏問曰：「桮何以多？」婦人曰：「家有客。」津吏曰：「客幾何？」婦人曰：「二人共飯，三人共羹，四人共肉，凡用桮六十五。不知客幾何？」}%

}{%

كانت امرأة تغسل أكوابًا على النهر. سألها عامل المَعْبَر: «لمَ الأكواب كثيرة؟» قالت: «إن في الدار ضيوفًا.» قال: «كم الضيوف؟» قالت: «اثنان يتقاسمان الطعام، وثلاثة يتقاسمان المرق، وأربعة يتقاسمان اللحم، ومجموع الأكواب خمسة وستون. لا أعلم كم الضيوف.»

}

\bianswer{%

\zh{答曰：六十人。}%

}{%

الجواب: ستون شخصًا.

}

\bimethod{%

\zh{術曰：置六十五桮，以一十二乘之，得七百八十，以十三除之，即得。}%

}{%

الطريقة: ضع خمسة وستين كوبًا، واضربها في اثني عشر فتصبح سبعمئة وثمانون. اقسمها على ثلاثة عشر، فتحصل على النتيجة.

}



\subsection*{Problem 18. Wood and Rope — \ar{المسألة الثامنة عشر: الخشب والحبل}}



\biproblem{%

\zh{今有木，不知長短。引繩度之，餘繩四尺五寸。屈繩量之，不足一尺。問木長幾何？}%

}{%

لدينا خشب لا نعرف طوله. سُمِكَ بِحَبْل فبقي من الحبل أربعة أقدام وخمسة \textit{تْسُون}. طُوِيَ الحبل لِقياسه فنقص قدمًا واحدة. كم طول الخشب؟

}

\bianswer{%

\zh{答曰：六尺五寸。}%

}{%

الجواب: ستة أقدام وخمسة \textit{تْسُون}.

}

\bimethod{%

\zh{術曰：置餘繩四尺五寸，加不足一尺，共五尺五寸。倍之，得一丈一尺。減餘四尺五寸，即得。}%

}{%

الطريقة: ضع بقية الحبل أربعة أقدام وخمسة \textit{تْسُون}، وأضف النقص وهو قدم واحدة، فتصبح خمسة أقدام وخمسة \textit{تْسُون}. اضربها في اثنين فتصبح \textit{جانْغ} واحد وقدمان. اطرح البقية أربعة أقدام وخمسة \textit{تْسُون}، فتحصل على النتيجة.

}



\subsection*{Problem 19. Rice in a Vessel — \ar{المسألة التاسعة عشر: الأرز في الوعاء}}



\biproblem{%

\zh{今有器中米，不知其數。前人取半，中人三分取一，後人四分取一，餘米一斗五升。問本米幾何？}%

}{%

في وعاءٍ أرزٌ لا نعرف كميته. أخذ الأول نصفه، وأخذ الثاني ثلث ما بقي، وأخذ الثالث ربع ما بقي بعد ذلك، فبقي من الأرز \textit{تُو} واحد وخمسة \textit{شَنْغ}. كم كان الأرز الأصلي؟

}

\bianswer{%

\zh{答曰：六斗。}%

}{%

الجواب: ستة أَدْوَال (تُو).

}

\bimethod{%

\zh{術曰：置餘米一斗五升，以六乘之，得九斗。以二除之，得四斗五升。以四乘之，得一斛八斗。以三除之，即得。}%

}{%

الطريقة: ضع بقية الأرز \textit{تُو} واحد وخمسة \textit{شَنْغ}، واضربها في ستة فتصبح تسعة أَدْوَال. اقسمها على اثنين فتصبح أربعة أَدْوَال وخمسة \textit{شَنْغ}. اضربها في أربعة فتصبح \textit{هُو} واحد وثمانية أَدْوَال. اقسمها على ثلاثة فتحصل على النتيجة.

}



\subsection*{Problem 31. Pheasants and Rabbits — \ar{المسألة الحادية والثلاثون: الدجاج والأرانب في القفص}}



\biproblem{%

\zh{今有雉、兔同籠，上有三十五頭，下九十四足。問雉、兔各幾何？}%

}{%

لدينا طيور وأرانب في قفص واحد: فوق خمس وثلاثون رأسًا، وتحت أربع وتسعون قدمًا. كم طائرًا وكم أرنبًا؟

}

\bianswer{%

\zh{答曰：雉二十三。兔一十二。}%

}{%

الجواب: ثلاثة وعشرون طائرًا. اثنا عشر أرنبًا.

}

\bimethod{%

\zh{術曰：上置三十五頭，下置九十四足。半其足，得四十七。以少減多，再命之，上三除下三，上五除下五。下有一除上一，下有二除上二，即得。}%

}{%

الطريقة: ضع في الأعلى خمس وثلاثين رأسًا، وفي الأسفل أربع وتسعين قدمًا. أنصف الأقدام فتصبح سبع وأربعين. اطرح الأصغر من الأكبر وكرِّر، ثم اقسم. فتحصل على النتيجة.

}



\subsection*{Problem 34. Nine Dikes — \ar{المسألة الرابعة والثلاثون: السدود التسعة}}



\biproblem{%

\zh{今有出門望見九隄。隄有九木，木有九枝，枝有九巢，巢有九禽，禽有九雛，雛有九毛，毛有九色。問各幾何？}%

}{%

خرج رجل من بابه فرأى تسع سدود. كل سد فيه تسع شجرات، وكل شجرة فيها تسع أغصان، وكل غصن فيه تسع أعشاش، وكل عش فيه تسع طيور، وكل طائر فيه تسع فراخ، وكل فرخ فيه تسع ريشات، وكل ريشة فيها تسع ألوان. كم كل منها؟

}

\bianswer{%

\zh{答曰：木八十一。枝七百二十九。巢六千五百六十一。禽五萬九千四十九。雛五十三萬一千四百四十一。毛四百七十八萬二千九百六十九。色四千三百四萬六千七百二十一。}%

}{%

الجواب: شجرات واحد وثمانون. أغصان سبعمئة وتسع وعشرون. أعشاش ستة آلاف وخمسمئة وواحد وستون. طيور تسعة وأربعون ألفًا وتسعة وأخمسون. فراخ واحد وثلاثون ألفًا وأربعمئة وواحد وأربعون. ريشات تسعمئة واثنتان وستون ألفًا وثمانمئة وتسعة وسبعون. ألوان واحد وعشرون مليونًا وستمئة وسبعة وسبعون ألفًا ومئتان وواحد.

}



\subsection*{Problem 35. Three Daughters — \ar{المسألة الخامسة والثلاثون: البنات الثلاث}}



\biproblem{%

\zh{今有三女，長女五日一歸，中女四日一歸，少女三日一歸。問三女幾何日相會？}%

}{%

لدينا ثلاث بنات: الكبرى تعود كل خمسة أيام، والوسطى كل أربعة أيام، والصغرى كل ثلاثة أيام. بعد كم يوم تلتقي البنات الثلاث معًا؟

}

\bianswer{%

\zh{答曰：六十日。}%

}{%

الجواب: ستون يومًا.

}

\bimethod{%

\zh{術曰：置長女五日、中女四日、少女三日，於右方。各列一算於左方。維乘之，各得所到數：長女十二到，中女十五到，少女二十到。又各以歸日乘到數，即得。}%

}{%

الطريقة: ضع أيام الكبرى خمسة، والوسطى أربعة، والصغرى ثلاثة. ضع لكل منها حسابًا على اليسار. اضربها: الكبرى تصل اثنتي عشرة مرة، والوسطى خمس عشرة، والصغرى عشرين. ثم اضرب أيام العودة في عدد مرات الوصول، فتحصل على النتيجة.

}



\bipassage{%

This is a least common multiple problem: $\text{lcm}(3,4,5) = 60$.

}{%

هذه مسألة المضاعف المشترك الأصغر: $\text{lcm}(3,4,5) = 60$.

}



\newpage



%% ========================================================================

%% THE CHINESE REMAINDER THEOREM

%% ========================================================================

\section*{\zh{物不知數} — Problem of the Unknown Quantity}

\addcontentsline{toc}{section}{The Chinese Remainder Theorem (\zh{物不知數}) — \ar{مُسَلَّمَةُ البَاقِي الصِّينِيَّة}}



\begin{crtbox}

\noindent\textbf{\large Problem 26, Volume III (\zh{卷下第26問}) — \ar{المسألة السادسة والعشرون}}



\vspace{0.3cm}



\biproblem{%

\zh{今有物，不知其數。三、三數之，賸二；五、五數之，賸三；七、七數之，賸二。問物幾何？}%

}{%

توجد أشياء لا نعرف عددها. إذا عددناها على ثلاثات بقيت اثنتان، وعلى خمسات بقيت ثلاث، وعلى سبعات بقيت اثنتان. كم عددها؟

}



\bianswer{%

\zh{答曰：二十三。}%

}{%

الجواب: ثلاث وعشرون.

}



\bimethod{%

\zh{術曰：「三、三數之，賸二」，置一百四十；「五、五數之，賸三」，置六十三；「七、七數之，賸二」，置三十。并之，得二百三十三。以二百一十減之，即得。凡三、三數之，賸一，則置七十；五、五數之，賸一，則置二十一；七、七數之，賸一，則置十五。一百六以上，以一百五減之，即得。}%

}{%

الطريقة: «إذا عددنا على ثلاثات بقيت اثنتان»، فضع مئة وأربعين؛ «إذا عددنا على خمسات بقيت ثلاث»، فضع ثلاثة وستين؛ «إذا عددنا على سبعات بقيت اثنتان»، فضع ثلاثين. اجمعها فتصبح مئتين وثلاثة وثلاثين. اطرح منها مئتين وعشرة، فتحصل على النتيجة. عامة: إذا عددنا على ثلاثات بقيت واحدة، فضع سبعين؛ وعلى خمسات بقيت واحدة، فضع واحدًا وعشرين؛ وعلى سبعات بقيت واحدة، فضع خمسة عشر. إذا كان المجموع فوق مئة وستة، فاطرح منه مئة وخمسة، فتحصل على النتيجة.

}



\end{crtbox}



\vspace{0.5cm}



\subsection*{Modern Mathematical Analysis — \ar{التحليل الرياضي الحديث}}



\bipassage{%

In modern notation, this problem asks us to find the smallest positive integer $N$ satisfying the system of congruences:

\begin{align*}

N &\equiv 2 \pmod{3} \\

N &\equiv 3 \pmod{5} \\

N &\equiv 2 \pmod{7}

\end{align*}

The method gives:

\begin{align*}

N &= 140 + 63 + 30 - 210 \\

  &= 233 - 210 \\

  &= 23

\end{align*}

}{%

بالصياغة الحديثة، تطلب هذه المسألة إيجاد أصغر عدد صحيح موجب $N$ يحقق نظام المتطابقات:

\begin{align*}

N &\equiv 2 \pmod{3} \\

N &\equiv 3 \pmod{5} \\

N &\equiv 2 \pmod{7}

\end{align*}

الطريقة تعطي:

\begin{align*}

N &= 140 + 63 + 30 - 210 \\

  &= 233 - 210 \\

  &= 23

\end{align*}

}



\bipassage{%

The three key numbers \textbf{70}, \textbf{21}, and \textbf{15} have special properties:

\begin{itemize}[leftmargin=1em,topsep=0pt,itemsep=2pt]

  \item $70 = 2 \times 5 \times 7 \equiv 1 \pmod{3}$, and $70 \equiv 0 \pmod{5}$, $70 \equiv 0 \pmod{7}$

  \item $21 = 1 \times 3 \times 7 \equiv 1 \pmod{5}$, and $21 \equiv 0 \pmod{3}$, $21 \equiv 0 \pmod{7}$

  \item $15 = 1 \times 3 \times 5 \equiv 1 \pmod{7}$, and $15 \equiv 0 \pmod{3}$, $15 \equiv 0 \pmod{5}$

\end{itemize}

Since $3 \times 5 \times 7 = 105$, the general solution is:

\[N = 70r_1 + 21r_2 + 15r_3 - 105k\]

}{%

الأرقام الثلاثة الرئيسية \textbf{٧٠} و\textbf{٢١} و\textbf{١٥} لها خصائص خاصة:

\begin{itemize}[leftmargin=1em,topsep=0pt,itemsep=2pt]

  \item $70 = 2 \times 5 \times 7 \equiv 1 \pmod{3}$، و $70 \equiv 0 \pmod{5}$، و $70 \equiv 0 \pmod{7}$

  \item $21 = 1 \times 3 \times 7 \equiv 1 \pmod{5}$، و $21 \equiv 0 \pmod{3}$، و $21 \equiv 0 \pmod{7}$

  \item $15 = 1 \times 3 \times 5 \equiv 1 \pmod{7}$، و $15 \equiv 0 \pmod{3}$، و $15 \equiv 0 \pmod{5}$

\end{itemize}

وبما أن $3 \times 5 \times 7 = 105$، فالحل العام هو:

\[N = 70r_1 + 21r_2 + 15r_3 - 105k\]

}



\subsection*{Verification — \ar{التحقق}}



\bipassage{%

\begin{align*}

23 \div 3 &= 7 \text{ remainder } \mathbf{2} \quad \checkmark \\

23 \div 5 &= 4 \text{ remainder } \mathbf{3} \quad \checkmark \\

23 \div 7 &= 3 \text{ remainder } \mathbf{2} \quad \checkmark

\end{align*}

}{%

\begin{align*}

23 \div 3 &= 7 \text{ و الباقي } \mathbf{2} \quad \checkmark \\

23 \div 5 &= 4 \text{ و الباقي } \mathbf{3} \quad \checkmark \\

23 \div 7 &= 3 \text{ و الباقي } \mathbf{2} \quad \checkmark

\end{align*}

}



\subsection*{Sunzi's Formula — \ar{قانون سونز}}



\bipassage{%

For general remainders $r_1, r_2, r_3$:

\[

\boxed{N = 70 \times r_1 + 21 \times r_2 + 15 \times r_3 - 105 \times k}

\]

where $M = 3 \times 5 \times 7 = 105$ is the least common multiple, and:

\begin{itemize}[leftmargin=1em,topsep=0pt,itemsep=1pt]

  \item $70 = 105/3 \times 2$ is the multiplier for remainder mod 3

  \item $21 = 105/5 \times 1$ is the multiplier for remainder mod 5

  \item $15 = 105/7 \times 1$ is the multiplier for remainder mod 7

\end{itemize}

}{%

للبقايا العامة $r_1, r_2, r_3$:

\[

\boxed{N = 70 \times r_1 + 21 \times r_2 + 15 \times r_3 - 105 \times k}

\]

حيث $M = 3 \times 5 \times 7 = 105$ هو المضاعف المشترك الأصغر، و:

\begin{itemize}[leftmargin=1em,topsep=0pt,itemsep=1pt]

  \item $70 = 105/3 \times 2$ هو المعامل للباقي من القسمة على ٣

  \item $21 = 105/5 \times 1$ هو المعامل للباقي من القسمة على ٥

  \item $15 = 105/7 \times 1$ هو المعامل للباقي من القسمة على ٧

\end{itemize}

}



\subsection*{Historical Significance — \ar{الأهمية التاريخية}}



\bipassage{%

This problem, known as the \textbf{``Problem of the Unknown Quantity''} (\zh{物不知數}), is the earliest known statement of what we now call the \textbf{Chinese Remainder Theorem}. While Sunzi's text gives only a specific numerical method without proof, it contains the essential idea:

\begin{enumerate}[leftmargin=1em,topsep=0pt,itemsep=2pt]

  \item Find numbers that are multiples of all but one modulus, and leave remainder 1 when divided by the remaining modulus

  \item Multiply each such number by the corresponding remainder

  \item Sum the results and subtract multiples of the LCM to get the smallest positive solution

\end{enumerate}

}{%

هذه المسألة، المعروفة بـ«مسألة العدد المجهول» (物不知數)، هي أقدم صياغة معروفة لما نسميه الآن \textbf{مسألة الباقي الصينية}. ورغم أن نص سونز يعطي طريقة عددية محددة دون برهان، إلا أنه يحتوي على الفكرة الأساسية:

\begin{enumerate}[leftmargin=1em,topsep=0pt,itemsep=2pt]

  \item أوجد أعدادًا هي مضاعفات الكل ما عدا معيار واحد، وتترك باقيًا واحدًا عند قسمتها على المعيار الباقي

  \item اضرب كل عدد منها في الباقي المقابل

  \item اجمع النتائج واطرح مضاعفات المضاعف المشترك الأصغر للحصول على أصغر حل موجب

\end{enumerate}

}



\bipassage{%

The general theorem was later fully developed by \textbf{Qin Jiushao} (\zh{秦九韶}, 1202--1261) in his \textit{Mathematical Treatise in Nine Sections} (\zh{數書九章}, 1247 CE), where he called it the \textbf{``Great Extension Seeking-Unity Method''} (\zh{大衍求一術}). This method could solve systems of arbitrary linear congruences.

}{%

المسألة العامة طوَّرها لاحقًا \textbf{تْشِن جْيوشَاو} (秦九韶، ١٢٠٢-١٢٦١) في كتابه «المعالجة الرياضية في تسعة أقسام» (數書九章، ١٢٤٧م)، وسماها \textbf{«طريقة الامتداد العظيم لطلب الواحد»} (大衍求一術). وكانت هذه الطريقة تستطيع حل أنظمة المتطابقات الخطية من أي نوع.

}



\bipassage{%

In Europe, similar methods were not developed until the work of Euler (1743) and Gauss (1801). In 1852, the British missionary Alexander Wylie published notes on Chinese science, introducing this problem to European mathematicians. In 1876, the German mathematician L. Matthiessen first pointed out that Sunzi's solution was consistent with Gauss's method.

}{%

في أوروبا، لم تُطوَّر طرق مماثلة حتى عمل أويلر (١٧٤٣) وغاوس (١٨٠١). في عام ١٨٥٢، نشر المبشر البريطاني ألكسندر وايلي ملاحظات عن العلوم الصينية، وقدَّم هذه المسألة إلى الرياضيين الأوروبيين. وفي عام ١٨٧٦، أشار الرياضي الألماني ل. ماتيسين لأول مرة إلى أن حل سونز يتوافق مع طريقة غاوس.

}



\newpage



\subsection*{The ``Sunzi Song'' — \ar{أغنية سونز} (\zh{孫子歌})}



\begin{songbox}

\zh{三人同行七十稀，五樹梅花廿一枝，}\\[4pt]

\zh{七子團圓正半月，除百零五便得知。}

\end{songbox}



\vspace{0.3cm}



\bipassage{%

\textit{``Three people walking together---seventy is rare; five plum trees---twenty-one branches; seven sons reunite---half a full moon (15); subtract one hundred and five, then you will know.''}

}{%

\ar{ثلاثة يمشون معًا — السبعون نادرة؛ خمس شجرات خوخ — واحد وعشرون غصنًا؛ سبعة أبناء يجتمعون — نصف قمر كامل (١٥)؛ اطرح مئة وخمسة، فستعرف.}

}



\bipassage{%

This gives the key numbers: \textbf{70} for mod 3, \textbf{21} for mod 5, \textbf{15} for mod 7, and \textbf{105} as the LCM.

}{%

هذا يعطي الأرقام الرئيسية: \textbf{٧٠} للقسمة على ٣، و\textbf{٢١} للقسمة على ٥، و\textbf{١٥} للقسمة على ٧، و\textbf{١٠٥} كمضاعف مشترك أصغر.

}



\newpage



%% ========================================================================

%% REMAINING PROBLEMS FROM VOLUME III

%% ========================================================================

\section*{Remaining Problems from Volume III — \ar{مسائل أخرى من الجزء الثالث}}

\addcontentsline{toc}{section}{Remaining Problems — \ar{مسائل أخرى}}



\subsection*{Problem 27. Birds and Beasts — \ar{المسألة السابعة والعشرون: الطيور والوحوش}}



\biproblem{%

\zh{今有獸六首四足，禽四首二足。上有七十六首，下有四十六足。問禽、獸各幾何？}%

}{%

لدينا وحش بستة رؤوس وأربع قوائم، وطائر بأربعة رؤوس وقدمين. فوق ستة وسبعون رأسًا، وتحت ست وأربعون قدمًا. كم طائرًا وكم وحشًا؟

}

\bianswer{%

\zh{答曰：八獸，七禽。}%

}{%

الجواب: ثمانية وحوش، وسبعة طيور.

}

\bimethod{%

\zh{術曰：倍足以減首，餘，半之，即禽。四乘禽數，減足，餘，半之，即獸。}%

}{%

الطريقة: اضعف الأقدام واطرح الرؤوس، ونصف الباقي فهو عدد الطيور. اضرب عدد الطيور في أربعة واطرح من الأقدام، ونصف الباقي فهو عدد الوحوش.

}



\subsection*{Problem 29. One Hundred Deer — \ar{المسألة التاسعة والعشرون: مئة أيل}}



\biproblem{%

\zh{今有百鹿入城，家取一鹿，不盡；又三家共一鹿，適盡。問城中家幾何？}%

}{%

دخلت مئة أيل المدينة، فأخذت كل بيت أيلاً واحداً ولم يكفِ؛ ثم اقتسمت ثلاثة بيوت كل أيل وكفى بالضبط. كم بيتًا في المدينة؟

}

\bianswer{%

\zh{答曰：七十五家。}%

}{%

الجواب: خمسة وسبعون بيتًا.

}

\bimethod{%

\zh{術曰：以盈不足取之。假令七十二家，鹿不盡四。令之九十家，鹿不足二十。置七十二於右上，盈四於右下。置九十於左上，不足二十於左下。維乘之，所得，并為實。并盈不足為法。除之，即得。}%

}{%

الطريقة: بالزيادة والنقص. لنفرض اثنتين وسبعين بيتًا، فالأيائل تبقى أربعة. ولنفرض تسعين بيتًا، فالأيائل تنقص عشرون. ضع اثنتين وسبعين في الأعلى الأيمن، والزيادة أربعة في الأسفل الأيمن. ضع تسعين في الأعلى الأيسر، والنقص عشرون في الأسفل الأيسر. اضرب متقاطعًا واجمع للمقسوم. واجمع الزيادة والنقص للقاسم. اقسم فتحصل على النتيجة.

}



\subsection*{Problem 30. Three Chickens — \ar{المسألة الثلاثون: الدجاج الثلاث}}



\biproblem{%

\zh{今有三雞共啄粟一千一粒。雛啄一，母啄二，翁啄四。主責本粟，三雞主各償幾何？}%

}{%

لدينا ثلاث دجاجات تلتهم ألفًا وواحدة حبة ذُرَة. الكتكوت يلتهم واحدة، والدجاجة اثنتين، والديك أربع. يريد صاحب الذرة ردَّها، فكم يدفع كل صاحب دجاجة؟

}

\bianswer{%

\zh{答曰：雞雛主一百四十三。雞母主二百八十六。雞翁主五百七十二。}%

}{%

الجواب: صاحب الكتكوت مئة وثلاث وأربعون. صاحب الدجاجة مئتان وست وثمانون. صاحب الديك خمسمئة واثنتان وسبعون.

}



\subsection*{Problem 33. Wheel Turns — \ar{المسألة الثالثة والثلاثون: دورات العجلة}}



\biproblem{%

\zh{今有長安、洛陽相去九百里。車輪一匝一丈八尺。欲自洛陽至長安，問輪匝幾何？}%

}{%

بين مدينتي تشْانْغآن ولويْيانْغ تسعمئة لِي. دورة عجلة العربة \textit{جانْغ} واحد وثمانية أقدام. إذا أردت السفر من لويْيانْغ إلى تشْانْغآن، كم دورة للعجلة؟

}

\bianswer{%

\zh{答曰：九萬匝。}%

}{%

الجواب: تسعون ألف دورة.

}



\newpage



%% ========================================================================

%% EPILOGUE

%% ========================================================================

\section*{Epilogue — \ar{خاتمة}}

\addcontentsline{toc}{section}{Epilogue — \ar{خاتمة}}



\bipassage{%

The \textit{Sunzi Suanjing} stands as a remarkable testament to the sophistication of ancient Chinese mathematics. While the text covers basic arithmetic---units of measurement, multiplication tables, and simple fractions---its greatest legacy lies in the brief but brilliant ``Problem of the Unknown Quantity'' that opens the door to modern number theory.

}{%

يقف كتاب سونز في الحساب شهادةً رائعة على تطور الرياضيات الصينية القديمة. ورغم أن النص يغطي الحساب الأساسي — وحدات القياس وجداول الضرب والكسور البسيطة — فإن أعظم إرث له يكمن في المسألة الموجزة والبراقة «مسألة العدد المجهول» التي تفتح الباب لنظرية الأعداد الحديثة.

}



\bipassage{%

The Chinese Remainder Theorem, as this method came to be known in the West, would not be fully formalized in Europe for over a millennium. The algorithm described by Sunzi contains the essential ingredients of the modern theorem: the construction of auxiliary numbers with specific divisibility properties, their combination with the given remainders, and reduction modulo the least common multiple.

}{%

مسألة الباقي الصينية، كما عُرِفت هذه الطريقة في الغرب، لم تُصاغ رسميًّا في أوروبا لأكثر من ألف عام. الخوارزمية التي وصفها سونز تحتوي على المكونات الأساسية للنظرية الحديثة: بناء أعداد مساعدة بخصائص قابلية قسمة محددة، ودمجها مع البقايا المعطاة، والتخفيض بالمضاعف المشترك الأصغر.

}



\bipassage{%

The influence of this text extended far beyond China. The problem circulated as a mathematical puzzle throughout East Asia, and the general method developed by Qin Jiushao represents one of the great achievements of medieval mathematics anywhere in the world.

}{%

تأثير هذا النص امتد بعيدًا خارج الصين. انتشرت المسألة كأحجية رياضية في جميع أنحاء شرق آسيا، والطريقة العامة التي طوَّرها تْشِن جْيوشَاو تمثل أحد أعظم الإنجازات الرياضية في العصور الوسطى في أي مكان في العالم.

}



\vspace{1cm}

\begin{center}

\longline

\end{center}



\vspace{0.5cm}



\begin{center}

\textit{Typeset from the classical Chinese text — \ar{مُنضَّدٌ من النص الصيني الكلاسيكي}}\\[4pt]

\textit{Chinese Text Project} (\texttt{ctext.org/sunzi-suan-jing})

\end{center}



\end{document}







% ============================================================

% Source: translations/ar/chinese/test_framed.tex

% ============================================================



\documentclass{article}

\usepackage{fontspec}

\usepackage{polyglossia}

\usepackage{xeCJK}

\usepackage{xcolor}

\usepackage{framed}

\usepackage{color}

\begin{document}

\begin{framed}

Hello framed

\end{framed}

\begin{shaded}

Hello shaded

\end{shaded}

Test content

\end{document}







% ============================================================

% Source: translations/ar/chinese/test_mdframe.tex

% ============================================================



\documentclass{article}

\usepackage{fontspec}

\usepackage{polyglossia}

\usepackage{xeCJK}

\usepackage{xcolor}

\usepackage{mdframed}

\usepackage{hyperref}

\begin{document}

\begin{mdframed}[backgroundcolor=red!5,linewidth=1pt,linecolor=red!75!black]

Hello

\end{mdframed}

Test content

\end{document}







% ============================================================

% Source: translations/ar/chinese/test_minimal.tex

% ============================================================



\documentclass{article}

\usepackage{fontspec}

\usepackage{polyglossia}

\usepackage{xeCJK}

\usepackage{hyperref}

\begin{document}

Test

\end{document}







% ============================================================

% Source: translations/ar/chinese/test_tcolor.tex

% ============================================================



\documentclass{article}

\usepackage{fontspec}

\usepackage{polyglossia}

\usepackage{xeCJK}

\usepackage{xcolor}

\usepackage{tcolorbox}

\usepackage{hyperref}

\begin{document}

\begin{tcolorbox}[colback=red!5,colframe=red!75!black,title=Test]

Hello

\end{tcolorbox}

Test content

\end{document}







% ============================================================

% Source: translations/ar/chinese/test_tcolor2.tex

% ============================================================



\documentclass{article}

\usepackage{fontspec}

\usepackage{polyglossia}

\usepackage{xeCJK}

\usepackage{xcolor}

\usepackage{tcolorbox}

\tcbuselibrary{most}

\usepackage{hyperref}

\begin{document}

\begin{tcolorbox}[colback=red!5,colframe=red!75!black,title=Test]

Hello

\end{tcolorbox}

Test content

\end{document}







% ============================================================

% Source: translations/ar/chinese/test_tcolor3.tex

% ============================================================



\documentclass{article}

\usepackage{fontspec}

\usepackage{polyglossia}

\usepackage{xeCJK}

\usepackage{xcolor}

\usepackage{tcolorbox}

\usepackage[colorlinks=false,draft]{hyperref}

\begin{document}

\begin{tcolorbox}[colback=red!5,colframe=red!75!black,title=Test]

Hello

\end{tcolorbox}

Test content

\end{document}







% ============================================================

% Source: translations/ar/chinese/test_tcolor4.tex

% ============================================================



\documentclass{article}

\usepackage{fontspec}

\usepackage{polyglossia}

\usepackage{xeCJK}

\usepackage{xcolor}

\usepackage{tcolorbox}

\begin{document}

\begin{tcolorbox}[colback=red!5,colframe=red!75!black,title=Test]

Hello

\end{tcolorbox}

Test content

\end{document}







% ============================================================

% Source: translations/ar/islamic/al-khwarizmi-al-jabr_arabic-enhanced.tex

% ============================================================



%% ========================================================================

%% Kitab al-Mukhtasar fi Hisab al-Jabr wa-l-Muqabala

%% by Muhammad ibn Musa al-Khwarizmi (c. 820 CE)

%% Enhanced Arabic Scholarly Edition

%% ========================================================================



\documentclass[11pt,a4paper]{book}



%% -------- Core Packages --------

\usepackage{fontspec}

\usepackage{polyglossia}

\usepackage{amsmath,amssymb,amsthm}

\usepackage{geometry}

\usepackage{graphicx}

\usepackage{xcolor}

\usepackage{titlesec}

\usepackage{fancyhdr}

\usepackage{enumitem}

\usepackage{array,booktabs,longtable,multirow}

\usepackage{hyperref}

\usepackage{lettrine}

\usepackage{microtype}

\usepackage{tocloft}

\usepackage{indentfirst}

\usepackage{tcolorbox}



%% -------- Page Geometry --------

\geometry{

  paperwidth=210mm,paperheight=297mm,

  left=22mm,right=22mm,top=28mm,bottom=28mm,

  headheight=14pt,footskip=30pt

}



%% -------- Language Setup --------

\setmainlanguage[numerals=maghrib]{arabic}

\setotherlanguage{english}



%% -------- Font Configuration --------

\setmainfont{Arial}



\newfontfamily\latinfont{Times New Roman}

\newfontfamily\latinfontsf{Arial}

\setmonofont{Courier New}



%% Arabic display font

\newfontfamily\arabicdisplayfont[Script=Arabic,Scale=1.1]{Arial}



%% -------- Colors --------

\definecolor{titlecolor}{RGB}{45,55,72}

\definecolor{sectioncolor}{RGB}{55,65,81}

\definecolor{notecolor}{RGB}{120,53,15}

\definecolor{rulecolor}{RGB}{180,160,140}

\definecolor{arabicbg}{RGB}{252,250,245}

\definecolor{dropgold}{RGB}{139,90,43}

\definecolor{commentarybg}{RGB}{245,248,250}

\definecolor{originalbg}{RGB}{255,253,245}



%% -------- Section Formatting --------

\titleformat{\chapter}[display]

  {\normalfont\huge\bfseries\color{titlecolor}}

  {\chaptertitlename\ \thechapter}{20pt}{\Huge}

\titleformat{\section}

  {\normalfont\Large\bfseries\color{sectioncolor}}

  {\thesection}{1em}{}

\titleformat{\subsection}

  {\normalfont\large\bfseries\color{sectioncolor}}

  {\thesubsection}{1em}{}



%% -------- Header/Footer --------

\pagestyle{fancy}

\fancyhf{}

\fancyhead[LE]{\small\thepage}

\fancyhead[RO]{\small\thepage}

\fancyhead[C]{\small Kitab al-Jabr wa-l-Muqabala -- al-Tab'a al-Arabiyya al-Muhaqqaqa}

\renewcommand{\headrulewidth}{0.4pt}

\renewcommand{\footrulewidth}{0pt}



%% -------- Custom Commands --------

\newcommand{\longline}{\noindent\rule{\textwidth}{0.4pt}}

\newcommand{\shortline}{\noindent\rule{0.5\textwidth}{0.4pt}\par}



%% -------- tcolorbox styles --------

\tcbuselibrary{skins,breakable}



\newtcolorbox{originaltext}[1]{

  colback=originalbg,

  colframe=rulecolor,

  arc=3pt,

  boxrule=0.6pt,

  left=12pt,right=12pt,top=8pt,bottom=8pt,

  fontupper=\normalsize,

  breakable,

  title={\small\textbf{al-Nass al-Asli li-l-Khwarizmi}},

  coltitle=dropgold,

  colbacktitle=arabicbg,

  fonttitle=\small,

  #1

}



\newtcolorbox{commentary}[1]{

  colback=commentarybg,

  colframe=sectioncolor,

  arc=2pt,

  boxrule=0.5pt,

  left=10pt,right=10pt,top=6pt,bottom=6pt,

  fontupper=\small,

  breakable,

  title={\small\textbf{Taqliq Tahqiqi}},

  coltitle=sectioncolor,

  fonttitle=\small,

  #1

}



\newtcolorbox{mathexplain}[1]{

  colback=white,

  colframe=notecolor,

  arc=2pt,

  boxrule=0.6pt,

  left=10pt,right=10pt,top=6pt,bottom=6pt,

  fontupper=\small,

  breakable,

  title={\small\textbf{al-Tawdih al-Riyadi}},

  coltitle=notecolor,

  fonttitle=\small,

  #1

}



\newtcolorbox{historicalbox}[1]{

  colback=arabicbg,

  colframe=dropgold,

  arc=2pt,

  boxrule=0.5pt,

  left=10pt,right=10pt,top=6pt,bottom=6pt,

  fontupper=\small,

  breakable,

  title={\small\textbf{Siyaq Tarikhi}},

  coltitle=dropgold,

  fonttitle=\small,

  #1

}



\newtcolorbox{geombox}[1]{

  colback=white,

  colframe=sectioncolor,

  arc=2pt,

  boxrule=0.6pt,

  left=10pt,right=10pt,top=6pt,bottom=6pt,

  fontupper=\small,

  breakable,

  title={\small\textbf{al-Tashrih al-Hindi}},

  coltitle=sectioncolor,

  fonttitle=\small,

  #1

}



%% -------- Hyperref Setup --------

\hypersetup{

  colorlinks=true,

  linkcolor=sectioncolor,

  citecolor=sectioncolor,

  urlcolor=sectioncolor,

  pdfencoding=auto,

  pdftitle={Kitab al-Jabr wa-l-Muqabala - Arabic Enhanced Edition},

  pdfauthor={al-Khwarizmi}

}



\setcounter{tocdepth}{2}

\setcounter{secnumdepth}{2}



%% ========================================================================

\begin{document}



%% ========================================================================

%% TITLE PAGE

%% ========================================================================

\begin{titlepage}

\begin{center}



\vspace*{1.5cm}



{\fontsize{28}{34}\selectfont\color{titlecolor} Kitab al-Mukhtasar fi Hisab al-Jabr wa-l-Muqabala}



\vspace{0.5cm}



{\LARGE\textbf{li-Muhammad ibn Musa al-Khwarizmi}}



\vspace{0.3cm}



{\normalsize (taqriban 780--850 M)}



\vspace{1.5cm}



\begin{tcolorbox}[colback=arabicbg,colframe=dropgold,arc=4pt,boxrule=1pt,width=0.85\textwidth]

\begin{center}

\large

\textbf{al-Tab'a al-Arabiyya al-Muhaqqaqa al-Mu'azzaza}



\vspace{0.2cm}



Enhanced Arabic Scholarly Edition



\vspace{0.2cm}



al-Nass al-Asli ma'a Ta'liqat Tahqiqiyya wa-Shuruh Riyadiyya

\end{center}

\end{tcolorbox}



\vspace{1.5cm}



{\large Allifa hadha al-Kitab 'inda Balat al-Khalifa al-Ma'mun fi Baghdad}



\vspace{0.5cm}



{\normalsize Hawali Sanat 200 H / 820 M}



\vspace{1.5cm}



\rule{0.3\textwidth}{0.5pt}



\vspace{0.8cm}



{\footnotesize\latinfont This enhanced edition prepared with:}



{\footnotesize al-Nass al-Asli al-Arabi -- Ta'liqat Tahqiqiyya -- Tawdihat Riyadiyya -- Siyaq Tarikhi}



\vspace{0.5cm}



{\footnotesize\latinfont Based on the critical edition of}



{\footnotesize 'Ali Mustafa Musharrafa wa-Muhammad Mursi Ahmad (al-Qahira, 1939)}



\end{center}

\end{titlepage}



%% ========================================================================

%% ABOUT THIS EDITION

%% ========================================================================

\thispagestyle{empty}

\begin{center}

\vspace*{2cm}



{\Large\textbf{'An Hadhihi al-Tab'a}}



\vspace{1cm}



\begin{tcolorbox}[colback=arabicbg,colframe=rulecolor,arc=3pt,boxrule=0.5pt,width=0.88\textwidth]

\small

Tuqaddimu hadhihi al-tab'a al-Kitab al-Asli fi 'ilm al-Jabr li-Muhammad ibn

Musa al-Khwarizmi, alladhi allafahu fi Balat al-Khalifa al-Ma'mun fi Baghdad

hawali sanat 820 M. Yashtamil hadha al-Isdar al-Mu'azzaz 'ala al-nass al-Arabi

al-asli ma'a ta'liqat tahqiqiya mufassala wa-sharh riyadi wa-tawdihat

hindisiya wa-siyaq tarikhi -- kulluha bi-l-lugha al-Arabiyya.



\vspace{0.3cm}



al-Nass al-asli ma'khudh min al-tab'a al-tahqiqiya li-'Ali Mustafa Musharrafa

wa-Muhammad Mursi Ahmad (al-Qahira, 1359 H / 1939 M), allati istanadat ila

al-makhtuta al-wahida al-mahfuza fi maktabat Bodleyan bi-Oxford (Hunt. 214,

nusikhat sanat 743 H / 1342 M).

\end{tcolorbox}



\end{center}



\newpage



%% ========================================================================

%% TABLE OF CONTENTS

%% ========================================================================

\tableofcontents

\newpage



%% ========================================================================

%% PREFACE

%% ========================================================================

\chapter*{Taqdimat al-Tab'a al-Mu'azzaza}

\addcontentsline{toc}{chapter}{Taqdimat al-Tab'a al-Mu'azzaza}

\markboth{Taqdimat al-Tab'a}{Taqdimat al-Tab'a}



\begin{historicalbox}

Tu'tabar ummat al-Islam min aghna al-umam turathan 'ilmiyyan, wa-qad

 Tatawaqqalat 'alayha hadarat mukhtalifa min fajr al-tarikh ila al-yawm.

Wa-fi kulli hadara minha qamat bi-qasT wafir min wajibiha al-'ilmi nahwa

al-usra al-bashariyya. Wa-laysa yakfi an natahaddath 'an majdina al-'ilmi

kama law kana usTura, bal yajib an yazhar hadha al-majd fi sura

malmusa tara'ha al-a'yun wa-tunaluha al-aydi. Li-dhalika kana min

al-muhimm an na'na bi-nashr al-kutub allati wada'aha aba'una

wa-ajdaduna, khususan idha kanat hadhihi al-kutub hamma al-athar fi

takyif al-tafkir al-bashari. Wa-la shakka anna fi muqaddimati hadhihi

al-kutub \textbf{Kitab al-Khwarizmi fi al-Jabr wa-l-Muqabala}.

\end{historicalbox}



\section*{Ahimmiyyat Hadha al-Kitab fi Tarikh al-'Ilm}



\begin{commentary}

\textbf{Kitab al-Khwarizmi fi al-Jabr wa-l-Muqabala} huwa awwal kitab

munaZZam fi 'ilm al-jabr 'ala al-itlaq. Wa-huwa yamthilu al-laHZa allati

taHawwala fiha al-jabr min majmu'at min al-qawa'id al-munfasila ila 'ilm

riyadi munazzam lahu usuluhu wa-qawa'iduhu al-khassa. Laqad adkhala

al-Khwarizmi niZaman lil-mu'adlat min al-daraja al-thaniya bina'an 'ala

thalathat anwa' min al-kamiyyat: \textbf{al-mal} (al-majhul al-tarb'i),

wa-\textbf{al-jadhr} (al-majhul al-khatti), wa-\textbf{al-'adad} (al-thabit).

Wa-min hadhihi al-anwa' al-thalatha istanbaTa \textbf{sitt haSalat}

kanuniyya taghatti kulli ashkal al-mu'adlat al-tarb'iyya dhat al-ma'amilat

al-mujiba.



\vspace{0.3cm}



\textbf{MustalaHat al-Jabr al-Asasiyya:}

\begin{itemize}[rightmargin=1em]

  \item \textbf{al-Jabr} (al-jabr): idafat ma yanquSu min ahadi

    Tarfay al-mu'adla ila al-Taraf al-akhar li-istikmal al-tawazun

  \item \textbf{al-Muqabala} (al-muqabala): hadf ma yatashabahu min

    al-Tarfayn, ay Tarq al-mutashabihat min kila Tarfi al-mu'adla

    li-tashil

  \item \textbf{al-Mal} (al-mal): mirab' al-majhul, ay $x^2$

  \item \textbf{al-Jadhr} (al-jadhr): al-majhul dhatuhu, ay $x$

  \item \textbf{al-'Adad} (al-'adad): al-hadd al-thabit fi al-mu'adla

\end{itemize}

\end{commentary}



\newpage



%% ========================================================================

%% AUTHOR'S PREFACE

%% ========================================================================

\chapter*{Kalamat al-Musannif: al-Muqaddima al-Asliyya}

\addcontentsline{toc}{chapter}{Kalamat al-Musannif: al-Muqaddima al-Asliyya}

\markboth{al-Muqaddima al-Asliyya}{al-Muqaddima al-Asliyya}



\begin{originaltext}

Bismi Allahi al-Rahmani al-Rahim.



Al-Hamdu lillahi 'ala ni'amihi li-ahli Ta'atihi, bi-ada'iha bi-qudratihi

'ala 'ibadihi al-mu'minina bihi, li-yuthayyibahum biha thawaban,

wa-yuHafiZahum biha 'an al-taghyir, wa-ya'rifahum qudratahu,

wa-yudhillilahum li-'iZmatihi. Arsala Muhammadan Salla Allahu 'alayhi

wa-Sallam bi-risalatihi ba'da inqiTa' al-rusul, wa-dhihab al-'adl,

wa-fasad al-nuSul.



Wa-qad kana al-'ulama' fi al-azmani al-khaliya wa-fi al-umam al-salifa

yu'allifuna al-kutub fi abwab al-'ulum wa-furu' al-ma'rifa, yatadhakkaruna

min ba'dihim, wa-yarjuna al-ajra bi-qadri ma kana minhum, wa-yathiquna

bi-ma yaHasil lahum min al-shukr wa-al-ijlal wa-al-i'Zam, raDiina

wa-law bi-qalilin min al-thina', qalilan muqarana bi-ma badhuluhu min

al-jahd wa-al-mushaqqa fi kashf khafa'a al-'ilmi wa-izalat satayrihi.



Wa-dhalika al-Hubb li-l-'ilmi alladhi mayyaza Allahu bihi al-Imam

al-Ma'mun Amir al-Mu'minin (ma'a ma manna bihi 'alayhi min al-khilafa

bi-al-ijma', wa-al-libas alladhi albasahu iyahu, wa-al-zina allati

zayyanahu biha), wa-dhalika al-laTf wa-al-tawaddud alladhi yubdihi

li-l-'ulama', wa-tilka al-sur'a allati yuHamiha wa-yunSiha fi tawDiH

al-mubhamat wa-izalat al-Su'ubat -- qad Hamalni 'ala an u'allif

mukhtaSan fi al-Jabr wa-l-Muqabala, aqtasiru fihi 'ala ma huwa aysar

wa-anfa' fi al-Hisab, mimma yaHtaju ilayhi al-nasu da'iman fi masa'il

al-mirath wa-al-'aTaya wa-al-qisma wa-al-khuSumat wa-al-tijara wa-kulli

ma yata'amaluna bihi fima baynahum, aw fi misaHat al-arDiyn wa-Hufri

al-anhar wa-masa'il al-handasa wa-ghayriha min anwa' al-ashya'

al-mukhtalifa al-aSnaf.

\end{originaltext}



\begin{commentary}

\textbf{al-Ta'liq al-Tahqiqi:}



Yastahillu al-Khwarizmi kitabahu bi-al-basmala al-taqlidiyya, thumma

yuthni 'ala al-Khalifa al-Ma'mun (198--218 H / 813--833 M), wa-huwa min

ashhura khulafa' Bani al-'Abbas tashji'an lil-'ilm. Laqad assasa

al-Ma'mun \textbf{Bayt al-Hikma} fi Baghdad, wa-ja'alahu markazan

'alamiyyan lil-baHth al-'ilmi, wa-jama'a fihi al-mutarjimina

wa-al-'ulama' min al-muslimina wa-ghayrihim.



Yubayyinu al-Khwarizmi dafi'ahu li-ta'lif hadha al-Kitab: al-Haja

al-mulahha li-uslub manZum fi Hal al-masa'il al-'amaliyya allati

yuwajihuha al-nas yawmiyyan -- min qismat al-tarakat ila misaHat

al-aradi wa-Hufri al-qanawat. Wa-hadha yuwiZZahu anna al-jabr 'inda

al-Khwarizmi lam yakun mujarrad niZariyya mujarrada, bal kana

\textbf{adatan 'amaliyya} muWajahatan nahwa Hall al-mushkilat

al-waqi'iyya.

\end{commentary}



\begin{historicalbox}

\textbf{al-Imam al-Ma'mun Amir al-Mu'minin (786--833 M):}



Kana al-Ma'mun min akthar al-khulafa' tashji'an lil-'ilm wa-al-'ulama'.

Amara bi-tarjamat al-kutub al-Yunaniyya wa-al-Hindiyya wa-al-Suryaniyya

ila al-Arabiyya, wa-assasa Bayt al-Hikma fi Baghdad allati kanat 'aSima

al-'ilmi fi al-'alam. Wa-fi 'ahdihi taqallada al-Khwarizmi mansaBan fi

Bayt al-Hikma, haythu allafa hadha al-kitab wa-ghayrahu min al-muSannafat

fi al-riyadiyyat wa-al-falak wa-al-jughrafiya.

\end{historicalbox}



\newpage



%% ========================================================================

%% CHAPTER: ON NUMBERS AND THE THREE KINDS

%% ========================================================================

\chapter{fi al-A'dad wa-al-Anwa' al-Thalatha}

\label{ch:intro}



\begin{originaltext}

Lamma naZartu fima yaHtaju ilayhi al-nasu fi Hisabatihim, wajadtu

anna dhalika kulluhu 'adad. Wa-wajadtu aydan anna kullu 'adadin

mu'allaf min al-waHadat, wa-annahu la yakhlu 'adad min an yanqasim

ila waHadat. Wa-'alimtu anna kulla 'adadin yu'adda min al-waHid ila

al-'ashara yazidu 'ala alladhi qablahu bi-waHid, thumma yuDa'afu

al-'ashru wa-yuthallathu, kama fu'ila bi-l-waHadat, fa-yakunu

'ishrun wa-thalathun wa-haka-dha ila al-ma'iyya, thumma yuDa'afu

al-ma'iyya wa-yuthallathuha kama fu'ila bi-l-'asharat wa-l-waHadat ila

al-alf, thumma yukarru al-alfu kadhaka fi kulli 'adadin murakkab,

wa-haka-dha ila aqSa Hadin fi al-'add.

\end{originaltext}



\begin{commentary}

\textbf{al-NiZam al-'Adadi 'inda al-Khwarizmi:}



Yabda'u al-Khwarizmi kitabahu bi-munaqasha falsafiyya li-l-'adad,

mustakhdiman al-niZam al-'adadi al-'ashari alladhi yu'rafu al-yawma

bi-``al-arqam al-Arabiyya'' (wa-hiya fi al-aSl min aSl Hind).

Yasifu al-Khwarizmi al-niZam al-marakaZi li-l-a'dad Haythu kullu

mawDi' yamthilu quwwa min 'ashara. Hadhihi al-muqaddima laysat

majarrad tamhid, bal hiya i'lan 'an \textbf{al-niZam al-'adadi

al-mawDi'i} alladhi sa-yaj'alu kull al-Hisabati al-taliqa mumkina.



Laqad istaw'abba al-Khwarizmi anna kulla masa'ala Hisabiyya -- sawa'

fi al-tijara aw al-misaHa aw al-mirath -- yumkinu irja'uha fi

al-nihaya ila mafhum al-'adad, wa-min thumma yumkinu Halluha bi-l-Turuq

al-riyadiyya.

\end{commentary}



\section{al-Anwa' al-Thalatha min al-A'dad fi al-Jabr}



\begin{originaltext}

Fa-wajadtu anna al-a'dada allati yaHtaju ilayha fi Hisab al-Jabr

wa-l-Muqabala thalathatu anwa': \textbf{juthur}, wa-\textbf{amwal},

wa-\textbf{a'dad mufarrada} la tat'allaqu bi-jadhrin wa-la bi-malin.



Amma al-jadhr fa-huwa kullu miqdarin yuDribu fi nafsihi mu'allafan

min waHadatin aw a'dadin sa'ida aw kusurin nazila. Wa-amma al-mal

fa-huwa jami'u ma Hasala min Darb al-jadhr fi nafsihi. Wa-amma

al-'adad al-mufrad fa-huwa kullu 'adadin yudhkaru bi-ghayri nisba ila

jadhrin wa-la malin.

\end{originaltext}



\begin{mathexplain}

\textbf{al-Anwa' al-Thalatha bi-l-Rumuz al-Haditha:}



\begin{center}

\begin{tabular}{|c|c|c|l|}

\hline

\textbf{al-Isem al-Arabi} & \textbf{al-Ramz al-Hadith} & \textbf{al-Mathal} & \textbf{al-Ta'rif} \\

\hline

al-mal ($x^2$) & mal & $x^2$ & mirab' al-majhul \\

\hline

al-jadhr ($x$) & jadhr & $x$ & al-kamiyya al-majhula \\

\hline

al-'adad ($c$) & 'adad & 39 & al-thabit al-'adadi \\

\hline

\end{tabular}

\end{center}



\vspace{0.3cm}



\textbf{Mulahadha:} kalimat \textbf{``mal''} fi al-iStilaH al-jabri

'inda al-Khwarizmi ta'ni \textit{mirab' al-majhul}. Wa-amma

\textbf{``al-shay' ''} (alladhi yumaththilu ``al-jadhr'')

fa-huwa al-majhul al-khatti. Wa-l-jadhr huwa alladhi tunsha'u

minhu al-amwal bi-l-Darb fi al-nafs. Wa-l-'adad al-mufrad huwa

kullu 'adadin la yatagayyaru bi-taghyyuri al-majhul.

\end{mathexplain}



\begin{commentary}

\textbf{Tafsir al-MustalaHat al-Jabriyya:}



Kalimat \textbf{``al-jadhr''} fi hadha al-siyaq tuKhTallifu 'an

al-jadhr al-tarb'i kama nafhamuhu al-yawm. Fa-l-jadhr 'inda

al-Khwarizmi huwa \textbf{al-majhul dhatuhu}, ay $x$ fi al-rumuz

al-Haditha. Wa-l-mal huwa \textbf{tarb' hadha al-majhul}, ay $x^2$.



Wa-qad ikhtara al-Khwarizmi kalimat \textbf{``mal''} rubbama

li-annana mirab' al-majhul yumaththilu \textit{al-tharwa} aw

\textit{al-mumtalakat} fi al-tafkir al-iqtiSadi li-l-'asr -- fa-huwa

al-kamiyya al-akbar wa-al-ahamm fi al-mu'adala.

\end{commentary}



\newpage



%% ========================================================================

%% CHAPTER: THE SIX CANONICAL CASES

%% ========================================================================

\chapter{al-HaSalat al-Sitt al-Kanuniyya lil-Mu'adlat}

\label{ch:six-cases}



\begin{originaltext}

Hadhi-hi hiya al-HaSalat al-sitt alladhi dhakartuha fi muqaddimati

hadha al-kitab. Thalathun minha la taHtaju ila an tunaSafa

al-juthuru, wa-qad bayyantu kayfa tujab. Wa-amma al-thalath

al-ukhra allati yajibu fiha naSf al-juthuri, fa-ara an ubayyinaha

aydan mufassaLan ma'a dhikri al-shakl li-kulli haSala yubayyin

sabab al-naSf.

\end{originaltext}



\begin{commentary}

YuQassimu al-Khwarizmi jami'a al-mu'adlat al-tarb'iyya ila

\textbf{sitt haSalat kanuniyya} -- thalath \textit{basita} wa-thalath

\textit{murakkaba}. Hadha al-taqsimu yughatti jami'a ashkal

al-mu'adlat al-tarb'iyya dhat al-ma'amilat al-mujiba. Wa-l-sabab fi

al-iqtiSir 'ala al-ma'amilat al-mujiba huwa anna al-Khwarizmi yafhamu

kulla Haddin fi al-mu'adala 'ala annahu \textbf{kamiyya hindisiyya}

(Tul aw misaHa), wa-la yumkinu an takuna al-misaHa saliba.



\vspace{0.3cm}



\textbf{al-Tamayyuz bayn al-HaSalat al-Basita wa-l-Murakkaba:}

\begin{itemize}

  \item al-HaSalat al-basita (al-ula wa-l-thaniya wa-l-thalitha): la

taHtaju ila \textbf{naSf al-juthur} -- yumkinu Halluha bi-l-tashil

al-mubashshar.

  \item al-HaSalat al-murakkaba (al-rabi'a wa-l-khamisa wa-l-sadisa):

tataTallaBu \textbf{naSf al-juthur}, wa-hadhihi hiya al-HaSalat

allati yaZharu fiha \textbf{takmil al-murabba'} (completing the square).

\end{itemize}

\end{commentary}



\section{al-HaSalat al-Basita al-Thalath}



\begin{mathexplain}

\textbf{al-HaSalat al-Basita (la tatallaBu naSf al-juthur):}



\begin{center}

\begin{tabular}{|c|c|c|l|}

\hline

\textbf{al-HaSala} & \textbf{al-Mu'adala} & \textbf{al-Hall} & \textbf{al-Isem al-Arabi} \\

\hline

al-Ula & $ax^2 = bx$ & $x = b/a$ & mal yusawi usuLan \\

\hline

al-Thaniya & $ax^2 = c$ & $x = \sqrt{c/a}$ & mal yusawi 'adadan \\

\hline

al-Thalitha & $bx = c$ & $x = c/b$ & usuL tusawi 'adadan \\

\hline

\end{tabular}

\end{center}

\end{mathexplain}



\subsection{al-HaSala al-Ula: al-mal yusawi usuLan}



\begin{originaltext}

\textbf{Mathal:} mal yusawi khamsata usuLihi. Jadhr al-mal khamsa,

wa-l-mal khamsa wa-'ishrun, wa-huwa yusawi khamsata a'da' jadhr.



\textbf{Mathal akhar:} thulthu mal yusawi arba'ata usuLin. Fa-l-mal

kulluh yusawi ithnay 'ashara jadhra; ay mi'a wa-arba'a wa-arba'in,

wa-jadhrhu ithna 'ashar.



\textbf{Mathal thalith:} khamsatu amwalin tusawi 'ashara usuLan.

Fa-mal wahid yusawi jadhrayn; jadhr al-mal ithnan, wa-malahu

arba'a.

\end{originaltext}



\begin{mathexplain}

\textbf{al-Tawdih al-Riyadi lil-HaSala al-Ula:}



\begin{center}

$\begin{aligned}

  x^2 &= 5x \quad\Rightarrow\quad x = 5, \quad x^2 = 25 \\

  \tfrac{1}{3}x^2 &= 4x \quad\Rightarrow\quad x^2 = 12x = 144, \quad x = 12 \\

  5x^2 &= 10x \quad\Rightarrow\quad x^2 = 2x = 4, \quad x = 2

\end{aligned}$

\end{center}



\textbf{al-Tariqa al-'Amma:} fi al-mu'adala $ax^2 = bx$, naqsimu

al-Tarfayn 'ala $ax$ (bi-afTariD $x \neq 0$) fa-nuHallu 'ala

$x = b/a$. Amma Hall $x = 0$ fa-kana yutajahalu taqlidiyan li-annahu

la ma'na 'amaliyya lahu fi muZamm al-masa'il.

\end{mathexplain}



\subsection{al-HaSala al-Thaniya: al-mal yusawi 'adadan}



\begin{originaltext}

\textbf{Mathal:} mal yusawi tis'a. Fa-hadha mal, wa-jadhrhu thalatha.



\textbf{Mathal akhar:} khamsatu amwalin tusawi thamanin. Fa-mal wahid

yusawi khumus thamanin, wa-huwa sitta 'ashar.



\textbf{Mathal thalith:} niSf mal yusawi thamaniya 'ashar. Fa-l-mal

sitta wa-thalathun, wa-jadhrhu sitta.

\end{originaltext}



\begin{mathexplain}

\textbf{al-Tawdih al-Riyadi lil-HaSala al-Thaniya:}



\begin{center}

$\begin{aligned}

  x^2 &= 9 \quad\Rightarrow\quad x = 3 \\

  5x^2 &= 80 \quad\Rightarrow\quad x^2 = 16, \quad x = 4 \\

  \tfrac{1}{2}x^2 &= 18 \quad\Rightarrow\quad x^2 = 36, \quad x = 6

\end{aligned}$

\end{center}



\textbf{al-Tariqa al-'Amma:} fi al-mu'adala $ax^2 = c$, naqsimu 'ala

$a$ fa-nuHallu 'ala $x^2 = c/a$, thumma na'khudhu al-jadhr al-tarb'i

li-kila al-Tarfayn: $x = \sqrt{c/a}$.

\end{mathexplain}



\subsection{al-HaSala al-Thalitha: al-usuL tusawi 'adadan}



\begin{originaltext}

\textbf{Mathal:} jadhr wahid yusawi thalathata 'adadan. Fa-l-jadhr

thalatha, wa-maluha tis'a.



\textbf{Mathal akhar:} arba'atu usuLin tusawi 'ishrina. Fa-jadhr wahid

yusawi khamsa, wa-l-mal alladhi yunsha'u minhu khamsa wa-'ishrun.



\textbf{Mathal thalith:} niSf jadhr yusawi 'ashara. Fa-l-jadhr kulluh

yusawi 'ishrina, wa-l-mal alladhi yunsha'u minhu arba'umi'a.

\end{originaltext}



\begin{mathexplain}

\textbf{al-Tawdih al-Riyadi lil-HaSala al-Thalitha:}



\begin{center}

$\begin{aligned}

  x &= 3 \\

  4x &= 20 \quad\Rightarrow\quad x = 5, \quad x^2 = 25 \\

  \tfrac{1}{2}x &= 10 \quad\Rightarrow\quad x = 20, \quad x^2 = 400

\end{aligned}$

\end{center}



\textbf{al-Tariqa al-'Amma:} fi al-mu'adala $bx = c$, naqsimu 'ala

$b$ fa-nuHallu 'ala $x = c/b$. Wa-hadhihi abSar al-HaSalat, li-annaha

mu'adala khattiyya.

\end{mathexplain}



\newpage



%% ========================================================================

%% THE THREE COMPOUND CASES

%% ========================================================================

\section{al-HaSalat al-Murakkaba al-Thalath}



\begin{mathexplain}

\textbf{al-HaSalat al-Murakkaba (tatallaBu naSf al-juthur -- takmil

al-murabba'):}



\begin{center}

\begin{tabular}{|c|c|l|}

\hline

\textbf{al-HaSala} & \textbf{al-Mu'adala} & \textbf{al-Isem al-Arabi} \\

\hline

al-Rabi'a & $ax^2 + bx = c$ & mal wa-usuL yusawi 'adadan \\

\hline

al-Khamisa & $ax^2 + c = bx$ & mal wa-'adad yusawi usuLan \\

\hline

al-Sadisa & $bx + c = ax^2$ & usuL wa-'adad yusawi malan \\

\hline

\end{tabular}

\end{center}



\textbf{limadha tatallaBu hadhihi al-HaSalat naSf al-juthur?}



li-annaha taShtamilu 'ala thalathat aTrafi mal + juthur + 'adad,

wa-la yumkinu 'azlu al-majhul bi-l-qisma al-mubashshara. badalan

min dhalika, nastakhdimu tariqat \textbf{takmil al-murabba'}

allati taHawwilu al-mu'adala al-tarb'iyya ila mu'adala min al-shakl

$(x + p)^2 = q$ yumkinu Halluha bi-akhdhi al-jadhr al-tarb'i.

\end{mathexplain}



\subsection{al-HaSala al-Rabi'a: mal wa-usuL yusawi 'adadan}



\begin{originaltext}

\textbf{Mathal:} mal wa-'asharatu usuLihi tusawi tis'a wa-thalathin

darhaman.



al-Hall: naSifu 'adada al-juthuri, fa-yakunu khamsa, wa-naDribuhu fi

nafsihi fa-yakunu khamsa wa-'ishrina, fa-nuDifuhu ila tis'a

wa-thalathin fa-yakunu tis'a wa-sittin, fa-na'khudh jadhrhu fa-yakunu

thamaniya, fa-nanquSu minhu naSf 'adadi al-juthuri wa-huwa khamsa,

Yabqa thalatha, wa-hadha huwa jadhr al-mal al-maTlub, wa-l-mal

tis'a.

\end{originaltext}



\begin{mathexplain}

\textbf{al-Hall al-'Am lil-HaSala al-Rabi'a:}



\begin{center}

$\boxed{x^2 + bx = c \quad\Longrightarrow\quad x = \sqrt{\left(\frac{b}{2}\right)^2 + c} - \frac{b}{2}}$

\end{center}



\textbf{al-Tatbiq 'ala al-mathal:}

$x^2 + 10x = 39 \Rightarrow x = \sqrt{25 + 39} - 5 = \sqrt{64} - 5 = 8 - 5 = 3$



\textbf{al-Tafsir:} al-Khwarizmi la yaktubu Sighatan 'amma -- bal

yuqaddimu \textbf{Khwarizmiyya} (silsilat khuTwat) yumkinu tatbiQuha

'ala ay mu'adala min hadha al-shakl. Wa-hadha huwa al-sabab fi

anna naSammiyanfusana al-yawma ``Khwarizmiyyat'' (algorithms)!

\end{mathexplain}



\begin{commentary}

\textbf{limadha ya'malu hadha al-Hall?}



Tariqat al-Khwarizmi ta'tamidu 'ala fikrat \textbf{takmil al-murabba'}.

'indama nuDifu $(b/2)^2$ ila kila Tarfi al-mu'adala $x^2 + bx = c$,

naHsalu 'ala:

$x^2 + bx + (b/2)^2 = c + (b/2)^2$



al-Taraf al-aysar huwa al-an mirab' kamil: $(x + b/2)^2 = c + (b/2)^2$



'inda akhdhi al-jadhr al-tarb'i: $x + b/2 = \sqrt{c + (b/2)^2}$



Wa-bi-Tarqi $b/2$: $x = \sqrt{(b/2)^2 + c} - b/2$



Wa-hadha bi-ZabT ma yaf'aluhi al-Khwarizmi, wa-in duna istikhdam

al-rumuz al-jabriyya al-Haditha!

\end{commentary}



\newpage



\subsection{al-HaSala al-Khamisa: mal wa-'adad yusawi usuLan}



\begin{originaltext}

\textbf{Mathal:} mal wa-wahid wa-'ishruna 'adadan yusawi 'asharata

usuLihi.



al-Hall: naSifu 'adada al-juthuri fa-yakunu khamsa, wa-naDribuhu fi

nafsihi fa-yakunu khamsa wa-'ishrina, fa-nanquSu minhu al-wahida

wa-l-'ishrina al-mulHaqqayni bi-l-mal, Yabqa arba'a, fa-nastakhrij

jadhraha fa-yakunu ithnayn, fa-nanquSuhu min naSf 'adadi al-juthuri

wa-huwa khamsa, Yabqa thalatha, wa-hadha jadhr al-mal al-maTlub, wa-l-mal

tis'a. Aw tuDifu al-jadhr ila naSf al-juthuri, fa-yakunu sab'a,

wa-hadha jadhr al-mal al-maTlub, wa-l-mal tis'a wa-arba'un.

\end{originaltext}



\begin{mathexplain}

\textbf{al-Hall al-'Am lil-HaSala al-Khamisa:}



\begin{center}

$\boxed{x^2 + c = bx \quad\Longrightarrow\quad x = \frac{b}{2} \pm \sqrt{\left(\frac{b}{2}\right)^2 - c}}$

\end{center}



\textbf{al-Tatbiq:}

$x^2 + 21 = 10x \Rightarrow x = 5 \pm \sqrt{25 - 21} = 5 \pm 2 = 3 \text{ aw } 7$



\textbf{mulahadha muhimma:} hadhihi hiya al-HaSala \textbf{al-wahida}

allati yakunu fiha Hallan -- wa-li-dhalika fa-inna al-Khwarizmi yadhkuru

SariHan annahu yumkinu istikhdam al-jam' aw al-TarH. Amma fi

al-HaSalat al-ukhra fa-yakunu al-Hall wahidan faqaT.

\end{mathexplain}



\begin{commentary}

\textbf{shuruT imkaniyyat al-Hall fi al-HaSala al-Khamisa:}



ba'da naSf al-juthur wa-Darb al-naSf fi nafsihi, idha kana al-natij

\textbf{aqall} min al-'adad al-mulHaqq bi-l-mal, fa-l-masa'la

\textbf{mumtani'a} (la Hall Haqiqiyyan). Amma idha kana al-natij

\textbf{musawiyan} lil-'adad, fa-jadhr al-mal yusawi naSf al-juthur

watfahu, bi-la ziyada wa-la nuqsan.



\vspace{0.2cm}



\textbf{bi-l-rumuz al-Haditha:} al-dakhil taHt al-jadhr yajib an la

yakun saliban: $(b/2)^2 - c \geq 0$. Wa-hadha ya'ni $b^2 \geq 4c$,

wa-huwa al-shart al-ma'ruf al-yawma bi-ism \textbf{al-mumayyiz}

(discriminant) lil-mu'adala al-tarb'iyya.

\end{commentary}



\newpage



\subsection{al-HaSala al-Sadisa: usuL wa-'adad yusawi malan}



\begin{originaltext}

\textbf{Mathal:} thalathatu usuLin wa-arba'atu a'dadin yusawi malan.



al-Hall: naSifu al-juthuri fa-yakunu wahidan wa-niSf, wa-naDribuhu fi

nafsihi fa-yakunu ithnayn wa-rub'an, fa-nuDifuhu ila al-arba'a fa-yakunu

sitta wa-rub'an, fa-nastakhrij jadhraha fa-yakunu ithnayn wa-niSfan,

fa-nuDifuhu ila naSf al-juthuri wa-huwa wahid wa-niS fa-yakunu

arba'a, wa-hadha jadhr al-mal, wa-l-mal sitta 'ashar.

\end{originaltext}



\begin{mathexplain}

\textbf{al-Hall al-'Am lil-HaSala al-Sadisa:}



\begin{center}

$\boxed{bx + c = x^2 \quad\Longrightarrow\quad x = \frac{b}{2} + \sqrt{\left(\frac{b}{2}\right)^2 + c}}$

\end{center}



\textbf{al-Tatbiq:}

$3x + 4 = x^2 \Rightarrow x = 1\tfrac{1}{2} + \sqrt{2\tfrac{1}{4} + 4} = 1\tfrac{1}{2} + \sqrt{6\tfrac{1}{4}} = 1\tfrac{1}{2} + 2\tfrac{1}{2} = 4$



\textbf{mulahadha:} fi hadhihi al-HaSala yakunu al-Hall \textbf{mujaban

da'iman}, li-annana najma'u al-jadhr al-tarb'i ma'a niSf al-ma'amil.

Hadha yatawaafaqu ma'a Tabi'at al-masa'il al-'amaliyya allati kana

yata'amalu ma'aha al-Khwarizmi.

\end{mathexplain}



\newpage



%% ========================================================================

%% CHAPTER: GEOMETRICAL DEMONSTRATIONS

%% ========================================================================

\chapter{al-Tashrihat al-Hindisiyya lil-HaSalat al-Murakkaba}

\label{ch:geom}



\begin{originaltext}

Hadhi-hi hiya al-HaSalat al-sitt alladhi dhakartuha fi muqaddimati

hadha al-kitab. Qad wuDDoHat al-an. Qad bayyantu anna thalathan minha

la taHtaju ila naSf al-juthuri, wa-'allamtu kayfa yajibu Halluha. Wa-amma

al-thalath al-ukhra allati fiha naSf al-juthuri Daruriyyun, fa-ara min

al-muSlaH an uuwiHHahu bi-fuSulin munfasila, yu'ta fi kulli minha

shakl yubayyin sabab al-naSf.

\end{originaltext}



\begin{commentary}

\textbf{al-manhaj al-Hindisi lil-Khwarizmi:}



Min ahamma ma yumayyizu kitab al-Khwarizmi annahu la yaktaSi bi-taqdim

\textbf{Khwarizmiyya} (al-ijra' al-Hisabi) li-kulli HaSala, bal

yuqaddimu aydan \textbf{al-tashrih al-Hindisi} (geometric demonstration)

alladhi yuZhir \textit{limadha} ta'malu hadhihi al-tariqa. hadha

al-manhaj ya'ksi al-ta'thir al-Yunani -- khususan Iqlidus -- 'ala

al-riyadiyyat al-Islamiyya.



al-Khwarizmi yastakhdimu \textbf{murabba'at wa-ajniha (mustatTilat)}

li-tamthil al-kamiyyat al-jabriyya. al-mal huwa murabba', wa-l-juthur

hiya mustatTilat mulHaqa bi-jawanib al-murabba', wa-l-a'dad hiya

murabba'at ma'rufa al-misaHa. hadha al-tafsir al-Hindisi huwa alladhi

yubirru qawa'idat ``naSf al-juthur'' -- fa-hiya al-tariqa allati

nuHawwilu biha al-shakl ghayr al-munTaZZim ila murabba' kamil.

\end{commentary}



\section{Tashrih al-HaSala al-Rabi'a: ``Mal wa-'asharatu usuL yusawi tis'a wa-thalathin''}



\begin{originaltext}

al-shakl al-mubayyin li-hadha huwa \textbf{murabba'}, aDla'uhu

majhula. Yumanthilu al-mal, jadhr alladhi turidu ma'rifatahu. Hadha

huwa al-shakl A B, wa-kullu Dila' min aDla'ihi yumkinu i'tibaruhu

ahada juthurihi; wa-idha Darabta ahada hadhihi al-aDla' fi ayyi

'adadin, fa-inna miqdar dhalika al-'adadi yumkinu an yunZar ilayhi

ka-'adadi al-juthuri allati tuDafu ila al-mal.



kullu Dila' min aDla' al-murabba' yumanthilu jadhr al-mal; wa-kama

fi al-misal kanat al-juthuru murtabiTa bi-l-mal, yumkinuna an

na'khudha rub'a al-'ashara, ay thnan wa-niSf, wa-nulHiqaha bi-kulli

min al-aDla' al-arba'a lil-shakl. Fa-ma'a al-murabba' al-asli A B,

tulHaqu arba'at ajniha mutawaziyya al-aDla', kullu minha lahu Dila'

min aDla' al-murabba' ka-Tul, wa-'adadi thnayni wa-niSf ka-'arD;

wa-hiya al-ajniha J, Z, T, K.



Ladayna al-an murabba' mutasawi al-aDla' majhul, walakin fi kulli

min zawaya-hi al-arba' yanquSu laHZa murabba'a min thnayni wa-niSf

maDruba fi thnayni wa-niSf. Li-nu'awiDa 'an hadha al-nuqsan

wa-nukammila al-murabba', yajib an nuDifa (ila ma ladayna) arba'

marat murabba'a thnayni wa-niSf, ay khamsa wa-'ishrun.



Wa-nahnu na'lamu (min al-mas'ala) anna al-shakl al-awwal, ayyuha

al-murabba' al-mumaththil lil-mal, ma'a al-ajniha al-arba'a al-muHiTa

bihi, al-mumaththila lil-'asharati juthur, yusawi tis'a wa-thalathina

'adadan. Fa-idha aSafnaha ila khamsa wa-'ishrina, allati hiya

musawiyya lil-murabba'at al-arba' 'inda zawaya al-shakl A B, allati

tukammilu al-murabba' al-kabir D H, fa-n'alamu anna hadha al-majmui'

ya'kul sitta wa-arba'in. ahadu aDla' hadha al-murabba' al-kabir huwa

jadhrhu, ay thamaniya. Fa-idha naqasna Di'fa rub'i al-'ashara, ay

khamsa, min thamaniya, kama min Tarafayyi al-Dil' al-murabba' al-kabir

D H, fa-yabqa min hadha al-Dil' thalatha, wa-hadha huwa jadhr

al-mal, aw Dil' al-shakl al-asli A B.

\end{originaltext}



\begin{geombox}

\textbf{al-Burhan al-Hindisi lil-HaSala al-Rabi'a:}



\begin{center}

\begin{tabular}{lcl}

\textbf{al-shakl} & $=$ & \textbf{al-misaHa} \\

\hline

Murabba' A B & $=$ & $x^2$ \\

+ 4 mutawaziyat al-aDla' & $=$ & $10x$ \\

\hline

al-majmui' al-juz' & $=$ & $39$ \\

+ 4 murabba'at zawiya ($2.5 \times 2.5$) & $=$ & $25$ \\

\hline

al-majmui' (Murabba' D H al-kabir) & $=$ & $64$ \\

Dila' al-Murabba' al-kabir & $=$ & $\sqrt{64} = 8$ \\

$x = 8 - 2 \times 2.5 = 8 - 5 =$ & \fbox{3} & \\

\end{tabular}

\end{center}



\vspace{0.3cm}



\textbf{al-tafsir al-mar'i:} nabda'u bi-murabba' janibuhu $x$

(yumanthil al-mal). NulHiq bi-kulli janib mustatTila Tuluhu $x$

wa-'arDuha $2.5$ (rub' al-'adad 10). al-arba' mustatTilat tumanthil

$10x$. al-misaHa al-kulliyya hatta al-an hiya $x^2 + 10x = 39$.



Li-nukammila al-murabba' al-kabir, naHtaju li-izafat arba' murabba'at

Sagira 'inda al-zawaya, kullu minha $2.5 \times 2.5 = 6.25$. al-majmui'

$4 \times 6.25 = 25$. al-an al-murabba' al-kabir lahu misaHa

$39 + 25 = 64$, fa-Dil'uhu huwa $8$. NunquS minhu $5$ ('arDu

mustatTilatayn) fa-nuHallu 'ala $x = 3$.

\end{geombox}



\begin{commentary}

\textbf{limadha ``rub' '' thumma ``Di'fa rub' ''?}



Yulahidhu al-qari'u anna al-Khwarizmi ya'khudhu \textbf{rub'a} 'adad

al-juthuri (rub' al-'ishrin = 2.5) li-insha' al-ajniha al-arba'a, thumma

yuTraHu \textbf{Di'fa} hadha al-rub'i min Dil' al-murabba' al-kabir.

limadha?



li-anna kulli janib min jawaniB al-murabba' yaHaSSulu 'alayhi \textit{janiHan}

(wahid min kulli janib). Fa-idha kana 'arDu kulli janiH huwa $b/4$,

fa-inna al-majmui' al-muTraHu min kila Tarafayi al-Dil' huwa

$2 \times (b/4) = b/2$, wa-huwa bi-ZabT ``naSf 'adadi al-juthur''

al-madhkur fi al-Khwarizmiyya!



Wa-haka-dha yuZhiru al-burhan al-Hindisi \textit{sabab} naSf al-juthur --

wa-huwa taHawwilu al-shakl ila murabba' kamil.

\end{commentary}



\newpage



\section{Tashrih al-HaSala al-Khamisa: ``Mal wa-21 yusawi 10 usuL''}



\begin{originaltext}

Numaththilu al-mal bi-murabba' A D, Tul Dil'ihi majhul ladayna.

Ila hadha nulHiqu mustatTila, 'arDuhu musawin li-ahadi aDla'

al-murabba' A D, kayna al-Dil' H N. Hadha al-mustatTil huwa H B.

Tul al-shaklayni ma'an yusawi al-KhaTT H J. Na'lamu anna Tuluha

'asharatu a'dadin; fa-kullu murabba'in lahu aDla' mutasawiya

wa-zawaya mutasawiya, wa-ahad aDla'ihi maDruban fi wahida huwa

jadhr al-murabba', aw maDruban fi ithnatayn fa-huwa Di'fu jadhr

al-shakl al-batt.



Fakama wada'at al-mas'ala anna malan wa-wahidan wa-'ishrina 'adadan

yusawi 'asharata juthuran, yumkinuna an nastanTiZu anna Tul al-KhaTT

H J yusawi 'asharata a'dadin, idha al-KhaTT J D yumanthilu jadhr

al-mal. Naqsimu al-an al-KhaTT J H ila niSfayn mutasawiyayn 'inda

al-nuqTa W: fa-l-KhaTT W J idhun yusawi H W.



Wa-na'lamu anna al-ruba'i H B yumanthilu al-wahida wa-l-'ishrina

'adadan allati zidat 'ala al-murabba'. Wa-wujida anna al-murabba'

M T kulluhu yusawi khamsa wa-'ishrina. Fa-idha naqasna min hadha

al-murabba' M T al-ruba'iyyayn H T wa-M R, alladhayn yusawiyan

wahidan wa-'ishrina, fa-yabqa murabba' Sagir K R, yumanthilu

al-furq bayna khamsa wa-'ishrina wa-wahid wa-'ishrina. Wa-huwa

arba'a; wa-jadhruh, al-mumaththal bi-l-KhaTT W A alladhi yusawi

A J, huwa ithnan.

\end{originaltext}



\begin{geombox}

\textbf{al-Burhan al-Hindisi lil-HaSala al-Khamisa:}



\begin{center}

\begin{tabular}{lcl}

$MT = (10/2)^2$ & $=$ & $25$ \\

$HT + MR = HB$ & $=$ & $21$ \\

$KR = 25 - 21$ & $=$ & $4$ \\

$\sqrt{KR}$ & $=$ & $2$ \\

$x = WJ - WA = 5 - 2 =$ & \fbox{3} & \\

aw $x = WJ + WA = 5 + 2 =$ & \fbox{7} & \\

\end{tabular}

\end{center}



\vspace{0.3cm}



\textbf{al-tafsir:} Nabni murabba'an janibuhu $x$ (al-mal). NulHiq

bihi mustatTilan Tuluhu $10$ wa-'arDuhu $x$ -- walakin al-misaHa

al-za'ida $21$ tumaththalu bi-l-mustatTil alladhi yulSaqu kharij

al-murabba'. Nukammilu al-murabba' al-kabir bi-Tul $10$ (jami'

al-juthur), wa-na'khudhu niSfahu $5$. al-murabba' al-kabir ($5 \times 5 = 25$)

naqisan al-mustatTil ($21$) yu'tina murabba'an Sagiran ($4$) janibuhu

$2$. al-Hallani ya'tiyani min imkaniyyati TarHi aw jam'i hadha

al-janib min al-niSf.

\end{geombox}



\newpage



\section{Tashrih al-HaSala al-Sadisa: ``3 usuL wa-4 yusawi malan''}



\begin{originaltext}

Li-yakun al-mal mumaththan bi-ruba'i aDla'uhu majhula ladayna,

wa-in kanat mutasawiya fima baynaha, kama al-zawaya. Hadha huwa

al-murabba' A D, alladhi yashtamilu 'ala al-thalathati juthurin

wa-al-arba'ati a'dadin al-madhkura fi hadhihi al-mas'ala.



Fi kulli murabba'in ahad aDla'ihi, maDruban fi wahida, huwa

jadhr. NaqTu' al-an min al-murabba' A D al-ruba'i H D, wa-na'khudh

ahad aDla'ihi H J li-thalatha, wa-huwa 'adad al-juthur. Wa-huwa

musawin li-A R. Yatba'u idhan anna al-ruba'i H B yumanthilu

al-arba'a a'dadan allati zidat 'ala al-juthur. al-an nunaSSifu

al-Dil' J H alladhi yusawi thalathata juthur, 'inda al-nuqTa W;

min hadha al-qisma nabni al-murabba' W T alladhi huwa hasilu Darb

niSf al-juthur (wahid wa-niSf) fi nafsihi, ay thnan wa-rub'.

\end{originaltext}



\begin{geombox}

\textbf{al-Burhan al-Hindisi lil-HaSala al-Sadisa:}



\begin{center}

\begin{tabular}{lcl}

Murabba' $WT = (3/2)^2$ & $=$ & $2.25$ \\

$AN + KL = AR = 4$ & & (al-a'dad al-muDafa) \\

$GM = WT + (AN + KL) = 2.25 + 4$ & $=$ & $6.25$ \\

$AG = \sqrt{GM} = \sqrt{6.25}$ & $=$ & $2.5$ \\

$x = AG + GW = 2.5 + 1.5 =$ & \fbox{4} & \\

\end{tabular}

\end{center}



\vspace{0.3cm}



\textbf{al-tafsir:} Nabda'u bi-murabba' kabir ($x \times x = x^2$)

yumanthilu al-mal. NaqTa'u minhu mustatTila ($3 \times x = 3x$)

yumanthilu al-juthur. ma yabqa huwa al-mustatTil alladhi yumanthilu

al-'adad 4. Nukammilu al-murabba' al-Sagir 'inda al-zawiya bi-niSf

al-juthur (1.5). Na'khudhu al-jadhr al-tarb'i lil-majmui' wa-nujmi'

'alayhi niSf al-juthur.

\end{geombox}



\newpage



%% ========================================================================

%% CHAPTER: MULTIPLICATION

%% ========================================================================

\chapter{fi Darb al-Ta'abirat al-Jabriyya}

\label{ch:multiplication}



\begin{originaltext}

Sa-u'allimuka al-an kayfa taDribu al-a'dada al-majhula -- ay

al-juthur -- ba'Duha fi ba'D, in kanat munfarida, aw idha uDifat

ilayha a'dad, aw idha nuqiSat minha a'dad, aw idha nuqiyat min

a'dad; aydan kayfa tujmi'uha ba'Duha ila ba'D, aw kayfa tanquSu

ba'Duha min ba'D.



Kullama urida Darbu 'adadin fi 'adadin, fayajib an yukarrara

ahaduhuma bi-'adadi ma yaHtawihi al-akhar min al-waHadat.



Idha kanat a'dadun akbaru mu'allafa ma'a waHadatin li-tuDafa aw

tunqas minhum, fa-inna \textbf{arba'a Darbatin} lazima; ay: al-a'dad

al-akbaru fi al-a'dad al-akbar, wa-l-a'dad al-akbar fi al-waHadat,

wa-l-waHadat fi al-a'dad al-akbar, wa-l-waHadat fi al-waHadat.



Idha kanat al-waHadat al-mu'allafa ma'a al-a'dadi al-akbar

mujibatayni kulltayhima, fa-l-Darbatu al-rabi'a mujiba; wa-idha

kanat salibatayni kulltayhima, fa-l-Darbatu al-rabi'a mujiba aydan.

Wa-idha kanat ihdahuma mujiba wa-l-ukhra saliba, fa-l-Darbatu

al-rabi'a saliba.

\end{originaltext}



\begin{commentary}

\textbf{Qawa'id al-Darb al-Arba'i 'inda al-Khwarizmi:}



Yuqaddimu al-Khwarizmi huna qawa'id Darb dhati al-Haddayn -- ay ma

nusammi al-yawma ``tawzi' al-Darb 'ala al-jam' ''. Idda aradna Darba

$(a + x)(b + y)$, fa-nahnu naHtaju li-arba' Darbat munfasila:



\begin{enumerate}

  \item al-akbar fi al-akbar: $a \times b$

  \item al-akbar fi al-waHadat: $a \times y$

  \item al-waHadat fi al-akbar: $x \times b$

  \item al-waHadat fi al-waHadat: $x \times y$

\end{enumerate}



Wa-hadhihi hiya bi-ZabT al-Sigha allati na'rifuha al-yawma:

$(a + x)(b + y) = ab + ay + bx + xy$



\textbf{Ahammiyyat qawa'id al-Isharat:} Yulahidu al-Khwarizmi anna

``al-salibi fi al-salibi = mujib'' wa-``al-salibi fi al-mujibi

= salib''. Wa-hadhihi awwalu Siyaghat ma'rufa li-\textbf{qawa'id

al-isharat} fi tarikh al-riyadiyyat!

\end{commentary}



\begin{mathexplain}

\textbf{Qawa'id al-Isharat fi al-Darb:}



\begin{center}

\begin{tabular}{|c|c|c|}

\hline

& \textbf{mujib} & \textbf{salib} \\

\hline

\textbf{mujib} & mujib & salib \\

\hline

\textbf{salib} & salib & mujib \\

\hline

\end{tabular}

\end{center}



\vspace{0.3cm}



$\begin{aligned}

  (+x)(+y) &= +xy \\

  (-x)(-y) &= +xy \\

  (+x)(-y) &= -xy \\

  (-x)(+y) &= -xy

\end{aligned}$



hadhihi al-qawa'id kanat thawriyya fi zamaniha -- fa-qad samahat

bi-l-ta'amul ma'a ``al-kamiyyat al-naqiSa'' (al-a'dad al-saliba)

fi siyaq al-Darb wa-l-qisma.

\end{mathexplain}



\section{Amthila 'ala al-Darb}



\begin{originaltext}

\textbf{al-mithal al-awwal:} 'ashara wa-waHid yuDribani fi 'ashara

wa-ithnayn.



'asharat 'asharat mi'a; waHid 'ashara 'asharatu 'asharat mujiba;

ithnan 'ashara 'ishruna mujiba; wa-waHid fi ithnayn ithnan mujiban;

al-majmui' mi'a wa-ithnan wa-thalathun.



\textbf{al-mithal al-thani:} 'ashara naqis wahid, yuDrabu fi 'ashara

naqis wahid.



Fa-'asharat 'asharat mi'a; al-salib wahid fi 'ashara 'asharatu

'asharat saliba; al-salib al-akhar fi 'ashara 'asharatu 'asharat

salibat ukhra, fa-yakunu thamanun; lakinna al-salib wahid fi

al-salib wahid wahid mujib, fa-yakunu al-natij wahidan wa-thamanin.



\textbf{al-mithal al-thalith:} 'ashara wa-shay' yuDribani fi 'ashara

wa-shay'.



Fa-'asharat 'asharat mi'a, wa-'asharat shay' 'asharat ashya'; wa

'asharat shay' 'asharat ashya'; wa-l-shay' maDruban fi al-shay'

mal mujib, fa-yakunu al-natij mi'at dirham wa-'ishruna shay'an

wa-malan mujiban wahidan.



\textbf{al-mithal al-rabi':} 'ashara naqis shay', fi 'ashara naqis

shay'.



Fa-'asharat 'asharat mi'a; wa-salib shay' fi 'ashara 'asharatu

ashya' saliba; wa-salib shay' akhar fi 'ashara 'asharatu ashya'

saliba. Lakin salib shay' maDruban fi salib shay' mal mujib.

Fa-l-natij mi'a wa-mal, naqis 'ishruna shay'an.



\textbf{al-mithal al-khamis:} 'ashara naqis shay', fi 'ashara wa-

ziyadat shay'.



Fa-taqul: 'asharat 'asharat mi'a; wa-salib shay' fi 'ashara

'asharat ashya' saliba; wa-shay' fi 'ashara 'asharat ashya'

mujiba; wa-salib shay' fi shay' mal mujib; fa-l-natij mi'at

dirham, naqis malin.

\end{originaltext}



\begin{mathexplain}

\textbf{al-Tawdih al-Riyadi li-l-amthila:}



\begin{center}

$\begin{aligned}

  (10 + 1)(10 + 2) &= 100 + 10 + 20 + 2 = 132 \\

  (10 - 1)(10 - 1) &= 100 - 10 - 10 + 1 = 81 \\

  (10 + x)(10 + x) &= 100 + 10x + 10x + x^2 = 100 + 20x + x^2 \\

  (10 - x)(10 - x) &= 100 - 10x - 10x + x^2 = 100 - 20x + x^2 \\

  (10 - x)(10 + x) &= 100 - 10x + 10x - x^2 = 100 - x^2

\end{aligned}$

\end{center}



\textbf{mulahadha:} al-mithal al-khamis yuZhiru natijatan mudhishatan:

$(10 - x)(10 + x) = 100 - x^2$ -- wa-huwa \textbf{farq al-mirabba'ayn}

(difference of squares), ihdha ahamm al-Sughar al-jabriyya fi

jami' al-riyadiyyat!

\end{mathexplain}



\newpage



%% ========================================================================

%% CHAPTER: ADDITION AND SUBTRACTION

%% ========================================================================

\chapter{fi al-Jam' wa-al-TarH}

\label{ch:addition}



\section{'Amaliyyat 'ala al-Juthur al-tarb'iyya}



\begin{originaltext}

i'lam anna jadhr mi'atayni naqis 'ashara, mudafan ila 'ishrina

naqis jadhr mi'atayni, huwa 'asharat.



Wa-jadhr mi'atayni naqis 'ashara, munqasan min 'ishrin naqis jadhr

mi'atayni, huwa thalathun naqis Di'fi jadhr mi'atayni; wa-Di'fu

jadhr mi'atayni yusawi jadhr thamanimi'a.

\end{originaltext}



\begin{mathexplain}

\textbf{al-Tawdih al-Riyadi:}



\begin{center}

$\begin{aligned}

  (\sqrt{200} - 10) + (20 - \sqrt{200}) &= 10 \\

  (20 - \sqrt{200}) - (\sqrt{200} - 10) &= 30 - 2\sqrt{200} = 30 - \sqrt{800}

\end{aligned}$

\end{center}



\textbf{al-mulahadha al-muhimma:} al-Khwarizmi ya'rifu anna

$2\sqrt{200} = \sqrt{800}$. Wa-hadha yatba'u min al-qawa'id al-'amma

$a\sqrt{b} = \sqrt{a^2 b}$. Fa-l-jadhr al-tarb'i yu'ammuqu kakyan

jabri mustaqill yumkinu jam'uhu wa-TarQuhu mitla al-mutaghayyirat.

\end{mathexplain}



\section{fi muDa'afat al-Juthur wa-niSfiha}



\begin{originaltext}

Idha aradta an tuDa'ifa jadhr ayi malin ma'lum aw majhul,

(bi-ma'na annaka taDribuhu fi ithnayn) fa-yakfi annaka taDriba

ithnayn fi ithnayn, thumma fi al-mal; fa-jadhr al-natij yusawi

Di'fa jadhr al-mal al-asli.



Wa-idha aradta an ta'khudhah thalathan, fa-Draba thalathata fi

thalatha thumma fi al-mal; fa-jadhr al-natij huwa thalathatu

a'da' jadhr al-mal al-asli.



Wa-idha aradta an tajida niSf jadhr al-mal, fama 'alayka illa an

naDriba niSf fi niSf, wa-huwa rub'; thumma hadha fi al-mal:

fa-jadhr al-natij sa-yakunu niSfa jadhr al-mal al-awwal.

\end{originaltext}



\begin{mathexplain}

\textbf{al-Tawdih al-Riyadi:}



\begin{center}

$\begin{aligned}

  2\sqrt{S} &= \sqrt{4S} \\

  3\sqrt{S} &= \sqrt{9S} \\

  \tfrac{1}{2}\sqrt{S} &= \sqrt{\tfrac{1}{4}S}

\end{aligned}$

\end{center}



\textbf{al-tafsir:} al-qawa'id hiY $\sqrt{a^2 \cdot S} = a\sqrt{S}$.

al-Khwarizmi yuqaddimuha ka-Khwarizmiyya 'amaliyya: li-tuDa'ifa

al-jadhr, aDrib al-mal (mal) fi 4 (2x2) thumma khudh al-jadhr.

hadhihi al-tariqa tatajannabu al-ta'amul ma'a ``naSf al-jadhr''

mubashsharatan.

\end{mathexplain}



\newpage



%% ========================================================================

%% CHAPTER: DIVISION

%% ========================================================================

\chapter{fi al-Qisma}

\label{ch:division}



\begin{originaltext}

Idha aradta an taqsima jadhr tis'a 'ala jadhr arba'a, fa-bda' bi-

qismati tis'a 'ala arba'a, fa-yu'tika ithnayn wa-rub'an; fa-jadhr

hadha huwa al-'adad alladhi taHtaju ilayhi -- wa-huwa wahid wa-niSf.



Wa-idha aradta an taqsima jadhr arba'a 'ala jadhr tis'a, fa-aqsim

arba'a 'ala tis'a; fa-huwa arba'u atsa' al-wahda; wa-jadhr hadha

ithnan maqsuman 'ala thalatha; ay thulthay al-wahda.



Wa-idha aradta an taDruba jadhr tis'a fi jadhr arba'a, fa-Drab

tis'a fi arba'a; fa-yu'tika sitta wa-thalathin; fa-khudh jadhr,

fa-huwa sitta; wa-hadha huwa jadhr tis'a maDruban fi jadhr arba'a.



Wa-idha aradta an taDruba Di'fa jadhr tis'a fi thalathati a'da'

jadhr arba'a, fa-khudh Di'fa jadhr tis'a Hasaba al-qawida al-

madhkura, Hatta ta'rifa jadhr ayi mal huwa. Waf'il al-mithla

bi-l-nisbati li-l-thalathati juthur al-arba'a li-ta'riffa ma yajib

an yakuna mal dhalika al-jadhr. Thumma aDrab hadhayni al-malayn

ahadahuma fi al-akhar, fa-jadhr al-natij yusawi Di'fa jadhr tis'a

maDruban fi thalathati a'da' jadhr arba'a.

\end{originaltext}



\begin{mathexplain}

\textbf{Qawa'id al-Juthur 'inda al-Khwarizmi:}



\begin{center}

$\begin{aligned}

  \frac{\sqrt{9}}{\sqrt{4}} &= \sqrt{\frac{9}{4}} = \frac{3}{2} = 1.5 \\

  \frac{\sqrt{4}}{\sqrt{9}} &= \sqrt{\frac{4}{9}} = \frac{2}{3} \\

  \sqrt{9} \times \sqrt{4} &= \sqrt{9 \times 4} = \sqrt{36} = 6 \\

  (2\sqrt{9}) \times (3\sqrt{4}) &= \sqrt{36 \times 36} = 36

\end{aligned}$

\end{center}



\textbf{al-qawa'id al-'Amma:}

\begin{itemize}

  \item $\dfrac{\sqrt{a}}{\sqrt{b}} = \sqrt{\dfrac{a}{b}}$

  \item $\sqrt{a} \times \sqrt{b} = \sqrt{a \times b}$

  \item $(m\sqrt{a}) \times (n\sqrt{b}) = mn\sqrt{ab}$

\end{itemize}



Hadhihi al-qawa'id hiya al-asas alladhi yabni 'alayhi al-Khwarizmi

jami' 'amaliyyatihi 'ala al-juthur.

\end{mathexplain}



\newpage



%% ========================================================================

%% CHAPTER: THE SIX ILLUSTRATIVE PROBLEMS

%% ========================================================================

\chapter{al-sittu Masa'il al-TawDiHiyya}

\label{ch:six-problems}



\begin{originaltext}

Qabla fuSuli al-Hisab wa-anwa'iha al-mukhtalifa, sa'uqaddimu al-an

sitta masa'il, kamithal 'ala al-HaSalat al-sitt allati tanaWalat

ha fi bidayat hadha al-'amal. Qad bayyantu anna thalathan min

hadhihi al-HaSalat, li-Halliha, la taHtaju ila naSf al-juthur, wa-

dhakartu aydanna anna al-Hisab bi-l-jabr wa-l-muqabala yajib ann

yaqulaka da'iman ila ihdha haa al-HaSalat. Sa-ulHaqu hadhihi al-

masa'il al-an, allati sa-takunu aqraba ila al-fahm, wa-tusahhilu

istiya' al-mawdii', wa-tubayyinu al-ijj b-iDaaa akbar.

\end{originaltext}



\begin{commentary}

\textbf{limadha sittu masa'il?}



Yuqaddimu al-Khwarizmi ba'da shariH al-HaSalat al-sitt \textbf{sitta

masa'il tatbiqiyya} -- mas'ala wahida li-kulli HaSala. Hadhihi al-

masa'il laysat majarrad tamariyn riyadiyya, bal hiya \textit{qiSas}

(word problems) ta'ksi mawaqif waqi'iyya. Wa-al-ahammu min dhalika

annaha tuZhiru kayfiyyat \textbf{nimuDajat} mas'ala waqi'iyya ila

mu'adala jabriyya -- wa-hiya al-mahara al-asasiyya fi 'ilm al-jabr.

\end{commentary}



\section*{al-mas'ala al-ula}

\addcontentsline{toc}{section}{al-mas'ala al-ula}



\begin{originaltext}

``Qasamtu 'ashara ila qismayn; Darabtu ahadahuma fi al-akhar; ba'da

hadha Darabtu ahada al-qismayn fi nafsihi, wa-hasil al-Darb fi

nafsihi arba'atu a'da'i ma hasala min Darbi ahadi al-qismayni fi

al-akhar.''



al-Hall: li-yakun ahad al-qismayn shay'an, wa-al-akhar 'ashara

naqis shay'. TaDrabu al-shay' fi 'ashara naqis shay'; fa-huwa

'asharat ashya' naqis mal. Thumma taDribuhu fi arba'a, li-anna al-

mas'ala taqulu ``arba'atu a'da' ''. Fa-l-natij arba'un 'asharat

ashya' naqis arba'ati amwalin. Ba'da hadha taDribu al-shay' fi al-

shay', ay ahada al-qismayni fi nafsihi. Wa-hadha mal, huwa musawin

li-arba'in shay'an naqis arba'ati amwalin. KhaffiDha al-an bi-l-

arba'ati amwalin, wa-azidha ila al-mali al-waHid. Fa-l-mu'adala:

arbau 'ishruna shay'an tusawi khamsata amwla; wa-malu wahid

yusawi thamaniyata juthuran, ay arba'a wa-sittin; wa-jadhr hadha

thamaniya.

\end{originaltext}



\begin{mathexplain}

\textbf{al-Hall al-Riyadi:}



\begin{center}

$\begin{aligned}

  x^2 &= 4x(10-x) = 40x - 4x^2 \\

  \Rightarrow\quad 5x^2 &= 40x \quad\Rightarrow\quad x^2 = 8x = 64 \quad\Rightarrow\quad x = 8

\end{aligned}$

\end{center}



al-qismaan huma: $x = 8$ wa $10 - x = 2$.\\

Hadhihi al-mas'ala taquluna ila ihdha al-HaSalat al-sitt, ay:

\textbf{mal yusawi usuLan}.

\end{mathexplain}



\section*{al-mas'ala al-thaniya}

\addcontentsline{toc}{section}{al-mas'ala al-thaniya}



\begin{mathexplain}

\textbf{al-mas'ala:} ``Qasamtu 'ashara ila qismayn; Darabtu kulla

qism fi nafsihi, wa-ba'dahu 'ashara fi nafsaha: hasilu 'asharat

fi nafsiha yusawi ahada al-qismayni maDruban fi nafsihi, wa-ba'dahu

fi ithnayn wa-sab'at atsa'.''



$100 = x^2 \times 2\tfrac{7}{9} = x^2 \times \tfrac{25}{9}$

$\Rightarrow x^2 = \tfrac{900}{25} = 36 \Rightarrow x = 6$



al-qismaan huma: $x = 6$ wa $10 - x = 4$.

\end{mathexplain}



\section*{al-mas'ala al-thalitha}

\addcontentsline{toc}{section}{al-mas'ala al-thalitha}



\begin{mathexplain}

\textbf{al-mas'ala:} ``Qasamtu 'ashara ila qismayn. Qasamtu al-

wahida 'ala al-ukhar, wa-al-natij arba'a.''



$\dfrac{10 - x}{x} = 4 \Rightarrow 10 - x = 4x \Rightarrow 10 = 5x \Rightarrow x = 2$



al-qismaan huma: $x = 2$ wa $10 - x = 8$.

\end{mathexplain}



\section*{al-mas'ala al-rabi'a}

\addcontentsline{toc}{section}{al-mas'ala al-rabi'a}



\begin{mathexplain}

\textbf{al-mas'ala:} ``Darabtu thulthi shay' wa-dirham fi rub'i

shay' wa-dirham, wa-al-natij 'ishrun.''



$\left(\dfrac{x}{3} + 1\right)\left(\dfrac{x}{4} + 1\right) = 20$

$\Rightarrow \dfrac{x^2}{12} + \dfrac{7x}{12} + 1 = 20$

$\Rightarrow x^2 + 7x = 228$



naSf al-juthur: $3.5$. al-murabba': $12.25$.

nuDiI ila 228: $240.25$. al-jadhr: $15.5$.

nanquS $3.5$: $x = 12$.

\end{mathexplain}



\section*{al-mas'ala al-khamisa}

\addcontentsline{toc}{section}{al-mas'ala al-khamisa}



\begin{mathexplain}

\textbf{al-mas'ala:} ``Qasamtu 'ashara ila qismayn; ba'da Darbi

kulli qism fi nafsihi, jamatuhuma; wa-al-majmu' thamaniya wa-

khamsun dirhama.''



$x^2 + (10-x)^2 = 58 \Rightarrow 2x^2 - 20x + 100 = 58$

$\Rightarrow x^2 + 28 = 11x$



$x = \dfrac{11}{2} \pm \sqrt{\dfrac{121}{4} - 28} = 5.5 \pm 2.5$.

al-qismaan: $x = 3$ aw $x = 7$.

\end{mathexplain}



\section*{al-mas'ala al-sadisa}

\addcontentsline{toc}{section}{al-mas'ala al-sadisa}



\begin{mathexplain}

\textbf{al-mas'ala:} ``Darabtu thultha jadhr fi rub'a jadhr, wa-al-

natij yusawi al-jadhr wa-arba'a wa-'ishrina dirhama.''



$\dfrac{x}{3} \cdot \dfrac{x}{4} = x + 24 \Rightarrow \dfrac{x^2}{12} = x + 24$

$\Rightarrow x^2 = 12x + 288$



$x = 6 + \sqrt{36 + 288} = 6 + 18 = 24$

\end{mathexplain}



\newpage



%% ========================================================================

%% CHAPTER: VARIOUS QUESTIONS

%% ========================================================================

\chapter{masa'il mukhtalifa}

\label{ch:various}



\begin{originaltext}

Wa-hadhihi masa'il iDafiyya li-tawsi' al-fahm wa-ta'kidi al-mabadi'

allati buniyat 'alayha al-HaSalat al-sitt al-kanuniyya.

\end{originaltext}



\section*{masa'il tawDiHiyya iDafiyya}

\addcontentsline{toc}{section}{masa'il tawDiHiyya iDafiyya}



\begin{originaltext}

\textbf{mas'ala 1:} qasamtu 'ashara ila qismayn, wa-Darabtu

ahadahuma fi al-akhar, fa-kana al-natij wahidan wa-'ishrina.



\textbf{mas'ala 2:} qasamtu 'ashara ila qismayn; ba'da Darbi kulla

qism fi nafsihi, naqastu al-asghar min al-akbar, wa-al-baqiyy

arba'un.



\textbf{mas'ala 3:} lana murabba', thulthahu min khumushihi yusawi

sub'i jadhr.



\textbf{mas'ala 4:} wujad jadhr murabba', idha Duriba fi arba'a

amthalihi kana al-majmui' 'ishrina.



\textbf{mas'ala 5:} murabba', idha Duribat arba' min juthur fi

khamsat min juthur antajat Di'f al-mal, ma'a fa'iD sitta wa-

thalathin dirhaman.



\textbf{mas'ala 6:} wujad murabba', idha uDifa thulthuhu ila

thalathati darahim, wa-nuqiSa al-majmu' min al-murabba', thumma

Duriba al-baqiyy fi nafsihi, astarada al-murabba'.

\end{originaltext}



\begin{mathexplain}

\textbf{Hullu-l-masa'il al-mukhtalifa:}



\textbf{mas'ala 1:} $x(10-x) = 21 \Rightarrow 10x - x^2 = 21 \Rightarrow x^2 + 21 = 10x$



$x = 5 \pm \sqrt{25 - 21} = 5 \pm 2 = 3 \text{ aw } 7$



\vspace{0.2cm}



\textbf{mas'ala 2:} $(10-x)^2 - x^2 = 40 \Rightarrow 100 - 20x = 40 \Rightarrow x = 3$



\vspace{0.2cm}



\textbf{mas'ala 3:} $\tfrac{2}{3} \cdot \tfrac{1}{5} x^2 = \tfrac{1}{7}x \Rightarrow \tfrac{2}{15}x^2 = \tfrac{1}{7}x \Rightarrow x = \tfrac{15}{14} = 1\tfrac{1}{14}$



\vspace{0.2cm}



\textbf{mas'ala 4:} $x \cdot 4x = 20 \Rightarrow 4x^2 = 20 \Rightarrow x^2 = 5 \Rightarrow x = \sqrt{5}$



\vspace{0.2cm}



\textbf{mas'ala 5:} $4x \cdot 5x = 2x^2 + 36 \Rightarrow 20x^2 = 2x^2 + 36 \Rightarrow 18x^2 = 36 \Rightarrow x^2 = 2$



\vspace{0.2cm}



\textbf{mas'ala 6:} $\left(x - \tfrac{x}{3} - 3\right)^2 = x \Rightarrow \left(\tfrac{2x}{3} - 3\right)^2 = x$



$\Rightarrow \tfrac{4x^2}{9} + 9 - 4x = x \Rightarrow x^2 + 20\tfrac{1}{4} = 11\tfrac{1}{4}x \Rightarrow x = 9 \text{ aw } 2\tfrac{1}{4}$

\end{mathexplain}



\newpage



%% ========================================================================

%% PART: PRACTICAL APPLICATIONS

%% ========================================================================

\part{al-Tatbiqat al-'amaliyya}

\label{part:practical}



\begin{historicalbox}

\textbf{al-jabr wa-al-Hayat al-yawmiyya:}



Min ahamma ma yumayyizu kitab al-Khwarizmi anna al-jabra 'indahu

laysa mujarrad lu'bat ithniyya aw niZariyya mujarrada, bal huwa

\textbf{adatun 'amaliyya} li-l-ta'amul ma'a al-masa'il al-yawmiyya.

YatanaWal al-juz' al-thani min kitabihi tatbiqat 'amaliyya fi

thalathat majalat ra'iyya:



\begin{enumerate}

  \item \textbf{al-mu'amalat al-tijariyya} -- al-bay' wa-l-shira'

    wa-l-Sarf wa-l-ijara

  \item \textbf{al-misaHa wa-l-handasa} -- misaHat al-ardiyn wa-Hufri

    al-qanawat

  \item \textbf{al-mirath wa-l-wasaya} -- ahkam al-fara'iD al-islamiyya

\end{enumerate}



Wa-hadha ya'ksi al-wiDifa \textbf{al-'amaliyya} li-l-'ilmi fi

al-HaDara al-islamiyya -- lam yakun al-'ilmu munfasilan 'an al-Hayat,

bal kana fi khidmati al-insan mubashsharatan.

\end{historicalbox}



\newpage



%% ========================================================================

%% CHAPTER: MERCANTILE TRANSACTIONS

%% ========================================================================

\chapter{fi al-mu'amalat al-tijariyya}

\label{ch:mercantile}



\begin{originaltext}

fi al-mu'amalat al-tijariyya: i'lam anna jami'a mu'amalat al-nas

min bay' wa-shira' wa-muSarafa wa-ijara, fa-innaha la takhlu min

qadrayni wa-arba'ati a'dadin yas'alu 'anha al-sa'il, wa-hiya al-

miqdar wa-l-thaman wa-l-kayfiyya wa-l-jami'. 'adad al-miqdar

yunasibu 'adad al-jami', wa-'adad al-thaman yunasibu 'adad al-

kayfiyya. thalathatu min hadhihi al-arba'a da'iman ma'luma, wa-

Wahid majhul, wa-huwa alladhi yatadammanuhu qawluhu ``kam?'' wa-

huwa al-maTlub.

\end{originaltext}



\begin{commentary}

\textbf{Qa'idat al-arba'a a'dad -- al-jabr al-tijari:}



Yasifu al-Khwarizmi huna ma nusammi al-yawma ``qawa'id al-thalath''

(Rule of Three) aw ``al-tanasiq''. Kullu mas'alat tijariyya tatadammanu

arba'a kamiyyat murtabiTa bi-l-tanasiq:



\begin{itemize}

  \item \textbf{al-miqdar} (al-miqdar): kamiyyat al-sila' fi al-Safqa al-ula

  \item \textbf{al-thaman} (al-thaman): si'r tilka al-kamiyya

  \item \textbf{al-kayfiyya} (al-kamiyya): kamiyyat al-sila' fi al-Safqa al-thaniya

  \item \textbf{al-jami'} (al-majmui'): al-si'r al-ijmali al-maTlub

\end{itemize}



thalathu min hadhihi al-kamiyyat ma'lumma wa-l-rabi'a majhula.

Wa-l-'alaqa hiya: $\text{al-miqdar} : \text{al-kayfiyya} :: \text{al-jami'} : \text{al-thaman}$

aw bi-l-mu'adala: $a \times b = c \times d$

\end{commentary}



\begin{mathexplain}

\textbf{Qa'idat al-arba'a a'dad:}



\begin{center}

\begin{tabular}{c|c}

\textbf{al-miqdar} & \textbf{al-thaman} \\

\hline

$a$ (ma'lum) & $b$ (ma'lum) \\

$c$ (ma'lum aw majhul) & $d$ (ma'lum aw majhul) \\

\end{tabular}



\vspace{0.3cm}



$\text{al-majhul} = \dfrac{\text{hasil Darb al-'adadayni al-mutanasibayn}}{\text{al-'adad al-ma'lum al-thalith}}$

\end{center}



\textbf{al-Amthila:}

\begin{enumerate}

  \item ``'ashara bi-sitta, kam bi-arba'a?''

  $\text{al-kamiyya} = \dfrac{10 \times 4}{6} = \dfrac{40}{6} = 6\tfrac{2}{3}$



  \item ``'ashara bi-thamaniya, ma yakunu al-jami' li-arba'a?''

  $\text{al-jami'} = \dfrac{8 \times 4}{10} = \dfrac{32}{10} = 3\tfrac{1}{5}$



  \item ``lil-'amili ujrat 'asharati darahim fi al-shahr; kam

    yajib an takunu ujratuhu fi sittati ayyaam?''

  $\text{al-ujra} = \dfrac{10 \times 6}{30} = 2 \text{ dirahm}$,

  ba-frD anna al-shahr thalathun yawman.

\end{enumerate}

\end{mathexplain}



\newpage



%% ========================================================================

%% CHAPTER: MENSURATION

%% ========================================================================

\chapter{al-misaHa wa-l-handasa}

\label{ch:mensuration}



\begin{originaltext}

i'lam anna ma'na qawlina ``wahid fi wahid'' huwa al-misaHa: ay

Dhiraa' fi Tuliha fi Dhiraa' fi 'arDiha.



kullu murabba' mutasawi al-aDla' wa-l-zawaya, lahu Dhiraa' wahida

li-kulli Dil', fa-lahu wahid fi misaHihi. Fa-in kana li-mithali

hadha al-murabba' Di'lan fi Dil'ihi, fa-inna misaHat al-murabba'

'arba'atu amthali misaHa murabba' Dil'uhu Dhiraa' wahida.

\end{originaltext}



\begin{commentary}

\textbf{manhaj al-Khwarizmi fi al-misaHa:}



Yabda'u al-Khwarizmi bi-ta'rifi al-misaHa 'ala annaha murabba' al-

wahda (Dhiraa' fi Dhiraa'). hadha al-ta'rif yaj'alu al-misaHa

\textbf{kamiyya mujiba da'iman}, wa-huwa ma yufassiru iqtiSarahu

'ala al-ma'amilat al-mujiba fi al-jabr.



al-mulahadha al-muhimma anna al-Khwarizmi yafhamu al-misaHa

\textit{bi-shakl hindisi} -- fa-l-misaHa laysat majarrad 'adad, bal

hiya murabba' lahu Tul wa-'arD. hadha al-tafakkur al-hindisi huwa

alladhi yumakkinuhu min i'itai al-burahin al-hindisiyya li-l-HaSalat

al-jabriyya.

\end{commentary}



\section{al-muthallathat}



\begin{originaltext}

Idha Darabta irtifa'a ayi muthallath mutasawi al-aDla' fi niSfi

al-qaa alladhi yaqumu 'alayha al-khaTT al-mamayyiz lil-irtifaa'

'amudiyyan, fa-l-natij yu'ti misaHa dhalika al-muthallath.

\end{originaltext}



\begin{mathexplain}

\textbf{misaHat al-muthallath:}



\begin{center}

$\text{misaHat al-muthallath} = \text{al-irtifaa'} \times \dfrac{\text{al-qa}}{2}$

\end{center}



hadhihi al-Sigha SaHiHa li-\textbf{kulli} muthallath, laysa faqaT

al-muthallath al-mutasawi al-aDla'. wa-l-Khwarizmi yastakhdimuha

fi masa'ili misaHa al-ardiyn.

\end{mathexplain}



\section{al-dawa'ir}



\begin{originaltext}

Fi kulli da'ira, hasilu Darbi qiTriha fi thalath wa-sabi'a wahida

yusawi al-muHit.



Wa-tujad misaHa ayi da'ira bi-Darbi niSfi al-muHit fi niSfi al-qiT.



Wa-idha Darabta qiTa ayi da'ira fi nafsihi, wa-naqasta min al-natij

sub'an wa-niSfa sub'in minhu, fa-l-baqiy yusawi misaHat al-da'ira.

\end{originaltext}



\begin{mathexplain}

\textbf{misaHa wa-muHit al-da'ira 'inda al-Khwarizmi:}



\begin{center}

$\begin{aligned}

  C &= \pi d \approx \dfrac{22}{7} \times d \quad\text{(al-muHit)} \\

  A &= \dfrac{C}{2} \times \dfrac{d}{2} = \dfrac{Cd}{4} \quad\text{(al-misaHa -- al-tariqa al-ula)} \\

  A &= d^2 - \dfrac{d^2}{7} - \dfrac{d^2}{14} = d^2 \times \left(1 - \dfrac{3}{14}\right) = d^2 \times \dfrac{11}{14} \quad\text{(al-misaHa -- al-tariqa al-thaniya)}

\end{aligned}$

\end{center}



\textbf{mulahadha:} qimmat $\pi \approx \dfrac{22}{7}$ hiya al-taqrib

al-shahir alladhi istakhdamahu al-riyadiyyun al-islamiyyun wa-l-

Yunaniyyun. Amma al-tariqa al-thaniya fa-yuHassibu fiha al-Khwarizmi

$\pi/4 \approx 11/14$, wa-huwa taqrib mumtaz.

\end{mathexplain}



\section{al-aHjam}



\begin{originaltext}

Hajam al-jism ruba'i al-aDla' yujad bi-Darbi al-Tul fi al-'arD,

thumma fi al-irtifaa'.



Wa-amma al-makhruTat wa-al-ahram, kayna al-muthallath aw al-ruba'i,

fa-tuHassabu bi-Darbi thuluth misaHa al-qa fi al-irtifaa'.

\end{originaltext}



\begin{mathexplain}

\textbf{Hisab al-aHjam:}



\begin{center}

$\begin{aligned}

  V_{\text{muwazi al-mustatiilat}} &= \text{al-Tul} \times \text{al-'arD} \times \text{al-irtifaa'} \\

  V_{\text{haram}} &= \dfrac{1}{3} \times \text{(misaHa al-qa)} \times \text{al-irtifaa'}

\end{aligned}$

\end{center}



hadhihi al-Sighar SaHiHa wa-la tazalu tudarRasu fi al-madaris Hatta

al-yawm!

\end{mathexplain}



\section{al-muthallath al-qaaim al-zawiya (mubarhanat fifaghorus)}



\begin{originaltext}

laHiz annahu fi kulli muthallath qaaim al-zawiya al-Dil'ayn al-

qSirayn kulla wahid minhuma maDruban fi nafsihi, wa-majmu' al-

hasilayn, yusawi hasila Darbi al-Dil' al-TawiI maDruban fi nafsihi.

\end{originaltext}



\begin{geombox}

\textbf{mubarhanat fifaghorus 'inda al-Khwarizmi:}



\begin{center}

$\boxed{a^2 + b^2 = c^2}$

\end{center}



Haythu $a$ wa $b$ huma al-Dil'aan al-qiSiran (al-saqan), wa $c$ huwa

al-watar (al-Dil' al-TawiI).



\textbf{al-burhan al-hindisi:} yuqaddimu al-Khwarizmi burhanan

hindisiyyan jamilan bi-istikhdam murabba' maqsum ila arba'a murabba'at

Sagira mutasawiya wa-thamani muthallathat mutasawiya, mubayyinan anna

majmu' murabba'ayi al-Dil'ayni al-qiSirayn yusawi murabba' al-Dil' al-

TawiI.

\end{geombox}



\newpage



%% ========================================================================

%% CHAPTER: ON LEGACIES

%% ========================================================================

\chapter{fi al-wasaya wa-'ilm al-fara'iD}

\label{ch:legacies}



\begin{originaltext}

YatanaWalu al-juz' al-akhir min hadha al-'amal ahkam al-mirath al-

islami (al-fara'iD). Wa-hadha al-mawDu', ma'a annahu yabdU min

jinsi al-fiqh la al-riyadiyyat, yaHtaju li-jami' adawat al-jabr fi

Hallihi Hina turakkabu al-wasaya 'ala al-as-ham al-shar'iyya.

\end{originaltext}



\begin{commentary}

\textbf{al-jabr wa-l-fara'iD:}



YuZhiru hadha al-faSl -- alladhi ya'tabiruhu ba'Du al-mu'arrikhin

akthara li-l-ih-timam fi al-kitab -- kayfa anna al-jabr kana

\textbf{adatan Darruriyya} Hatta fi al-'ulum al-shar'iyya. fi al-

qanun al-islami li-l-mirath, hunaka \textit{as-ham thabita}

muHaddada fi al-qur'an al-karim (mithla: li-l-zawja al-thumnu,

li-l-abi al-sudus, ila akhirihi). Lakin 'indama yuDafu waSiyya

(taTaww'iyya) ila hadhihi al-as-ham al-thabita, taSbiH al-mas'ala

\textit{jabriyya}: naHtaju li-majhul ($x$) li-yumaththila al-mablagh

al-ijmali aw al-waSiyya.



Hadha huwa bi-ZabT ma daafa' al-Khwarizmi li-ta'lifi kitabihi --

fa-l-quDat wa-l-fuqaha' kanu bi-Haja li-Tariqa munazZama li-Halli

hadhihi al-masa'il.

\end{commentary}



\section*{fi ras al-mal wa-l-duyun}

\addcontentsline{toc}{section}{fi ras al-mal wa-l-duyun}



\begin{originaltext}

\textbf{mas'ala:} ``tawaffa rajulun wa-taraka ibnayn, wa-awSa

bi-thulthi ras mali-hi li-ajjanibi. Taraka 'asharati darahim

malan wa-'asharati darahim dainan 'ala ahadi al-ibnayn.''



al-Hall: tusammiya al-mablagh alladhi yustakhraj min al-dayni

shay'an. tuziIha ila ras al-mal alladhi huwa 'asharatu darahim.

a-yakunu al-majmuu' $10 + x$. TanquSu minhu thulthan, li-annahu

awSa bi-thulthi mali-hi: $3.33 + 0.33x$. al-Baqiy $6.67 + 0.67x$.

Taqsima hadha bayna al-ibnayn: yakulu kulla wahid minhuma $3.33 +

0.33x$. Wa-hadha yusawi al-shay' al-maTlub.

\end{originaltext}



\begin{mathexplain}

\textbf{al-Hall al-Riyadi:}



\begin{center}

$3.33 + 0.33x = x \quad\Rightarrow\quad 3.33 = 0.67x \quad\Rightarrow\quad x = 5$

\end{center}



YaHulu al-ajjanibi 'ala 5 darahim; wa-yaHulu kulla walad 'ala 5

darahim (yabqii al-ibnu al-madynun bi-daynihi kamilan ka-muqasSa

Hissatuhu).



\textbf{al-mu'adala:} nahnu nabHathu 'an $x$ alladhi yumaththilu

al-mablagh alladhi yuKHsamu min al-dayni li-yuzafa ila al-tarika

qabla qismat al-waSiyya.

\end{mathexplain}



\section*{fi naw' akhar min al-wasaya}

\addcontentsline{toc}{section}{fi naw' akhar min al-wasaya}



\begin{originaltext}

\textbf{mas'ala:} ``tawaffa rajulun tarikan ummahu wa-zawjatahu

wa-akhawayn wa-ukhtayn lahu min ab wa-umm, wa-awSa li-ajnanibi

bi-tsa'i ras mali-hi.''



al-Hall: tuj'alu as-hamu al-tarika min thamaniya wa-arba'in juz'an.

Wa-ta'lamu annaka idha akhadhta ts'an min ayyi ras mali yabqa

thamaniya atsa'ihi. azid al-an ila al-thamaniya atsa' thumnaha minha,

wa-ila al-thamaniya wa-arba'ayn aydan thumnahun, ay sitta, li-ikmal

ras mali-ka. Fa-hadha yu'ti arba'a wa-khamsin. yaHulu al-shakhS

alladhi yu'Sa lahu bi-ts'i 'ala sittatin min hadha. al-baqiyya

thamaniya wa-arba'un tuqSamu bayna al-warathah Hasaban li-as-hamihim

al-shar'iyya.

\end{originaltext}



\begin{mathexplain}

\textbf{al-Hall al-Riyadi:}



\begin{center}

\begin{tabular}{l|c}

\textbf{al-warith} & \textbf{al-sahm (min 48)} \\

\hline

al-umm & 8 \\

al-zawja & 6 \\

kullu akh & 14 \\

kullu ukht & 7 \\

\hline

al-majmu' & $8 + 6 + 2(14) + 2(7) = 56$ juz' mu'addal \\

\end{tabular}

\end{center}



\textbf{al-Khwarizmiyya:} ras al-mal = 54. al-waSiyya = 6 (ts' al-

54). al-baqiy = 48 yuqSamu bayna al-warathah. Hadhihi al-mas'ala

tataTallabu wijda maDami' mushtarak (48) yumkinu an yastaWi 'ala

jami' al-as-ham al-shar'iyya ba'da izalat al-waSiyya.

\end{mathexplain}



\section*{masa'il al-wasaya al-murakkaba}

\addcontentsline{toc}{section}{masa'il al-wasaya al-murakkaba}



\begin{originaltext}

\textbf{mas'ala:} ``tawaffat imra'atun tarikatan zawjaha wa-bnan

laha wa-ummaha. awSatat li-shakhsin bi-thuluthay maliha, wa-li-

akhar bi-rub'ihi. j'alat al-waSiyyatayn ma'an 'ala ibniha, wa-'ala

ummiha niSf (Hissati al-umm min al-baqiyya); 'ala zawjiha la

tulzimuhu illa al-thuluth (alladhi yajib an yusam bihi, Hasaba al-

qanun).''

\end{originaltext}



\begin{mathexplain}

\textbf{al-Hall al-Riyadi:}



\begin{center}

\begin{tabular}{l|c|l}

\textbf{al-warith} & \textbf{al-sahm al-asli} & \textbf{al-taSaruf} \\

\hline

al-ibn & $7/12$ & mukallaT bi- $\tfrac{2}{5} + \tfrac{1}{4} = \tfrac{13}{20}$ \\

al-zawj & $3/12 = 1/4$ & yusamim bi-thulthi sahmihi = $1/12$ \\

al-umm & $2/12 = 1/6$ & mukallafat bi-niSfi sahmihi = $1/12$ \\

\end{tabular}

\end{center}



\textbf{al-mu'adala:} al-musahimat al-ijmaliYya =

$\tfrac{13}{20} \times \tfrac{7}{12} + \tfrac{1}{12} + \tfrac{1}{12} = \text{muHassabat tanasubiyyan}$



Hadhihi al-mas'ala al-murakkaba tuZhiru al-quwwa al-Haqiqiyya li-l-

jabr fi Hall al-masa'il al-shar'iyya al-murakkaba allati kana

yuwajihuha al-quDat.

\end{mathexplain}



\newpage



%% ========================================================================

%% APPENDICES

%% ========================================================================

\appendix



\chapter{mulakHHas al-HaSalat al-sitt}

\label{app:six-cases}



\begin{mathexplain}

\textbf{al-HaSalat al-sitt al-kanuniyya li-l-mu'adlat al-tarb'iyya:}



\begin{center}

\renewcommand{\arraystretch}{1.8}

\begin{longtable}{|c|c|c|c|}

\hline

\textbf{al-HaSala} & \textbf{al-ism al-Arabi} & \textbf{al-mu'adala} & \textbf{al-Hall} \\

\hline

\endfirsthead

\hline

\textbf{al-HaSala} & \textbf{al-ism al-Arabi} & \textbf{al-mu'adala} & \textbf{al-Hall} \\

\hline

\endhead

\hline

\endfoot

al-ula & mal yusawi usuLan & $ax^2 = bx$ & $x = b/a$ \\

\hline

al-thaniya & mal yusawi 'adadan & $ax^2 = c$ & $x = \sqrt{c/a}$ \\

\hline

al-thalitha & usuL tusawi 'adadan & $bx = c$ & $x = c/b$ \\

\hline

al-rabi'a & mal wa-usuL yusawi 'adadan & $ax^2 + bx = c$ & $x = \sqrt{(b/2a)^2 + c/a} - b/2a$ \\

\hline

al-khamisa & mal wa-'adad yusawi usuLan & $ax^2 + c = bx$ & $x = \tfrac{b}{2a} \pm \sqrt{(b/2a)^2 - c/a}$ \\

\hline

al-sadisa & usuL wa-'adad yusawi malan & $bx + c = ax^2$ & $x = \tfrac{b}{2a} + \sqrt{(b/2a)^2 + c/a}$ \\

\hline

\end{longtable}

\end{center}



\vspace{0.5cm}



\textbf{mulahadha tarikhiyya:} Hadhihi al-HaSalat al-sitt tastanfudhu

jami'a imkaniyyat al-mu'adlat al-tarb'iyya dhat \textbf{al-ma'amilat

al-mujiba}. al-iqtiSir 'ala al-ma'amilat al-mujiba kana Darruriyyan

bi-sabab al-tafsir al-hindisi: fa-l-Khwarizmi yafhamu kulla Haddin

fi al-mu'adala 'ala annahu miqdar hindisi (masaHat wa-Tul), la

yumkinu an takuna al-misaHa saliba. wa-l-mu'alaja al-kamila li-jami'

al-mu'adlat al-tarb'iyya bima fiha al-ma'amilat al-saliba lam ta'ti

illa fi 'asr al-nahDa al-urubbiyya.

\end{mathexplain}



\newpage



\chapter{mu'jam al-muSTalaHat al-jabriyya al-Arabiyya}

\label{app:glossary}



\begin{center}

\renewcommand{\arraystretch}{1.6}

\begin{longtable}{|c|c|p{6cm}|}

\hline

\textbf{al-muSTalaH al-Arabi} & \textbf{al-nuskha al-latiniyya} & \textbf{al-ma'na wa-l-tafsir} \\

\hline

\endfirsthead

\hline

\textbf{al-muSTalaH al-Arabi} & \textbf{al-nuskha al-latiniyya} & \textbf{al-ma'na wa-l-tafsir} \\

\hline

\endhead

\hline

\endfoot

al-jabr (al-jabr) & algebra & idafat ma yanquSu min ahadi Tarfay al-mu'adala ila al-Taraf al-akhar li-istikmal al-tawazun \\

\hline

al-muqabala (al-muqabala) & almuchabala & hadf ma yatashabahu min al-Tarfayn li-tashil al-mu'adala \\

\hline

al-jadhr (al-jadhr) & radix & al-majhul dhatuhu ($x$) -- ``al-shay' '' al-majhul \\

\hline

al-mal (al-mal) & census & mirab' al-majhul ($x^2$) -- ``al-tharwa'' aw ``al-mumtalakat'' \\

\hline

al-'adad al-mufrad (al-'adad al-mufrad) & numerus simplex & al-Hadd al-thabit fi al-mu'adala (al-'adad al-ma'lum) \\

\hline

al-shay' (al-shay') & res/cosa & al-majhul ($x$) -- kalimat badIla li-l-jadhr \\

\hline

al-dirham (al-dirham) & drachma & wahidat al-nuqd/ al-qiyas -- tustakhdam kawwahida li-l-a'dad \\

\hline

al-waSiyya (al-waSiyya) & legatum & al-tarika aw al-waSiyya fi qanun al-mirath \\

\hline

al-fara'iD (al-fara'iD) & legata & al-as-ham al-ilzamiyya fi al-mirath al-islami \\

\hline

misaHa (misaHa) & mensura & misaHa al-ard -- Hisab al-masaaHat \\

\hline

Darb (Darb) & multiplicatio & 'amaliyyat al-Darb \\

\hline

qisma (qisma) & divisio & 'amaliyyat al-qisma \\

\hline

jam' (jam') & additio & 'amaliyyat al-jam' \\

\hline

TarH (TarH) & subtractio & 'amaliyyat al-TarH \\

\hline

jadhr (jadhr) & radix & al-jadhr al-tarb'i \\

\hline

muka''ab (muka''ab) & cubus & muka''ab al-majhul ($x^3$) \\

\hline

\end{longtable}

\end{center}



\begin{commentary}

\textbf{tatawwur al-muSTalaHat al-jabriyya:}



'inda tarjamat al-latiniyyin li-kutub al-riyadiyyat al-arabiyya fi

al-qarn al-thani 'ashar al-miladi (fi TulayTula wa-Siqilliyya),

Hawwalu al-muSTalaHat al-arabiyya ila latiniyya:



\begin{itemize}

  \item \textbf{al-jabr} $\rightarrow$ \textbf{algebra} (min ``al-jabr'')

  \item \textbf{al-shay'} $\rightarrow$ \textbf{res} (al-shay'/al-jawhar)

    $\rightarrow$ \textbf{cosa} (fi al-iTaliyya)

    $\rightarrow$ \textbf{coss} (fi al-almaniyya)

  \item \textbf{al-mal} $\rightarrow$ \textbf{census} (al-mumtalakat)

  \item \textbf{al-jadhr} $\rightarrow$ \textbf{radix} (al-jadhr)

  \item \textbf{al-dirham} $\rightarrow$ \textbf{drachma} (al-'umla)

\end{itemize}



Wa-kalimat ``\textbf{jabr}'' (algebra) hiya al-waHIda allati waSalat

ila al-lugha al-inkliZiyya mubashsharatan min al-arabiyya -- wa-hiya

iHda al-kalimat al-qalIla dhati al-aSl al-arabi fi al-riyadiyyat

al-inkliZiyya ila janib ``\textbf{Sifr}'' (zero, min al-Sifr)

wa-``\textbf{jayb}'' (sine, min al-jayb) wa-``\textbf{Khwarizmiyya}''

(algorithm, min al-Khwarizmi).

\end{commentary}



\newpage



%% ========================================================================

%% BIBLIOGRAPHIC NOTE

%% ========================================================================

\chapter*{mulahadha bIblIyughrafiyya}

\addcontentsline{toc}{chapter}{mulahadha bIblIyughrafiyya}

\markboth{mulahadha bIblIyughrafiyya}{mulahadha bIblIyughrafiyya}



\begin{historicalbox}

\textbf{al-maSadir al-makhTUTa:}



al-makhTUta al-ra'Isiyya li-kitab al-Khwarizmi \textit{al-mukhTasar fi

Hisab al-Jabr wa-l-Muqabala} maHfUZa fi maktabat Bodleyan bi-Oxford

(makhTUta Hunt. 214, awraq 1--123), nusikhat sanat 743 H / 1342 M.

Wa-hiya juz' min majlud CMXXVIII alladhi yaHtawi ayDan 'ala thalath

muSannafat ukhra fi al-Hisab wa-l-jabr.



al-makhTUta maktuba bi-KhaT wadiH wa-maqru', lakinnaha taniqu ila

Haddin kabiirin min al-nuqaT al-i'jamiyya allati tumayyizu al-Huruf

al-mutashabiha.

\end{historicalbox}



\vspace{0.3cm}



\begin{historicalbox}

\textbf{al-Taba'at wa-l-tarajim:}



\begin{enumerate}[label={[\arabic*]},rightmargin=1em]

  \item Frederic Rosen, \textit{The Algebra of Mohammed ben Musa},

  London: Oriental Translation Fund, 1831. (al-nas al-arabi ma'a

  al-tarjama al-inkliZiyya)



  \item Louis Charles Karpinski, \textit{Robert of Chester's Latin

  Translation of the Algebra of Al-Khowarizmi}, New York: Macmillan,

  1915. (al-nas al-latini ma'a al-nuskha al-inkliZiyya)



  \item 'Ali MuSTafa Musharrafa wa-Muhammad Mursi Ahmad (muhaqqiqan),

  \textit{Kitab al-Jabr wa-l-Muqabala}, al-Qahira: al-Jami'a al-

  miSriyya / maTba'at Fath Allah Ilyas Nuri, 1359 H / 1939 M.

  (al-Taba'a al-'arabiyya al-tahqiqiyya, maSr al-taDnId al-Hali)

\end{enumerate}

\end{historicalbox}



\vspace{0.3cm}



\begin{historicalbox}

\textbf{li-l-mazid min al-qira'a:}



\begin{itemize}[rightmargin=1em]

  \item J. Ruska, ``Zur aeltesten arabischen Algebra und Rechenkunst,''

  \textit{Sitzungsberichte der Heidelberger Akademie der Wissenschaften},

  Philosophisch-historische Klasse, Jahrgang 1917, 8. Abhandlung.



  \item H. Suter, ``Die Mathematiker und Astronomen der Araber und ihre

  Werke,'' \textit{Abhandlungen zur Geschichte der mathematischen

  Wissenschaften}, 10. Heft, Leipzig, 1900.



  \item R. Rashed (ed.), \textit{Al-Khwarizmi: Le commencement de

  l'algebre}, Paris: Blanchard, 2007. (Edition critique avec traduction

  francaise)



  \item B. van Dalen et al., ``Al-Khwarizmi's Astronomical Tables in

  a Tenth-Century Context,'' \textit{SCIAMVS}, 18 (2017), pp. 49--93.

\end{itemize}

\end{historicalbox}



\vfill



\begin{center}

\rule{0.3\textwidth}{0.4pt}



\vspace{0.3cm}



{\small\latinfont Modern LaTeX edition}\\[0.1cm]

{\small\latinfont Via XeLaTeX with polyglossia and fontspec}

\end{center}



\vspace{0.3cm}



\begin{tcolorbox}[colback=arabicbg,colframe=dropgold,arc=3pt,boxrule=0.6pt,width=0.5\textwidth,halign=center]

{\Large Tamma bi-Hamdi Allahi}

\end{tcolorbox}



\end{document}



